\chapter{The Modular Swiss-Cheese Operad: Local Approximation}
\label{ch:modular-sc-operad}
\label{ch:modular-swiss-cheese-operad}%% alias resolving legacy \ref sites

%% New macros for this chapter (providecommand to avoid clashes)
\providecommand{\cF}{\mathcal{F}}
\providecommand{\cG}{\mathcal{G}}
\providecommand{\SCmod}{\mathsf{SC}^{\mathrm{ch,top}}_{\mathrm{mod}}}
\providecommand{\SCmodcl}{\mathsf{SC}^{\mathrm{ch,top}}_{\mathrm{mod,cl}}}
\providecommand{\SCmodop}{\mathsf{SC}^{\mathrm{ch,top}}_{\mathrm{mod,op}}}
\providecommand{\SCmodmix}{\mathsf{SC}^{\mathrm{ch,top}}_{\mathrm{mod,mix}}}
\providecommand{\FTmod}{\mathrm{FT}_{\mathcal{M}\!od}}
\providecommand{\Commod}{\mathsf{Com}_{\mathrm{mod}}}
\providecommand{\Modop}{\mathcal{M}\!od}
\providecommand{\Sgp}{\mathfrak{S}}

%%%%%%%%%%%%%%%%%%%%%%%%%%%%%%%%%%%%%%%%%%%%%%%%%%%%%%%%%%%%%%%%%%%%%%%%%%%%%
% HONEST FRAMING
%%%%%%%%%%%%%%%%%%%%%%%%%%%%%%%%%%%%%%%%%%%%%%%%%%%%%%%%%%%%%%%%%%%%%%%%%%%%%

\noindent
All results in this chapter apply to logarithmic $\SCchtop$-algebras (Definition~\ref{def:log-SC-algebra}). Beyond the affine and $\cW$-algebra lineages, the question of which chiral algebras admit a logarithmic $\SCchtop$-algebra structure remains open.

The modular Swiss-cheese operad~$\SCmod$ extends $\SCchtop$ by
genus-raising clutching maps on the closed color while leaving the
open color at genus~$0$ (the interval~$\R$ has no genus). The
Feynman transform of $\SCmod$ produces the modular
Maurer--Cartan equation. Modular homotopy-Koszulity gives the flat
local/formal-disc modular bar--cobar equivalence; the curved
genus-$g$ coderived equivalence requires the factorization
Koszulity and chiral Riemann--Hilbert comparison of
Chapter~\ref{ch:factorization-swiss-cheese}. The emphasis
throughout is on honesty about scope: the topological operad
$\SCchtop$ and its modular extension $\SCmod$ model what happens
in a formal disc around a collision point, where the curve~$X$
is indistinguishable from~$\C$. They do not see global
information: $\cD$-module monodromy around cycles of~$\Sigma_g$,
period corrections from $H^1(\Sigma_g, \Z)$, or the Arakelov
propagator's non-holomorphic correction that produces the
curvature $\dfib^{\,2} = \kappa(\cA) \cdot \omega_g$. The global
story requires the factorization framework of
Chapter~\ref{ch:factorization-swiss-cheese}.

The local story is nonetheless essential. It explains \emph{why}
the corrected holomorphic differential
$\Dg{g}^{\,2} = 0$ is flat (the Fay trisecant identity is a
\emph{local} identity on configuration spaces), while the curvature
$\dfib^{\,2} = \kappa(\cA) \cdot \omega_g$ is a \emph{global} phenomenon.
The corrected holomorphic model lives in the operadic world; the curved
geometric model lives in the factorization world; the modular bar
complex mediates between them.

\begin{remark}[What the local model sees and what it misses]
\label{rem:local-model-scope}
\index{Swiss-cheese operad!local vs.\ global}
The formal-disc approximation captures:
\begin{itemize}
\item collision data: OPE coefficients, the Arnold relation and
 its Fay trisecant generalisation,
\item the Koszul dual $\cA^!$ (defined by bar-cobar on $\FM$
 compactifications),
\item the Swiss-cheese directionality (no open-to-closed),
\item the formal-disc propagator (Cauchy kernel times Heaviside).
\end{itemize}
It \emph{misses}:
\begin{itemize}
\item $\cD$-module monodromy around cycles of $\Sigma_g$,
\item period corrections from $H^1(\Sigma_g)$,
\item the Arakelov propagator's non-holomorphic correction,
\item the curvature $\kappa(\cA) \cdot \omega_g$.
\end{itemize}
Of the three models in Remark~\ref{rem:three-models}:
Model~1 (flat associated graded, $D_0^2 = 0$) and
Model~2 (corrected holomorphic, $\Dg{g}^{\,2} = 0$) are
fully captured by the local operad, because the Fay
trisecant identity is local. Model~3 (curved geometric,
$\dfib^{\,2} = \kappa \cdot \omega_g$) is \emph{not} captured:
its curvature is a global datum, invisible to any operadic
model. The derived--coderived comparison connecting the flat
models to the curved model requires the factorization
framework of Chapter~\ref{ch:factorization-swiss-cheese}.
\end{remark}


%%%%%%%%%%%%%%%%%%%%%%%%%%%%%%%%%%%%%%%%%%%%%%%%%%%%%%%%%%%%%%%%%%%%%%%%%%%%%
\section{Worked examples}
\label{sec:modular-sc-examples}
%%%%%%%%%%%%%%%%%%%%%%%%%%%%%%%%%%%%%%%%%%%%%%%%%%%%%%%%%%%%%%%%%%%%%%%%%%%%%

\providecommand{\Heis}{\mathsf{H}}
\providecommand{\slhat}{\widehat{\mathfrak{sl}}}
\providecommand{\Vir}{\mathsf{Vir}}
\providecommand{\hvee}{h^{\!\vee}}

Three examples of increasing complexity ($\Heis_\kappa$, $\slhat_2$ at level~$k$, $\Vir_c$) computing bar differentials, genus-raising maps, and curvature classes. Sections~\ref{sec:partially-modular-sc}--\ref{sec:modular-sc-consequences} abstract the structures visible here.

\subsection{The Heisenberg algebra $\Heis_\kappa$}
\label{subsec:example-heisenberg}

Let $\Heis_\kappa$ denote the rank-one Heisenberg vertex algebra
generated by a current $J(z)$ with OPE
\begin{equation}\label{eq:heis-ope}
 J(z)\, J(w) \;\sim\; \frac{\kappa}{(z - w)^2}\,.
\end{equation}
The bar complex $B(\Heis_\kappa)$ is an $E_1$ dg coassociative coalgebra over $(\mathrm{ChirAss})^!$ with:
\begin{itemize}
\item \emph{Holomorphic differential:}
 collision differential $D_0$ governed by
 $C_\ast(\FM_k(\C))$, extracting leading OPE coefficients
 from~\eqref{eq:heis-ope}.
\item \emph{Topological coproduct:}
 the $\Eone$ deconcatenation coproduct.
\end{itemize}

\begin{verification}[Arnold relation at degree~$3$, genus~$0$]
\label{ver:heis-arnold-3}
\index{Arnold relation!explicit verification at degree 3}
At degree $k = 3$, the bar element
$s^{-1}J \otimes s^{-1}J \otimes s^{-1}J \in B_3(\Heis_\kappa)$
parametrises configurations $(z_1, z_2, z_3) \in \FM_3(\C)$.
The differential $D_0$ extracts leading OPE coefficients at the three
pairwise collision strata $\partial_{12}$, $\partial_{23}$,
$\partial_{13}$ of $\FM_3(\C)$:
\begin{align*}
 D_0^{(12)} &\;=\; \kappa \cdot
 \bigl[\text{collision } z_1 \to z_2\bigr]
 = \kappa \cdot (s^{-1}\mathbf{1} \otimes s^{-1}J), \\
 D_0^{(23)} &\;=\; \kappa \cdot
 \bigl[\text{collision } z_2 \to z_3\bigr]
 = \kappa \cdot (s^{-1}J \otimes s^{-1}\mathbf{1}), \\
 D_0^{(13)} &\;=\; \kappa \cdot
 \bigl[\text{collision } z_1 \to z_3\bigr]
 = \kappa \cdot (s^{-1}\mathbf{1} \otimes s^{-1}J).
\end{align*}
(The collision $z_i \to z_j$ extracts the leading coefficient
of $J(z_i)\, J(z_j) \sim \kappa/(z_i - z_j)^2$,
replacing the pair by $\kappa \cdot \mathbf{1}$.)

The condition $D_0^2 = 0$ at degree~$3$ requires that the
sum of pairwise collision compositions vanishes. Explicitly,
$D_0^2$ on a $3$-fold bar element involves
$D_0^{(12)} \circ D_0^{(23)} + D_0^{(23)} \circ D_0^{(13)}
+ D_0^{(13)} \circ D_0^{(12)}$. Each double collision produces
$\kappa^2 \cdot s^{-1}\mathbf{1}$, and the three contributions
have signs $+1$, $+1$, $-2$ arising from the Arnold relation
\[
 \omega_{12} \wedge \omega_{23}
 + \omega_{23} \wedge \omega_{13}
 + \omega_{13} \wedge \omega_{12} = 0
\]
on $\FM_3(\C)$
(equation~\eqref{eq:genus-g-arnold} at $g = 0$), where
$\omega_{ij} = d\log(z_i - z_j)$.
The leading OPE coefficient at each boundary stratum of $\FM_3(\C)$
is controlled by the coefficient of~$\omega_{ij}^{\wedge 2}$, and the
Arnold relation forces the total boundary contribution to
vanish. Hence $D_0^2 = 0$.
\end{verification}

\begin{verification}[Genus-$1$ corrected differential]
\label{ver:heis-genus-1}
\index{Fay identity!genus 1 verification}
At genus~$1$, the curve $\Sigma_1$ is an elliptic curve
$\C / (\Z + \tau\Z)$ with period~$\tau$. The propagator
is the Weierstrass zeta function
$\zeta(z; \tau) = \frac{1}{z}
+ \sum_{\omega \in \Lambda \setminus 0}
\bigl(\frac{1}{z - \omega} + \frac{1}{\omega}
+ \frac{z}{\omega^2}\bigr)$,
which plays the role of $f(z,w) = (z-w)^{-1} + O(1)$.

The corrected genus-$1$ differential on~$B(\Heis_\kappa)$ is
\[
 \Dg{1} = D_0 + \tau \cdot D_{\mathrm{ell}},
\]
where $D_{\mathrm{ell}}$ uses the elliptic propagator in place
of the rational one. The flatness $\Dg{1}^{\,2} = 0$ follows
from the Fay trisecant identity at genus~$1$.

In the elliptic case, the Fay identity
(equation~\eqref{eq:fay-trisecant}) for
the Weierstrass $\sigma$-function
$\sigma(z;\tau)$ (whose logarithmic derivative is
$\zeta(z;\tau) = \sigma'(z)/\sigma(z)$) specialises to:
for distinct $u, v \in \C / \Lambda$,
\begin{equation}\label{eq:fay-genus1}
 \zeta(u + v) - \zeta(u) - \zeta(v)
 = \frac{1}{2}\,
 \frac{\wp'(u) - \wp'(v)}{\wp(u) - \wp(v)}\,,
\end{equation}
where $\wp(z) = -\zeta'(z)$ is the Weierstrass $\wp$-function.
(This is the classical addition formula for the Weierstrass
zeta function, which is the genus-$1$ degeneration of the
Fay trisecant.)

The three-term identity
$f(u,v) f(w,v) + f(v,w) f(u,w) + f(w,u) f(v,u) = 0$
where $f(u,v) = \zeta(u-v) + \eta_1(u - v)$
(with $\eta_1 = \zeta(1/2)$ the quasi-period) follows
from~\eqref{eq:fay-genus1} by the same mechanism as the
genus-$0$ Arnold relation: the three double-residue
contributions to $\Dg{1}^{\,2}$ at degree~$3$ are
proportional to $f(z_1,z_2) f(z_2,z_3)$,
$f(z_2,z_3) f(z_3,z_1)$, and $f(z_3,z_1) f(z_1,z_2)$
respectively, and their sum vanishes by the Fay identity.
\end{verification}

\begin{verification}[Curvature at genus~$1$]
\label{ver:heis-curvature}
\index{curvature!Heisenberg algebra}
The curved geometric model on a fixed elliptic
curve~$\Sigma_1$ uses the Arakelov propagator
\[
 G_{\mathrm{Ar}}(z,w)
 = \log\bigl|\vartheta_1(z - w; \tau)\bigr|^2
 - \frac{2\pi\, (\operatorname{Im}(z-w))^2}{\operatorname{Im} \tau}
 + C_1(\tau),
\]
where $\vartheta_1$ is the odd Jacobi theta function and
$C_1(\tau)$ is a normalisation constant. The curvature is
\[
 \dfib^{\,2}
 = \kappa \cdot \omega_1,
\]
where $\omega_1 = \frac{i}{2\operatorname{Im}\tau}\,
dz \wedge d\bar{z}$
is the Arakelov $(1,1)$-form on $\Sigma_1$, normalised so that
$\int_{\Sigma_1} \omega_1 = 1$
(since $\int_{\C/\Lambda} \frac{i}{2} dz \wedge d\bar{z}
= \operatorname{Im}\tau$). The factor of $\kappa$ comes
from the OPE level~\eqref{eq:heis-ope}: each propagator line
carries a factor of~$\kappa$.

This is a global datum: $\omega_1$ depends on the period
$\tau \in \mathfrak{h}$, hence on the moduli of~$\Sigma_1$.
The formal-disc functor $\mathrm{Loc}$
(Construction~\ref{con:extraction-functor}) sees only the local
part $\log|z-w|^2$ of the propagator, so it misses the
$(\operatorname{Im}(z-w))^2/\operatorname{Im}\tau$ correction
that produces the curvature.
\end{verification}


\subsection{Genus spectral sequence for $\Heis_\kappa$: explicit computation}
\label{subsec:heisenberg-spectral-sequence}

\providecommand{\Deltacyc}{\Delta_{\mathrm{cyc}}}

Since $\Heis_\kappa$ is Koszul (bar complex generated in degree~$2$) and
formal ($m_k = 0$ for $k \ge 3$), the genus-$0$ bar complex is completely understood and the genus-$1$ correction computable. The dichotomy: at $\kappa \neq 0$, genus-$1$ cohomology is killed by the curvature; at $\kappa = 0$, the spectral sequence degenerates.

\begin{computation}[Genus-$0$ bar complex of $\Heis_\kappa$]
\label{comp:heis-genus0-bar}
\index{Heisenberg algebra!bar complex!genus 0}
The genus-$0$ bar complex $B(\Heis_\kappa)$ with differential
$D_0$ is computed as follows. The chiral algebra $\Heis_\kappa$
is generated by the current $J(z)$ of conformal weight~$1$;
write $a = s^{-1}J$ for the desuspended bar generator.
The bar element of degree~$k$ is
$a^{\otimes k} \in B_k(\Heis_\kappa)$, and the
differential $D_0$ extracts pairwise leading OPE coefficients:
\[
 D_0(a^{\otimes k})
 = \sum_{1 \le i < j \le k}
 \pm\, \kappa \cdot
 a^{\otimes (i-1)} \otimes
 s^{-1}\mathbf{1} \otimes
 a^{\otimes (k-i-1)},
\]
where the collision $J(z_i)\, J(z_j) \sim \kappa/(z_i - z_j)^2$
produces the leading coefficient
$\kappa \cdot \mathbf{1}$ and the sign is the Koszul sign
from the desuspension.

Since $\Heis_\kappa$ is Koszul (the bar complex is generated
in degree~$2$; cf.\ the one-loop exactness criterion of
Theorem~\ref{thm:one-loop-koszul} applied to the free-field
BV theory) and formal ($m_k = 0$ for $k \ge 3$, since the OPE
of currents has no poles beyond the double pole), the
cohomology of $(B(\Heis_\kappa), D_0)$ is concentrated in
degree~$1$:
\begin{equation}\label{eq:heis-genus0-cohomology}
 H^\ast(B(\Heis_\kappa), D_0)
 \;\cong\; \Bbbk[a]
 \;=\; \Bbbk \cdot 1 \oplus \Bbbk \cdot a
 \oplus \Bbbk \cdot a^2 \oplus \cdots,
\end{equation}
the free graded-commutative algebra on one generator~$a$ of
degree~$1$. This is the Koszul dual coalgebra of $\Heis_\kappa$:
the commutative algebra $\Bbbk[a]$ is Koszul dual to the
free algebra on one generator, which is the underlying
associative structure of the Heisenberg current. The polynomial
$a^n$ represents the class of $s^{-1}a^{\otimes n}$ in bar
cohomology.
\end{computation}

\begin{computation}[Genus filtration and the $E_1$ page]
\label{comp:heis-genus-filtration}
\index{Heisenberg algebra!spectral sequence!E1 page@$E_1$ page}
The modular bar complex $\Bmod(\Heis_\kappa)$ carries the genus
filtration $F^{\le g}$ of
Corollary~\ref{cor:genus-spectral-convergence}. The
associated graded pieces are:
\begin{align*}
 \gr^0 &= (B(\Heis_\kappa),\, D_0)
 &&\text{(the genus-$0$ bar complex)}, \\
 \gr^1 &= (B^{(1)}(\Heis_\kappa),\, D_0)
 &&\text{(the genus-$1$ correction)},
\end{align*}
where $B^{(1)}(\Heis_\kappa)$ denotes the genus-$1$ summand
of the modular bar complex. The genus-$1$ correction arises
from the non-separating clutching
$D_1 \colon B_k^{(g)} \to B_{k-2}^{(g+1)}$,
which contracts a pair of bar entries using the invariant
bilinear form. For $\Heis_\kappa$, the bilinear form is
$(a, a) = \kappa$ (the level), and the contraction operator is
\[
 \Deltacyc :=
 \sum_{1 \le i < j \le k}
 \text{(contract entries $i$ and $j$ using $({\cdot},{\cdot})$)}.
\]

The $E_1$ page of the genus spectral sequence is the
cohomology of each associated graded piece with respect
to $D_0$:
\begin{equation}\label{eq:heis-E1-page}
\begin{aligned}
 E_1^{0,\ast}
 &= H^\ast(\gr^0, D_0) = \Bbbk[a]
 &&\text{(genus-$0$ cohomology)}, \\
 E_1^{1,\ast}
 &= H^\ast(\gr^1, D_0) = \Bbbk[a]
 &&\text{(genus-$1$ cohomology at $E_1$)}.
\end{aligned}
\end{equation}
The identification $E_1^{1,\ast} = \Bbbk[a]$ holds because
the genus-$1$ associated graded $\gr^1$ is, as a chain complex
under $D_0$, isomorphic to $B(\Heis_\kappa)$ (the same
underlying complex, with the genus label shifted by~$1$).
This uses the product decomposition
(Proposition~\ref{prop:mixed-product-decomposition},
Remark~\ref{rem:product-strictness}): the
non-separating clutching $D_1$ raises genus but does not change
the $D_0$-differential, which depends only on the OPE data
(a local datum, unaffected by genus). More precisely, the
genus-$1$ bar complex $B^{(1)}(\Heis_\kappa)$ consists of
stable graphs of total genus~$1$, and the $D_0$-differential
on each such graph acts by the \emph{same} OPE collision
residues as at genus~$0$: the genus label of the graph
does not enter the $D_0$-formula. The homotopy equivalence
$C_\ast(\FM_{k|m}(\Sigma_1, \partial)) \simeq
C_\ast(\FM_k(\Sigma_1)) \otimes C_\ast(\Eone(m))$
ensures that the chain-level $D_0$-cohomology at genus~$1$
is isomorphic to the genus-$0$ cohomology after passing
through the quasi-isomorphism of
Remark~\ref{rem:product-strictness}.
\end{computation}

\begin{computation}[The $d_1$ differential and spectral sequence collapse]
\label{comp:heis-d1-differential}
\index{Heisenberg algebra!spectral sequence!d1 differential@$d_1$ differential}
The $d_1$-differential of the genus spectral sequence maps
$E_1^{0,\ast} \to E_1^{1,\ast}$ and is induced by the
genus-raising operator $D_1 = \kappa \cdot \Deltacyc$ on bar
cohomology. We compute it explicitly.

\medskip\noindent
\emph{Action on generators.}
On the generator $a \in E_1^{0,1} = \Bbbk \cdot a$
(represented by $s^{-1}a \in B_1$), the operator $D_1$ is
zero: there is only one bar entry, so no pair can be contracted.
Hence $d_1(a) = 0$.

On $a^2 \in E_1^{0,2}$, represented by
$s^{-1}a \otimes s^{-1}a \in B_2$, the contraction gives
\[
 D_1(s^{-1}a \otimes s^{-1}a)
 = (a, a) \cdot \xi_{\mathrm{nsep}}(\mathbf{1})
 = \kappa \cdot \xi_{\mathrm{nsep}}(\mathbf{1})
 \;\in\; B_0^{(1)}.
\]
Passing to cohomology: $\xi_{\mathrm{nsep}}(\mathbf{1})$
represents the unit $1 \in E_1^{1,0} = \Bbbk$. Hence
$d_1(a^2) = \kappa \cdot 1 \in E_1^{1,0}$.

More generally, $d_1(a^n)$ is computed by contracting all
pairs in $s^{-1}a^{\otimes n}$, yielding
\begin{equation}\label{eq:heis-d1-formula}
 d_1(a^n) = \binom{n}{2}\, \kappa \cdot a^{n-2}
 \;\in\; E_1^{1, n-2},
\end{equation}
since there are $\binom{n}{2}$ pairs to contract, each
contributing $(a,a) = \kappa$. The factor
$\omega_1$ (the Arakelov form on $\Sigma_1$, normalised to
$\int_{\Sigma_1}\omega_1 = 1$) is absorbed into the
identification of genus-$1$ classes.

\medskip\noindent
\emph{The $E_2$ page: two regimes.}

\medskip\noindent
\textbf{Case 1: $\kappa \neq 0$ (generic level).}
The $d_1$-differential~\eqref{eq:heis-d1-formula} is
$d_1(a^n) = \binom{n}{2}\kappa \cdot a^{n-2}$, which
is multiplication by $\binom{n}{2}\kappa$ composed with the
degree-shift $a^n \mapsto a^{n-2}$. Since $\kappa \neq 0$,
this map is an isomorphism in the appropriate degrees:
$d_1 \colon E_1^{0,n} \to E_1^{1,n-2}$ is an isomorphism for
all $n \ge 2$, and $d_1 = 0$ on $E_1^{0,0}$ and $E_1^{0,1}$.
The $E_2$ page is:
\begin{equation}\label{eq:heis-E2-generic}
 E_2^{p,q} = 0 \quad \text{for } p \ge 1 \text{ or } q \ge 2.
\end{equation}
More precisely, the only surviving entries are
$E_2^{0,0} = \Bbbk$ and $E_2^{0,1} = \Bbbk \cdot a$.
We verify each vanishing.
\begin{itemize}
\item \emph{Column $p = 0$, row $q \ge 2$:} the
 differential $d_1 \colon E_1^{0,q} \to E_1^{1,q-2}$
 sends $a^q \mapsto \binom{q}{2}\kappa\, a^{q-2}$.
 Since $\kappa \neq 0$ and $\binom{q}{2} \ge 1$ for $q \ge 2$,
 the map is injective on the one-dimensional space $\Bbbk \cdot a^q$,
 so $\ker(d_1|_{E_1^{0,q}}) = 0$ and $E_2^{0,q} = 0$.
\item \emph{Column $p = 0$, rows $q = 0,1$:}
 $d_1(1) = 0$ and $d_1(a) = 0$ (no pair to contract),
 so both survive: $E_2^{0,0} = \Bbbk$ and $E_2^{0,1} = \Bbbk \cdot a$.
\item \emph{Column $p = 1$, all rows $q \ge 0$:}
 the map $d_1 \colon E_1^{0,q+2} \to E_1^{1,q}$ sends
 $a^{q+2} \mapsto \binom{q+2}{2}\kappa\, a^q$, which is
 surjective onto $E_1^{1,q} = \Bbbk \cdot a^q$ (since
 $\binom{q+2}{2} \ge 1$). Meanwhile $d_1 \colon E_1^{1,q} \to E_1^{2,q-2} = 0$
 (there is no genus-$2$ column in this truncation),
 so $\ker = E_1^{1,q}$ and
 $E_2^{1,q} = E_1^{1,q}/\im(d_1) = 0$.
\item \emph{Columns $p \ge 2$:} these are the genus-$\ge 2$
 summands of the modular bar complex, which do not participate
 in the $d_1$-differential between columns $0$ and~$1$.
 The analysis at higher genus requires additional differentials
 $d_2, d_3, \ldots$ of the spectral sequence. For the
 Heisenberg algebra, these all vanish because the invariant form
 is $2$-step nilpotent (each contraction reduces degree by~$2$),
 and the spectral sequence collapses at $E_2$ in the complete
 modular bar complex.
\end{itemize}

At $\kappa \neq 0$, the curvature $\kappa \cdot \omega_1$ kills genus-$1$ cohomology beyond the lowest degrees: $D_1$ provides the connecting differential, making $D = D_0 + D_1$ acyclic in high degrees.

\medskip\noindent
\textbf{Case 2: $\kappa = 0$ (critical level).}
At $\kappa = 0$, the bilinear form vanishes:
$(a, a)_{\kappa=0} = 0$. Hence $D_1 = 0$ identically, and
$d_1 = 0$. The spectral sequence degenerates at $E_1$:
\begin{equation}\label{eq:heis-E2-critical}
 E_\infty^{p,\ast} = E_1^{p,\ast} = \Bbbk[a]
 \quad \text{for all } p \ge 0.
\end{equation}
The genus-$1$ cohomology $H^\ast(\Bmod(\Heis_0)^{(1)})$ is
isomorphic to the genus-$0$ cohomology $\Bbbk[a]$. The
spectral sequence degenerates because there is no genus-raising
differential to connect the filtration levels.

At $\kappa = 0$, the OPE is trivial: the modular bar complex is a direct sum $\bigoplus_{g \ge 0} B^{(g)}$ with no inter-genus coupling, each summand having cohomology $\Bbbk[a]$.

The dichotomy is the spectral-sequence shadow of curvature: $\kappa(\Heis_\kappa) = \kappa$, and $\kappa \cdot \omega_g$ obstructs degeneration. When $\kappa = 0$ the genus tower decouples; when $\kappa \neq 0$ it is connected by
the curvature into an acyclic complex.
\end{computation}


\subsection{The affine algebra $\slhat_2$ at level $k$}
\label{subsec:example-sl2}

\begin{example}[$\slhat_2$ at level $k$: the degree-$2$ bar differential]
\label{ex:sl2-bar}
\index{affine Lie algebra!sl2@$\widehat{\mathfrak{sl}}_2$!bar differential}
Let $\{e, f, h\}$ be the Chevalley generators of
$\mathfrak{sl}_2$, with $[h,e] = 2e$, $[h,f] = -2f$,
$[e,f] = h$. The affine vertex algebra $\slhat_2$ at
level~$k$ has OPEs
\begin{equation}\label{eq:sl2-opes}
\begin{aligned}
 e(z)\, f(w) &\;\sim\; \frac{k}{(z-w)^2} + \frac{h(w)}{z-w}\,, \\
 h(z)\, e(w) &\;\sim\; \frac{2\, e(w)}{z-w}\,,
 \qquad
 h(z)\, f(w) \;\sim\; \frac{-2\, f(w)}{z-w}\,,
 \qquad
 h(z)\, h(w) \;\sim\; \frac{2k}{(z-w)^2}\,.
\end{aligned}
\end{equation}

\medskip\noindent
\emph{The genus-$0$ differential $D_0$.}
Consider the degree-$2$ bar element
$s^{-1}e \otimes s^{-1}f \in B_2(\slhat_2)$.
The differential $D_0$ extracts the residue at $z_1 \to z_2$
from the OPE $e(z_1)\, f(z_2)$.
The simple-pole residue gives the commutator $[e, f] = h$:
\[
 D_0(s^{-1}e \otimes s^{-1}f)
 = s^{-1} h
 \quad \text{(from the $(z-w)^{-1}$ term)}.
\]
The second-order pole $k/(z-w)^2$ does \emph{not} contribute
to $D_0$ at degree~$2$: it produces a constant (proportional
to $k \cdot \mathbf{1}$), which lives in $B_1 = s^{-1}\cA$ as
$k \cdot s^{-1}\mathbf{1}$. In the bar complex, this
contribution appears as part of the curvature datum, not as
a differential.

\medskip\noindent
\emph{The genus-raising differential $D_1$.}
The non-separating clutching
$D_1 \colon B_2^{(g)} \to B_0^{(g+1)}$ contracts the two
entries of a degree-$2$ bar element using the invariant
bilinear form. For $\slhat_2$ at level~$k$, the
normalised Killing form is
$(x, y)_k = k \cdot \operatorname{Tr}(\operatorname{ad}(x)
\operatorname{ad}(y)) / (2\hvee)$
with $\hvee = 2$ for $\mathfrak{sl}_2$. Thus:
\begin{equation}\label{eq:sl2-D1}
 D_1(s^{-1}e \otimes s^{-1}f)
 = (-1)^{\degree{e}+1}\,
 (e, f)_k \cdot \xi_{\mathrm{nsep}}(\mathbf{1})
 = (-1)^{1}\, \frac{k \cdot \operatorname{Tr}(\operatorname{ad}(e)
 \operatorname{ad}(f))}{2 \cdot 2} \cdot
 \xi_{\mathrm{nsep}}(\mathbf{1}).
\end{equation}
Since $\operatorname{Tr}(\operatorname{ad}(e)\operatorname{ad}(f))
= \operatorname{Tr}\bigl(\begin{smallmatrix} 0 & -2 \\ 0 & 0
\end{smallmatrix}\bigr)
\bigl(\begin{smallmatrix} 0 & 0 \\ 2 & 0
\end{smallmatrix}\bigr)
= \operatorname{Tr}\bigl(\begin{smallmatrix} -4 & 0 \\ 0 & 0
\end{smallmatrix}\bigr)
= -4$
(computed in the adjoint representation on
$\{e, f, h\}$, where
$\operatorname{ad}(e)(f) = h$,
$\operatorname{ad}(e)(h) = -2e$,
$\operatorname{ad}(f)(e) = -h$,
$\operatorname{ad}(f)(h) = 2f$),
we obtain
\[
 D_1(s^{-1}e \otimes s^{-1}f)
 = (-1)^{1}\,\frac{k \cdot (-4)}{4}\,
 \xi_{\mathrm{nsep}}(\mathbf{1})
 = k \cdot \xi_{\mathrm{nsep}}(\mathbf{1}).
\]
The output $k \cdot \xi_{\mathrm{nsep}}(\mathbf{1})$
lives in $B_0^{(g+1)}$, a genus-$(g+1)$ scalar weighted
by the level~$k$.
\end{example}

\begin{verification}[Bicomplex relation $D_0 D_1 + D_1 D_0 = 0$
for $\slhat_2$]
\label{ver:sl2-bicomplex}
\index{bicomplex!sl2@$\widehat{\mathfrak{sl}}_2$ verification}
We verify the bicomplex identity on
$s^{-1}e \otimes s^{-1}f \in B_2(\slhat_2)$.

\medskip\noindent
\emph{The term $D_1 D_0$.}
From above, $D_0(s^{-1}e \otimes s^{-1}f) = s^{-1}h$.
Now $D_1$ on a degree-$1$ bar element $s^{-1}h$ is zero:
the non-separating clutching requires two inputs to contract,
but $s^{-1}h \in B_1$ has only one tensor factor. Hence
\[
 D_1 D_0(s^{-1}e \otimes s^{-1}f) = D_1(s^{-1}h) = 0.
\]

\medskip\noindent
\emph{The term $D_0 D_1$.}
From~\eqref{eq:sl2-D1},
$D_1(s^{-1}e \otimes s^{-1}f)
= k \cdot \xi_{\mathrm{nsep}}(\mathbf{1})$.
This is a scalar in $B_0^{(g+1)}$ (bar degree zero, genus $g+1$).
The differential $D_0$ on a bar-degree-zero element is zero
(there are no collision strata for zero points). Hence
\[
 D_0 D_1(s^{-1}e \otimes s^{-1}f) = 0.
\]

\medskip\noindent
Therefore $D_0 D_1 + D_1 D_0 = 0 + 0 = 0$ on this element,
as required by Proposition~\ref{prop:D0D1}.
\end{verification}

\begin{verification}[Codimension-$2$ cancellation at degree~$3$
for $\slhat_2$]
\label{ver:sl2-codim2}
\index{bicomplex!sl2@$\widehat{\mathfrak{sl}}_2$!codimension-2 cancellation}
We verify $D_0^2 = 0$ on the degree-$3$ bar element
$s^{-1}e \otimes s^{-1}f \otimes s^{-1}h \in B_3(\slhat_2)$
by exhibiting a cancelling pair. The differential $D_0$ acts
by extracting pairwise OPE residues at the three collision
strata $\partial_{12}$, $\partial_{23}$, $\partial_{13}$ of
$\FM_3(\C)$. Applying $D_0$ twice, consider the two routes
through the $(e,f)$ and $(f,h)$ collisions:

\medskip\noindent
\emph{Route 1: $(e,f)$ first, then $h$.}
The collision $z_1 \to z_2$ extracts $[e,f] = h$ from the
simple-pole residue, yielding $s^{-1}h \otimes s^{-1}h$.
Then $D_0$ on $s^{-1}h \otimes s^{-1}h$ extracts
$[h,h] = 0$ (since $h(z)\,h(w) \sim 2k/(z-w)^2$ has no
simple pole). This route contributes zero.

\medskip\noindent
\emph{Route 2: $(f,h)$ first, then $e$.}
The collision $z_2 \to z_3$ extracts $[h,f] = -2f$ from the
OPE $h(z)\,f(w) \sim -2f(w)/(z-w)$, yielding
$s^{-1}e \otimes s^{-1}(-2f) = -2\,(s^{-1}e \otimes s^{-1}f)$.
Then $D_0$ on $s^{-1}e \otimes s^{-1}f$ extracts $[e,f] = h$,
giving $-2\,s^{-1}h$.

\medskip\noindent
\emph{Route 3: $(e,h)$ first, then $f$.}
The collision $z_1 \to z_3$ (with Koszul sign
$(-1)^{\degree{f}-1} = (-1)^{0} = +1$ from passing
$s^{-1}f$) extracts $[h,e] = 2e$, yielding
$s^{-1}f \otimes s^{-1}(2e) = 2\,(s^{-1}f \otimes s^{-1}e)$.
Then $D_0$ extracts $[f,e] = -h$, giving $2 \cdot (-s^{-1}h) =
-2\,s^{-1}h$.

\medskip\noindent
\emph{Cancellation.}
The total contribution is
$0 + (-2\,s^{-1}h) \cdot \omega_{23}\wedge\omega_{12}
+ (-2\,s^{-1}h) \cdot \omega_{13}\wedge\omega_{12}$.
By the Arnold relation
$\omega_{12}\wedge\omega_{23} + \omega_{23}\wedge\omega_{13}
+ \omega_{13}\wedge\omega_{12} = 0$
on $\FM_3(\C)$, these two terms cancel, confirming $D_0^2 = 0$.
This cancellation is the codimension-$2$ instance of the general
mechanism: the Arnold relation on propagator forms kills the
double-residue contributions at every degree.
\end{verification}


\subsection{The Virasoro algebra $\Vir_c$: non-formality
and shadow depth}
\label{subsec:example-virasoro}

\begin{example}[Non-formality of $\Vir_c$ and shadow depth]
\label{ex:virasoro-nonformality}
\index{Virasoro algebra!non-formality}
\index{shadow depth!Virasoro algebra}
The Virasoro vertex algebra $\Vir_c$ is generated by the
stress-energy tensor $T(z)$ with OPE
\begin{equation}\label{eq:virasoro-ope}
 T(z)\, T(w) \;\sim\;
 \frac{c/2}{(z-w)^4}
 + \frac{2\, T(w)}{(z-w)^2}
 + \frac{\partial T(w)}{z-w}\,.
\end{equation}
The degree-$4$ pole $c/2 \cdot (z-w)^{-4}$ is the source of
non-formality. In the language of the modular bar complex:

\medskip\noindent
\emph{The ternary operation $m_3 \neq 0$.}
The $\Ainf$-structure on $B(\Vir_c)$ has operations
$m_k$ for all $k \ge 2$, defined by residues on $\FM_k(\C)$.
The binary operation $m_2$ extracts the simple-pole
coefficient: $m_2(T, T) = \partial T$ (the $(z-w)^{-1}$ term
in~\eqref{eq:virasoro-ope}).
The ternary operation $m_3$ is determined by the
higher-pole data. Concretely, $m_3$ arises from the boundary
strata of $\FM_3(\C)$ where two of the three points collide
and the collision produces a second-order pole, which then
collides with the third point. The $(z-w)^{-2}$ coefficient
$2T(w)$ in~\eqref{eq:virasoro-ope} feeds into the composition
$m_2(m_2^{(2)}(T, T), T)$, where $m_2^{(2)}$ denotes the
second-order residue.

The degree-$4$ pole contributes directly to $m_3$: the residue
of $c/2 \cdot (z_1 - z_2)^{-4}$ integrated against the
propagator form on $\FM_3(\C)$ yields a nonzero ternary
operation. Hence $m_3 \neq 0$ whenever $c \neq 0$.

\medskip\noindent
\emph{The modular characteristic.}
The curvature of $\Vir_c$ is
$\kappa(\Vir_c) = c/2$,
reading off the coefficient of the highest-order pole
in~\eqref{eq:virasoro-ope}.
This equals the leading coefficient in Volume~I's
Theorem~D: the curvature $\dfib^{\,2} = (c/2) \cdot \omega_g$
governs the genus tower.

\medskip\noindent
\emph{Shadow depth $d_{\mathrm{alg}}(\Vir_c) = \infty$ (class~$\mathbf{M}$).}
The $A_\infty$ tower does not terminate: $m_k \neq 0$ on
cohomology for all $k \ge 3$. Although $m_3$ is the
only \emph{independent} higher generator (all $m_k$ for
$k \ge 4$ are determined by $m_2$ and $m_3$ via the
Stasheff relations), the determined operations are generically
nonzero. This is class~$\mathbf{M}$ in the Volume~I
classification.
\end{example}

\begin{verification}[Generating depth: determination of $m_4$ from $m_2$ and $m_3$]
\label{ver:virasoro-m4}
\index{Virasoro algebra!Stasheff relations at degree 4}
We verify that $m_4$ is determined by $m_2$ and $m_3$ at
degree~$4$. The Stasheff relation at degree~$4$ reads:
\begin{equation}\label{eq:stasheff-4}
 \sum_{i+j = 5} \sum_{\ell}
 (-1)^{\epsilon}\, m_i(a_1, \ldots, a_{\ell-1},
 m_j(a_\ell, \ldots, a_{\ell+j-1}),
 a_{\ell+j}, \ldots, a_4) = 0,
\end{equation}
where the sign $(-1)^{\epsilon}$ is the Koszul sign from
desuspension. Expanding for $(a_1, a_2, a_3, a_4) =
(T, T, T, T)$, the terms involving $m_4$ are:
\[
 m_1(m_4(T,T,T,T))
 + m_4(m_1(T), T, T, T) + \cdots = 0.
\]
Since $m_1 = 0$ (there is no linear differential in
$\Vir_c$), the $m_1 \circ m_4$ terms vanish, and the
relation becomes
\begin{multline}\label{eq:stasheff-4-virasoro}
 m_2(m_3(T,T,T), T) - m_3(m_2(T,T), T, T)
 + m_3(T, m_2(T,T), T) \\
 - m_3(T, T, m_2(T,T))
 + m_2(T, m_3(T,T,T)) = 0,
\end{multline}
up to Koszul signs (which we suppress for readability, since
$T$ has even degree $\degree{T} = 2$, so all signs are $+1$).

All terms in~\eqref{eq:stasheff-4-virasoro} involve only
$m_2$ and $m_3$. The equation is \emph{not} a constraint on
$m_2$ and $m_3$ (which are already determined by the OPE) but
rather a \emph{definition} of $m_4$: the Stasheff relation
determines $m_4$ as the unique operation such that the full
$\Ainf$ relation at degree~$4$ is satisfied. Explicitly,
$m_4(T,T,T,T)$ is the element forced by requiring that
$D_0^2 = 0$ at degree~$4$ on $\FM_4(\C)$, and the boundary
strata of $\FM_4(\C)$ contributing to $D_0^2$ decompose
entirely into compositions of $m_2$ and $m_3$. No new
information from the OPE is required.

The same argument applies inductively at all degrees $k \ge 5$:
the Stasheff relation at degree~$k$ determines $m_k$ in terms
of $\{m_2, m_3, \ldots, m_{k-1}\}$, and since $m_3$ is the
only independent operation beyond $m_2$, the generating depth is
$d_{\mathrm{gen}}(\Vir_c) = 3$ (all higher operations determined
by $m_3$), while the algebraic depth is
$d_{\mathrm{alg}}(\Vir_c) = \infty$ (class~$\mathbf{M}$: the
determined operations are generically nonzero).

The Virasoro OPE~\eqref{eq:virasoro-ope} has poles of order at
most~$4$, and all higher $\Ainf$ operations are
\emph{derived} from the collision residues of the poles
of orders $1$, $2$, and~$4$ via the $\FM$ compactification.
The degree-$4$ pole is the single datum beyond the Lie bracket,
and the shadow depth records this.
\end{verification}


%%%%%%%%%%%%%%%%%%%%%%%%%%%%%%%%%%%%%%%%%%%%%%%%%%%%%%%%%%%%%%%%%%%%%%%%%%%%%
\section{The partially modular Swiss-cheese operad}
\label{sec:partially-modular-sc}
%%%%%%%%%%%%%%%%%%%%%%%%%%%%%%%%%%%%%%%%%%%%%%%%%%%%%%%%%%%%%%%%%%%%%%%%%%%%%

\subsection{Partially modular two-coloured operads}
\label{subsec:partially-modular-def}

\begin{definition}[Partially modular two-coloured operad]
\label{def:partially-modular}
\index{operad!partially modular}
A \emph{partially modular two-coloured operad}
$\cP = (\cP_{\mathrm{cl}}, \cP_{\mathrm{op}}, \cP_{\mathrm{mix}})$
over a symmetric monoidal category $(\cV, \otimes, \mathbf{1})$
consists of the following data, subject to the axioms below.

\medskip\noindent\textbf{Data.}
\begin{enumerate}[label=\textup{(D\arabic*)}]
\item \emph{Closed-colour spaces.}
 For each $g \ge 0$ and $n \ge 0$ with $2g - 2 + n > 0$,
 an object $\cP_{\mathrm{cl}}(g,n) \in \cV$ carrying a right
 action of the symmetric group $\Sgp_n$.

\item \emph{Open-colour spaces.}
 For each $m \ge 0$, an object
 $\cP_{\mathrm{op}}(m) \in \cV$.
 (No genus parameter; no symmetric group action;
 the open inputs are \emph{ordered}.)

\item \emph{Mixed-colour spaces.}
 For each $g \ge 0$, $k \ge 0$, $m \ge 1$,
 an object $\cP_{\mathrm{mix}}(g,k,m) \in \cV$
 with $k$ closed inputs, $m$ open inputs, and one open output.
 The closed inputs carry a right $\Sgp_k$-action;
 the open inputs are ordered (no symmetric group action).

\item \emph{Units.}
 A closed unit $\eta_{\mathrm{cl}} \colon \mathbf{1} \to \cP_{\mathrm{cl}}(0,1)$
 and an open unit $\eta_{\mathrm{op}} \colon \mathbf{1} \to \cP_{\mathrm{op}}(1)$.
\end{enumerate}

\medskip\noindent\textbf{Composition maps.}
\begin{enumerate}[label=\textup{(C\arabic*)}]
\item \emph{Closed--closed composition.}
 For each stable graph $\Gamma$ of genus~$g$ with $n$ external legs,
 a map
 \[
 \gamma_\Gamma^{\mathrm{cl}}
 \colon
 \bigotimes_{v \in V(\Gamma)} \cP_{\mathrm{cl}}(g(v), n(v))
 \longrightarrow
 \cP_{\mathrm{cl}}(g, n),
 \]
 where $V(\Gamma)$ is the vertex set, $g(v)$ and $n(v)$ are the
 genus and valence of vertex~$v$, and
 $g = \sum_v g(v) + b_1(\Gamma)$ is the total genus.
 In particular, this includes the non-separating clutching
 (self-edge at a single vertex, raising genus by~$1$):
 \[
 \xi_{\mathrm{nsep}}
 \colon \cP_{\mathrm{cl}}(g, n+2)
 \longrightarrow \cP_{\mathrm{cl}}(g+1, n).
 \]

\item \emph{Open--open composition.}
 For each planar tree $T$ with $m$ leaves and one root, a map
 \[
 \gamma_T^{\mathrm{op}}
 \colon
 \bigotimes_{v \in V(T)} \cP_{\mathrm{op}}(m(v))
 \longrightarrow
 \cP_{\mathrm{op}}(m),
 \]
 where $m(v)$ is the number of incoming edges at vertex~$v$.
 (Trees are planar: compositions respect the ordering of open inputs.)

\item \emph{Closed-into-mixed composition.}
 For a closed operation inserted into a closed input slot of a mixed
 operation:
 \[
 \gamma^{\mathrm{cl} \to \mathrm{mix}}_i
 \colon
 \cP_{\mathrm{mix}}(g', k, m) \otimes \cP_{\mathrm{cl}}(g'', n)
 \longrightarrow
 \cP_{\mathrm{mix}}(g' + g'', k + n - 1, m),
 \]
 substituting the closed operation into the $i$-th closed input
 ($1 \le i \le k$), and summing genera.

\item \emph{Open-into-mixed composition.}
 For an open operation inserted into an open input slot of a mixed
 operation:
 \[
 \gamma^{\mathrm{op} \to \mathrm{mix}}_j
 \colon
 \cP_{\mathrm{mix}}(g, k, m) \otimes \cP_{\mathrm{op}}(m')
 \longrightarrow
 \cP_{\mathrm{mix}}(g, k, m + m' - 1),
 \]
 substituting the open operation into the $j$-th open input
 ($1 \le j \le m$), preserving the ordering.

\item \emph{Mixed-into-open composition.}
 For a mixed operation (with open output) inserted into an
 open input slot of an open operation:
 \[
 \gamma^{\mathrm{mix} \to \mathrm{op}}_j
 \colon
 \cP_{\mathrm{op}}(m) \otimes \cP_{\mathrm{mix}}(g, k, m')
 \longrightarrow
 \cP_{\mathrm{mix}}(g, k, m + m' - 1).
 \]

\item \emph{Mixed-into-mixed composition.}
 For a mixed operation inserted into an open input slot of
 another mixed operation:
 \[
 \gamma^{\mathrm{mix} \to \mathrm{mix}}_j
 \colon
 \cP_{\mathrm{mix}}(g', k', m')
 \otimes \cP_{\mathrm{mix}}(g'', k'', m'')
 \longrightarrow
 \cP_{\mathrm{mix}}(g' + g'', k' + k'', m' + m'' - 1),
 \]
 substituting the second into the $j$-th open slot of the first.

\item \emph{Non-separating clutching on mixed operations.}
 Contracting two closed inputs of a mixed operation raises
 the genus:
 \[
 \xi^{\mathrm{mix}}_{\mathrm{nsep}}
 \colon
 \cP_{\mathrm{mix}}(g, k+2, m)
 \longrightarrow
 \cP_{\mathrm{mix}}(g+1, k, m).
 \]
\end{enumerate}

\medskip\noindent\textbf{Axioms.}
\begin{enumerate}[label=\textup{(A\arabic*)}]
\item \emph{Equivariance.}
 The closed--closed compositions $\gamma_\Gamma^{\mathrm{cl}}$
 are $\Sgp_n$-equivariant in the sense of
 Getzler--Kapranov~\cite{GK98}.
 The closed-into-mixed compositions $\gamma^{\mathrm{cl} \to \mathrm{mix}}_i$
 are $\Sgp_k$-equivariant on the closed inputs and preserve the
 ordering on the open inputs.
 The open--open compositions $\gamma_T^{\mathrm{op}}$ respect
 the planar (ordered) structure of the tree.

\item \emph{Associativity.}
 All iterated compositions agree:
 \begin{enumerate}[label=\textup{(\alph*)}]
 \item Closed--closed: the compositions indexed by stable graphs
 satisfy the Getzler--Kapranov associativity: composition
 along a subgraph $\Gamma' \subset \Gamma$ followed by
 composition along $\Gamma / \Gamma'$ equals composition
 along~$\Gamma$.
 \item Open--open: for nested planar trees $T' \subset T$,
 $\gamma_{T}^{\mathrm{op}}
 = \gamma_{T/T'}^{\mathrm{op}} \circ
 (\id \otimes \gamma_{T'}^{\mathrm{op}})$.
 \item Mixed compositions are associative: inserting a
 closed operation into a mixed operation that is then composed
 with another yields the same result as composing first and
 inserting second (whenever both routes are defined).
 Formally, the diagrams
 expressing associativity of (C3)--(C7) commute.
 \end{enumerate}

\item \emph{Unitality.}
 The closed unit acts as the identity:
 $\gamma_i^{\mathrm{cl} \to \mathrm{mix}}(\mu, \eta_{\mathrm{cl}})
 = \mu$ for all $\mu \in \cP_{\mathrm{mix}}(g,k,m)$ and
 $1 \le i \le k$;
 $\gamma_\Gamma^{\mathrm{cl}}(\ldots, \eta_{\mathrm{cl}}, \ldots)$
 acts as the identity on the corresponding edge.
 The open unit acts as the identity:
 $\gamma_j^{\mathrm{op} \to \mathrm{mix}}(\mu, \eta_{\mathrm{op}})
 = \mu$ and
 $\gamma_T^{\mathrm{op}}(\ldots, \eta_{\mathrm{op}}, \ldots)$
 acts as the identity.

\item \emph{Genus-raising compatibility.}
 The non-separating clutching on mixed operations (C7) is
 compatible with the closed-colour clutching (C1):
 if a mixed operation $\mu$ has a closed suboperation
 $\alpha \in \cP_{\mathrm{cl}}(g'', n)$
 plugged into it via (C3), and
 the clutching $\xi_{\mathrm{nsep}}$ is applied to two
 inputs of~$\alpha$, then the result agrees with first
 clutching~$\alpha$ to obtain
 $\xi_{\mathrm{nsep}}(\alpha) \in \cP_{\mathrm{cl}}(g''+1, n-2)$
 and then substituting. In symbols:
 \[
 \xi^{\mathrm{mix}}_{\mathrm{nsep}} \circ
 \gamma^{\mathrm{cl} \to \mathrm{mix}}_i
 \;=\;
 \gamma^{\mathrm{cl} \to \mathrm{mix}}_i \circ
 (\id \otimes \xi_{\mathrm{nsep}}).
 \]

\item \emph{Directionality.}
 The category of operations with at least one open input
 and closed output is the zero object:
 \[
 \cP(\ldots, \mathrm{op}, \ldots; \mathrm{cl}) = 0.
 \]
 No composition (C1)--(C7) can produce an operation with open
 inputs and closed output.
 This is not merely an absence of data: it is required that
 no sequence of compositions circumvents the constraint.
\end{enumerate}
The mixed operations carry a genus label $g$ inherited from the
closed factor; the open colour carries no genus. This asymmetry
reflects the geometry: the holomorphic direction supports curves of
all genera, while the topological interval has no genus.
\end{definition}

\begin{remark}[Relation to Vallette's homotopy operads]
\label{rem:vallette-comparison}
\index{operad!partially modular!comparison with Vallette}
Loday--Vallette~\cite{LV12} develop a comprehensive theory of
bar-cobar duality for (coloured) operads, including a notion of
``homotopy operad'' (op.~cit., Ch.~10 and~13). Their framework
treats \emph{tree-level} operads; the composition maps are
indexed by trees, not by stable graphs of positive genus.
Our partially modular operads differ in two essential respects:
\begin{enumerate}[label=\textup{(\roman*)}]
\item The closed colour carries \emph{modular} composition maps
 indexed by stable graphs $\Gamma \in \StGraph(g,n)$ with
 arbitrary genus, including the non-separating clutching
 $\xi_{\mathrm{nsep}}$ that raises genus. This structure is
 governed by Getzler--Kapranov's theory of modular
 operads~\cite{GK98}, not by the tree-level framework of~\cite{LV12}.
\item The directionality axiom~(A5) has no analogue in the
 symmetric-operad setting: it is specific to the two-coloured
 geometry of Swiss-cheese operads, where the bulk--boundary
 asymmetry forbids open-to-closed compositions.
\end{enumerate}
In the genus-$0$ reduction $\SCmod\big|_{g=0} = \SCchtop$,
the closed colour becomes tree-level and the partially modular
structure collapses to an ordinary two-coloured operad. In
this case, the Loday--Vallette bar-cobar theory applies directly,
and our homotopy-Koszulity theorem
(Theorem~\ref{thm:modular-hkoszul-SC}) specialises to their
framework. At genus $g \ge 1$, the modular structure goes
genuinely beyond the tree-level theory: the non-separating
clutching generates new relations (the genus-$g$ Arnold
relation~\eqref{eq:genus-g-arnold}) that have no tree-level
analogue.
\end{remark}


\subsection{The modular chiral Swiss-cheese operad}
\label{subsec:sc-mod-def}

\begin{definition}[The modular chiral Swiss-cheese operad $\SCmod$]
\label{def:SC-mod}
\index{Swiss-cheese operad!modular|textbf}
The \emph{modular chiral Swiss-cheese operad} $\SCmod$ is the
partially modular two-coloured operad with:
\begin{enumerate}[label=\textup{(\roman*)}]
\item \emph{Closed colour (modular):}
 $\SCmodcl(g,k) = C_\ast(\FM_k(\Sigma_g))$ for all
 $g \ge 0$, $k \ge 0$ with $2g - 2 + k > 0$, where $\FM_k(\Sigma_g)$
 denotes the Fulton--MacPherson compactification of the configuration
 space of $k$ points on a smooth genus-$g$ surface~$\Sigma_g$.

\item \emph{Open colour (genus~$0$):}
 $\SCmodop(m) = \Eone(m)$, the chains on the
 Stasheff associahedron (equivalently, the little intervals
 operad).

\item \emph{Mixed operations:}
 $\SCmodmix(g, k, m)
 = C_\ast(\FM^{\mathrm{HT}}_{k|m}(\Sigma_g, I))$
 for $k$ closed (bulk) insertions on the holomorphic surface
 $\Sigma_g$ and $m$ open insertions along a chosen topological
 interval $I\subset \mathbb R$ in the boundary screen
 $\Sigma_g\times\{0\}\times I$ of the HT half-space.  This is the
 formal-collision screen of the mixed Ran geometry, not a claim that
 the closed surface $\Sigma_g$ has a boundary.

\item \emph{Directionality:}
 $\SCmod(\ldots, \mathrm{top}, \ldots; \mathrm{ch})
 = \varnothing$: no open inputs produce closed outputs.

\item \emph{Genus-$0$ reduction:}
 At $g = 0$, $\FM_k(\Sigma_0) = \FM_k(\C)$, and
 $\SCmod\big|_{g=0} = \SCchtop$
 (Definition~\ref{def:SC-operations}).
\end{enumerate}
\end{definition}

\begin{proposition}[Associated graded of mixed operations]
\label{prop:mixed-product-decomposition}
\index{Swiss-cheese operad!mixed operations!associated graded}
Filter $\FM^{\mathrm{HT}}_{k|m}(\Sigma_g,I)$ by mixed-bubbling
depth. The depth-zero associated graded piece satisfies
\begin{equation}\label{eq:mixed-product}
 \gr^0_{\mathrm{mix}} C_\ast(\FM^{\mathrm{HT}}_{k|m}(\Sigma_g,I))
 \cong
 C_\ast(\FM_k(\Sigma_g)) \otimes C_\ast(\Eone(m)).
\end{equation}
The full mixed complex also contains positive-depth Type~III
mixed bubbling faces; their boundary contribution is the mixed
differential $d_{\mathrm{mix}}$. Thus
\eqref{eq:mixed-product} is an associated-graded product, not a
homotopy equivalence from the full open--closed compactification
to a product.
\end{proposition}

\begin{proof}
The space $\FM^{\mathrm{HT}}_{k|m}(\Sigma_g,I)$ parametrizes the
formal-collision screen of $k$ holomorphic points on~$\Sigma_g$
and $m$ ordered topological points on a boundary interval
$I \cong [0,1]\subset\mathbb R$ in the HT half-space
$\Sigma_g\times\mathbb R_{\ge0}$.
We identify the depth-zero associated graded by forgetting the
positive-depth mixed faces.

\medskip\noindent
\emph{Step~1: The projection.}
Define $\pi \colon \FM^{\mathrm{HT}}_{k|m}(\Sigma_g,I)
\to \FM_k(\Sigma_g)$ by forgetting the boundary points.
Define $\rho \colon \FM^{\mathrm{HT}}_{k|m}(\Sigma_g,I)
\to \Eone(m)$ by forgetting the interior points and
retaining only the ordered boundary configuration.
The pair $(\pi, \rho)$ defines a continuous map
\[
 (\pi, \rho) \colon \FM^{\mathrm{HT}}_{k|m}(\Sigma_g,I)
 \longrightarrow \FM_k(\Sigma_g) \times \Eone(m).
\]

\medskip\noindent
\emph{Step~2: Geometric input.}
The open interval~$I$ lies in the topological boundary direction of
the HT half-space, while the closed insertions live on the
holomorphic surface~$\Sigma_g$. In the mixed FM screen, the strata of
$\FM^{\mathrm{HT}}_{k|m}(\Sigma_g,I)$ correspond to
nested screens. Interior collision screens are parametrized by
$\FM$-data on~$\Sigma_g$; boundary collision screens are
parametrized by $\FM$-data on~$I \cong \R$
(i.e., Stasheff associahedra).
A mixed screen, in which interior and boundary points meet at a
boundary collision face, is a genuine open--closed boundary face.
These faces are precisely the positive-depth part of the
mixed-bubbling filtration. After passing to depth zero, only the
pure interior and pure boundary screens remain, and their face
poset is the product of the face posets of $\FM_k(\Sigma_g)$ and
$\Eone(m)$.

\medskip\noindent
\emph{Step~3: The associated-graded product.}
The cellular chain complex of the depth-zero stratum is therefore
the tensor product of the pure interior and pure boundary chain
complexes, giving \eqref{eq:mixed-product}. The full differential
has one additional component: the boundary of a positive-depth
mixed face. We denote this component by $d_{\mathrm{mix}}$. It is
not killed by a collar contraction; it is the local
bulk-to-boundary operation of the Swiss-cheese structure.

\medskip\noindent
\emph{Step~4: Compatibility with operad structure.}
The associated-graded product~\eqref{eq:mixed-product} is
compatible with the depth-zero parts of the composition maps of
Definition~\ref{def:partially-modular}: closed--closed composition acts on the
$C_\ast(\FM_k(\Sigma_g))$ factor (via the FM composition maps),
open--open composition acts on the $C_\ast(\Eone(m))$ factor
(via the associahedron composition maps), and the mixed
compositions are the perturbations recorded by $d_{\mathrm{mix}}$.
\end{proof}

\begin{remark}[Product shadow versus mixed faces]
\label{rem:product-strictness}
The product in Proposition~\ref{prop:mixed-product-decomposition}
is the depth-zero shadow. It is useful for first approximations,
including the lowest-page genus-filtration computations, but it is
not the full mixed Swiss-cheese compactification. The full object
contains mixed faces, and any bar-cobar or Feynman-transform
argument using the product shadow must either work on associated
graded or explicitly control the $d_{\mathrm{mix}}$ perturbation.
\end{remark}

\begin{remark}[HT mixed FM screen and associated graded]
\label{rem:FM-boundary}
\index{Fulton--MacPherson compactification!manifolds with boundary}
The associated-graded product
$\gr^0_{\mathrm{mix}}\FM^{\mathrm{HT}}_{k|m}(\Sigma_g,I)
\simeq \FM_k(\Sigma_g) \times \Eone(m)$
(Proposition~\ref{prop:mixed-product-decomposition}) is the depth-zero
shadow of the mixed HT compactification. It does not identify the full
mixed open--closed space with a product. We record the geometric input,
following the Fulton--MacPherson compactification for manifolds with
boundary as in Sinha~\cite{Sinha2004} (Remark~5.1).

\medskip\noindent
\emph{(a) Definition of $\FM_k(M)$ for $M$ with boundary.}
Let $M$ be a smooth manifold, possibly with boundary
$\partial M$. The FM compactification $\FM_k(M)$ is the
real oriented blow-up of the iterated diagonals in $M^k$:
one blows up successively the thin diagonal
$\{(x, \ldots, x) : x \in M\} \subset M^k$, then the
pairwise diagonals $\{x_i = x_j\}$, then the triple
diagonals, etc., in order of increasing codimension. For
$M$ without boundary, this is the construction of
Fulton--MacPherson~\cite{FM94} (algebraic version) and
Axelrod--Singer (smooth version). For $M$ with
boundary, the blow-up procedure is identical, but the
resulting space $\FM_k(M)$ acquires two types of boundary
strata:
\begin{itemize}
\item \emph{Interior collision strata}: arising from two or
 more points in $M \setminus \partial M$ approaching the same
 limit. These are parametrised by $\FM$-data in the interior,
 exactly as for manifolds without boundary.
\item \emph{Boundary collision strata}: arising from two or
 more points on $\partial M$ approaching the same limit.
 These are parametrised by $\FM$-data on $\partial M$.
\end{itemize}

\medskip\noindent
\emph{(b) Collision types for
$M = \Sigma_g \times [0,1]$.}
For $M = \Sigma_g \times [0,1]$, a manifold with boundary
$\partial M = \Sigma_g \times \{0\} \sqcup \Sigma_g \times \{1\}$,
pure interior and pure boundary collisions are geometrically
separated, but there are also mixed boundary faces.
A point in the interior has the form $(z, t)$ with
$z \in \Sigma_g$, $t \in (0,1)$; a point on the boundary has
$t = 0$ or $t = 1$. The collar neighbourhood
$\Sigma_g \times [0, \epsilon)$ provides a buffer of width
$\epsilon > 0$ between boundary points (at $t = 0$) and any
interior point (at $t \ge \epsilon$). As two interior points
approach a common limit, they collide in the interior; as two
boundary points approach a common limit, they collide on the
boundary. A mixed open--closed bubbling face records the limiting
normal approach of an interior point to a boundary collision
configuration together with the tangential boundary configuration.
These Type~III faces are the geometric source of the mixed
differential $d_{\mathrm{mix}}$.

\medskip\noindent
\emph{(c) The associated-graded product.}
Removing the positive-depth mixed faces leaves the product
stratum
\[
 \gr^0_{\mathrm{mix}}\FM^{\mathrm{HT}}_{k|m}(\Sigma_g,I)
 \;\cong\;
 \FM_k(\Sigma_g) \times \FM_m([0,1]),
\]
where $\FM_k(\Sigma_g)$ compactifies the interior collisions
and $\FM_m([0,1])$ compactifies the boundary collisions.
The space $\FM_m([0,1])$ is the Stasheff associahedron $K_m$,
which is the cell model for the $\Eone$-operad:
$\FM_m([0,1]) = K_m = \Eone(m)$. The full compactification is
the product shadow plus the mixed faces. Those faces are
load-bearing: they are the local bulk-to-boundary operations of
the Swiss-cheese structure.
\end{remark}

\begin{remark}[Why ``partially modular'']
\label{rem:partially-modular}
The closed colour of $\SCmod$ is a full modular operad: it
carries non-separating clutching maps (contracting an edge to
raise genus) and separating clutching maps (gluing along marked
points), indexed by stable graphs of all genera. The open colour
is \emph{not} modular: the interval $\R$ has no genus, so there
is no genus-raising operation on the open colour. The mixed
operations inherit a genus label from the closed factor but have
no independent genus of their own. This asymmetry is the
operadic shadow of the geometric fact that the HT manifold
$\Sigma_g \times \R$ has genus only in the holomorphic direction.
\end{remark}

\begin{construction}[The extraction functor $\mathrm{Loc}$]
\label{con:extraction-functor}
\index{factorization algebra!extraction of local operad|textbf}
\index{Loc functor@$\mathrm{Loc}$ functor|textbf}
Let $\cF$ be a BD factorization Swiss-cheese algebra on the mixed
open--closed holomorphic--topological Ran geometry
$\Ran^{\mathrm{oc}}_{\mathrm{HT}}(\Sigma_g,I)$
(Chapter~\ref{ch:factorization-swiss-cheese}). We construct an
$\SCmod$-coalgebra $\mathrm{Loc}(\cF)$ as follows.

\medskip\noindent
\emph{Step~1: Formal completion along infinitesimal neighbourhoods.}
For each point $x \in \Sigma_g$, let $D_x^{(n)}$ denote the
$n$-th infinitesimal neighbourhood of~$x$ in~$\Sigma_g$
($D_x^{(n)} = \operatorname{Spec} \cO_{\Sigma_g,x}/\mathfrak{m}_x^{n+1}$).
The factorization $\cD$-module $\cF$ restricts to each
$D_x^{(n)}$, and we define the formal completion
\[
 \hat{\cF}_x
 \;:=\;
 \varprojlim_n \, \cF\big|_{D_x^{(n)}}
\]
as a pro-object in the category of chain complexes. The
factorization structure of~$\cF$ (the chiral product
on colliding points) restricts to the formal disc,
giving operations
\[
 \mu_k^{\mathrm{loc}} \colon
 \hat{\cF}_x^{\otimes k}
 \longrightarrow
 \hat{\cF}_x
\]
for each $k \ge 2$, parametrised by collision data in the
formal disc $\hat{D}_x \cong \operatorname{Spf} \C\llbracket z \rrbracket$.

\medskip\noindent
\emph{Step~2: Identification with FM chains.}
The formal-completion data determines chains on
$\FM_k(\C)$ by the $\FM$-to-$\Ran$ comparison: the
Fulton--MacPherson compactification $\FM_k(\C)$ is the
real-oriented blow-up of the diagonals in $\C^k$, and
its strata record precisely the collision patterns that
the formal-disc operations $\mu_k^{\mathrm{loc}}$ detect.
Concretely, for a collection
$(z_1, \ldots, z_k) \in \FM_k(\C)$ approaching a
boundary stratum $\partial_S \FM_k(\C)$
(where $S \subset \{1, \ldots, k\}$ is a colliding subset),
the formal-completion data provides the OPE residue
along that stratum, yielding a chain-level map
\[
 C_\ast(\FM_k(\C))
 \longrightarrow
 \Hom(\hat{\cF}_x^{\otimes k},\, \hat{\cF}_x).
\]
This identification uses the comparison between the
Ran space model and the FM compactification model for
factorisation algebras
(cf.~Ayala--Francis~\cite{AF15} for the general
comparison; in the chiral setting,
$\Ran^{\le k}(\hat{D}_x)$ and $\FM_k(\C)$ give
equivalent chain-level data over a field of
characteristic zero).

\medskip\noindent
\emph{Step~3: The open colour is unchanged.}
The $\Eone$-structure on the open colour (the topological
direction $\R$) is locally constant: the factorization
structure on $\Ran(\R)$ is an $\Eone$-algebra, and
restriction to any subinterval $I \subset \R$ yields the
same $\Eone$-algebra up to canonical isomorphism. Hence
the open-colour data passes through $\mathrm{Loc}$
unchanged.

\medskip\noindent
\emph{Assembly.} The closed-colour data from Step~2 and
the open-colour data from Step~3, together with the mixed
operations (formal completion applied to both bulk and
boundary insertions simultaneously, using the associated-graded
product and the mixed differential of
Proposition~\ref{prop:mixed-product-decomposition}),
assemble into an $\SCmod$-coalgebra $\mathrm{Loc}(\cF)$.
The coalgebra structure on the closed colour is the bar-dual
of the operadic algebra structure: the coproduct is
ordered deconcatenation on the $\Eone$ factor, and the
differential is the OPE residue on the
$C_\ast(\FM_k(\C))$ factor.
\end{construction}

\begin{proposition}[Formal-collision extraction;
\ClaimStatusProvedHere]
\label{prop:extraction-functor}
\index{Loc functor@$\mathrm{Loc}$ functor!properties}
The extraction $\mathrm{Loc}$ is a formal-collision assignment
\[
 \mathrm{Loc} \colon
 \mathrm{FactAlg}_{\mathrm{SC}}(\Ran^{\mathrm{oc}}_{\mathrm{HT}}(\Sigma_g,I))
 \longrightarrow {\SCmod}\text{-}\mathrm{CoAlg}.
\]
It records formal OPE residues, the locally constant open
$\Eone$ data, and the associated mixed collision maps. At
genus~$0$, on translation-invariant vacuum-generated data on the
formal disc, this is an equivalence onto its essential image. No
full faithfulness or essential-surjectivity statement is asserted
for factorization algebras on $\mathbb{P}^1$ or on
$\Sigma_g$ as global curves. At $g\ge 1$, periods, monodromy,
Arakelov corrections, and the curvature class
$\kappa(\cA)\cdot\omega_g$ are outside the formal-collision
skeleton.
\end{proposition}

\begin{proof}
\emph{Functoriality.}
A morphism $\phi \colon \cF \to \cG$ of BD factorization
Swiss-cheese algebras is a natural transformation of the
underlying factorizable $\cD$-modules, compatible with the
$\Eone$-structure on the open colour. Restricting to formal
neighbourhoods and applying the FM comparison preserves
the chain-level maps, so
$\mathrm{Loc}(\phi) \colon \mathrm{Loc}(\cF)
\to \mathrm{Loc}(\cG)$
is a well-defined morphism of $\SCmod$-coalgebras.

\medskip\noindent
\emph{Genus~$0$ local essential image.}
On the formal disc, translation-invariant vacuum-generated
factorization data is determined by its OPE residues. The
FM-to-Ran comparison identifies those residues with the local
$\SCchtop$ operations, so $\mathrm{Loc}$ is an equivalence onto
this formal-disc essential image. This does not reconstruct
arbitrary factorization $D$-modules on the global curve
$\mathbb{P}^1$.

\medskip\noindent
\emph{Positive genus.}
The Arakelov Green's function $G_{\mathrm{Ar}}(z,w)$ on
$\Sigma_g$ (the propagator of the curved geometric model)
has the local expansion
$G_{\mathrm{Ar}}(z,w) = \log|z - w|^2 + O(1)$ near the diagonal,
which matches the formal-disc propagator
$G_0(z,w) = \log|z-w|^2$. The global correction
\[
 G_{\mathrm{Ar}}(z,w) - G_0(z,w)
 = -2\pi \operatorname{Im}
 \Bigl(\int_w^z \boldsymbol{\omega}\Bigr)^T
 \cdot (\operatorname{Im} \boldsymbol{\tau})^{-1}
 \cdot \operatorname{Im}
 \Bigl(\int_w^z \boldsymbol{\omega}\Bigr)
 + C_g
\]
(where $\boldsymbol{\omega} = (\omega_1, \ldots, \omega_g)$
is the basis of holomorphic differentials and
$\boldsymbol{\tau}$ is the period matrix) is a \emph{global}
datum involving periods and the Hodge structure of~$\Sigma_g$.
This correction is a smooth function on
$\Sigma_g \times \Sigma_g$ (no singularity on the diagonal),
so it disappears only after formal completion followed by projection
to the singular/OPE principal part. Its Taylor expansion does not
vanish as a formal function. But it is
precisely this correction that produces the curvature:
applying $\dbar_z \dbar_w$ to the global correction gives
\[
 \dbar_z \dbar_w \bigl(G_{\mathrm{Ar}} - G_0\bigr)
 = -\frac{1}{g}\, \mu_{\mathrm{Ar}}(z) \wedge \mu_{\mathrm{Ar}}(w),
\]
where $\mu_{\mathrm{Ar}} = \frac{i}{2}
\sum_{j=1}^g \omega_j \wedge \bar{\omega}_j / \int_{\Sigma_g}
\omega_j \wedge \bar{\omega}_j$ is the Arakelov $(1,1)$-form.
Integrating this over the surface produces the curvature class
$\kappa(\cA) \cdot \omega_g \in H^2(\overline{\cM}_g)$.
Formal completion detects the singular OPE part but not this
global smooth correction. Thus the construction gives a local
shadow functor; it is not claimed to be faithful, full, or
essentially surjective on positive-genus factorization data.
\end{proof}


%%%%%%%%%%%%%%%%%%%%%%%%%%%%%%%%%%%%%%%%%%%%%%%%%%%%%%%%%%%%%%%%%%%%%%%%%%%%%
\section{Modular homotopy-Koszulity}
\label{sec:modular-homotopy-koszulity}
%%%%%%%%%%%%%%%%%%%%%%%%%%%%%%%%%%%%%%%%%%%%%%%%%%%%%%%%%%%%%%%%%%%%%%%%%%%%%

\subsection{Definition}
\label{subsec:mod-hkoszul-def}

\begin{definition}[Modular homotopy-Koszulity for partially modular operads]
\label{def:modular-hkoszul}
\index{homotopy-Koszulity!modular|textbf}
A partially modular two-coloured operad $\cP$ is \emph{modular
homotopy-Koszul} if the canonical counit
\[
 \epsilon \colon \Omega\mathbf{B}(\cP) \xrightarrow{\;\sim\;} \cP
\]
is a quasi-isomorphism of partially modular dg operads, where the bar
and cobar constructions are taken in the category of partially modular
operads (closed colour uses stable-graph bar; open colour uses
tree-level bar; mixed operations use the associated-graded product
together with the mixed-face perturbation).
Equivalently: the Feynman transform of the closed colour is homotopy
involutive ($\FTmod^2 \simeq \id$ on closed-colour algebras), the
open-colour bar-cobar is a Quillen equivalence, and the mixed bar-cobar
is compatible with both.
\end{definition}


\subsection{The theorem}
\label{subsec:mod-hkoszul-thm}

\begin{theorem}[Modular homotopy-Koszulity of $\SCmod$ under
formality and mixed-face convergence; \ClaimStatusConditional]
\label{thm:modular-hkoszul-SC}
\index{Swiss-cheese operad!modular homotopy-Koszulity|textbf}
The genus-$0$ formal-collision specialization of $\SCmod$ is
homotopy-Koszul by Theorem~\ref{thm:homotopy-Koszul}. The
all-genus modular statement for $\SCmod$ is conditional on:
\begin{enumerate}[label=\textup{(\roman*)}]
\item global genus-$g$ operadic formality for the closed-colour
 operation spaces, compatible with stable-graph clutching;
\item control of the mixed-face perturbation $d_{\mathrm{mix}}$
 in the two-coloured bar-cobar construction;
\item convergence of the completed modular genus filtration.
\end{enumerate}
Under these hypotheses, the counit
\[
 \epsilon \colon
 \Omega\mathbf{B}(\SCmod) \xrightarrow{\;\sim\;} \SCmod
\]
is a quasi-isomorphism of partially modular two-coloured dg operads.
\end{theorem}

\begin{proof}
The proof below verifies the genus-$0$ and local-screen inputs,
then records the additional positive-genus hypotheses under which
the standard modular bar-cobar argument applies.

\medskip\noindent
\textbf{Step~1: Genus-$0$ homotopy-Koszulity (established).}
At genus~$0$, the closed colour of $\SCmod$ is
$C_\ast(\FM_k(\C))$ and the full genus-$0$ operad is $\SCchtop$.
Theorem~\ref{thm:homotopy-Koszul} proves that $\SCchtop$ is
homotopy-Koszul, assembling three inputs: Koszulity of the
classical Swiss-cheese operad $\mathsf{SC}$~\cite{Hoe09,HL12},
Kontsevich formality of $C_\ast(\FM_\bullet(\C)) \simeq \Etwo$
\cite{Kon03}, and transfer of homotopy-Koszulity through
quasi-isomorphisms of dg operads~\cite{BM03,Val07}.

\medskip\noindent
\textbf{Step~2: Closed-colour modular Koszulity.}
We must show that the closed-colour modular operad
$\SCmodcl = \{C_\ast(\FM_k(\Sigma_g))\}_{g,k}$
is modular homotopy-Koszul in the sense of
Getzler--Kapranov~\cite{GK98}. This requires three
ingredients.

\medskip
\emph{(a) Fay trisecant implies the Arnold relation at genus~$g$.}
We give the complete derivation in five steps: the Fay identity,
its differentiation, the three-point specialisation, the wedge
with antisymmetric form factors, and the genus-$0$ and genus-$1$
verifications.

\medskip
\emph{Step~(a.i): The Fay trisecant identity.}
Let $E(x,y)$ denote the prime form on $\Sigma_g$, the
unique $(-\tfrac{1}{2}, -\tfrac{1}{2})$-differential on
$\Sigma_g \times \Sigma_g$ with a simple zero on the diagonal
and no other zeros. Let $\theta$ denote the Riemann theta
function with characteristics associated to a choice of
odd spin structure on $\Sigma_g$, and let $\Delta$ denote
the Riemann divisor class. The Fay trisecant identity
(Fay~\cite{Fay73}, Theorem~1) states: with integration paths
specified via the Abel--Jacobi map
$\cA \colon \Sigma_g \to \operatorname{Jac}(\Sigma_g)$,
\begin{multline}\label{eq:fay-trisecant}
 E(x,y)\, E(z,w)\;
 \theta\!\bigl(\cA(x) - \cA(y)
 + \cA(z) - \cA(w) + \Delta\bigr) \\
 = E(x,z)\, E(y,w)\;
 \theta\!\bigl(\cA(x) - \cA(w)
 + \cA(z) - \cA(y) + \Delta\bigr) \\
 \quad{}-\; E(x,w)\, E(y,z)\;
 \theta\!\bigl(\cA(x) - \cA(z)
 + \cA(w) - \cA(y) + \Delta\bigr).
\end{multline}
This is an identity of sections of a line bundle on
$\Sigma_g^4$.

\medskip
\emph{Step~(a.ii): Logarithmic differentiation.}
Define the holomorphic propagator
\begin{equation}\label{eq:hol-propagator}
 f(x,y) := \partial_x \log E(x,y).
\end{equation}
Since $E(x,y)$ has a simple zero at $x = y$ with unit
leading coefficient in any local coordinate, $f(x,y)$ has
a simple pole at $x = y$ with residue~$1$:
\[
 f(x,y) = \frac{1}{x - y} + r(x,y),
\]
where $r(x,y)$ is holomorphic on $\Sigma_g \times \Sigma_g$.
Apply $\partial_x \log$ to both sides
of~\eqref{eq:fay-trisecant}. On the left side,
$\partial_x \log(E(x,y)\, E(z,w)\, \theta(\cdots))
= f(x,y) + \partial_x \log \theta(\cdots)$,
since $\partial_x \log E(z,w) = 0$.
On the right side, the logarithmic derivative
produces $f(x,z)$ and $f(x,w)$ from the prime forms
$E(x,z)$ and $E(x,w)$, plus $\partial_x \log$ of the
theta functions. The theta-function derivatives on both
sides are \emph{symmetric} functions of the four variables
(by the symmetry properties of the Riemann theta function
and its heat equation), so they contribute a symmetric
remainder upon taking differences.

The identity underlying the Arnold relation is the
\emph{three-term quadratic identity}, obtained by
multiplying the differentiated Fay identity by $f(z,w)$
and symmetrising. After cancellation of the symmetric
theta-function contributions, one obtains:
\begin{equation}\label{eq:fay-quadratic}
 f(x,y)\, f(z,w) + f(y,z)\, f(x,w) + f(z,x)\, f(y,w)
 = S(x,y,z,w),
\end{equation}
where $S(x,y,z,w)$ is a \emph{symmetric function} of its
four arguments, expressible in terms of the Bergman kernel
$B(x,y) = \partial_x \partial_y \log \theta(\cA(x) - \cA(y)
+ \Delta)$ and the holomorphic $1$-forms on~$\Sigma_g$.
The function $S$ is symmetric. We verify this
explicitly. By definition,
\begin{equation}\label{eq:fay-S-explicit}
 S(x,y,z,w)
 = f(x,y)\,f(z,w)
 + f(y,z)\,f(x,w)
 + f(z,x)\,f(y,w),
\end{equation}
where $f(x,y) = \partial_x \log E(x,y)$ as
in~\eqref{eq:hol-propagator}. Symmetry of $S$ in all four
arguments follows from two observations. First, each summand
$f(a,b)\,f(c,d)$ is a product of scalars, so it equals
$f(c,d)\,f(a,b)$: simultaneous exchange of the two pairs
$(a,b) \leftrightarrow (c,d)$ is a symmetry of each term.
Second, the three summands are cyclically related under
$(x,y,z) \mapsto (y,z,x)$ with $w$ held fixed. Thus $S$ is
symmetric in $\{x,y,z\}$ by cyclic invariance, and the pair
exchange shows $S(x,y,z,w) = S(z,w,x,y)$, so $S$ is symmetric
in all four variables.

\smallskip\noindent\emph{Genus-$0$ verification.}
At genus~$0$, $E(x,y) = x - y$ and $f(x,y) = 1/(x - y)$.
Then~\eqref{eq:fay-S-explicit} gives
\[
 S(x,y,z,w)
 = \frac{1}{(x-y)(z-w)}
 + \frac{1}{(y-z)(x-w)}
 + \frac{1}{(z-x)(y-w)}.
\]
This vanishes identically by the partial fraction identity:
clearing the common denominator $(x-y)(y-z)(z-x)(x-w)(y-w)(z-w)$
yields a numerator that cancels termwise, as one verifies by
writing each fraction over the common denominator and applying
the identity $(y-z)(y-w)(z-x) + (x-y)(z-w)(z-x)
+ (x-y)(y-z)(x-w) = 0$.
Thus $S \equiv 0$ at genus~$0$, since
$\mathbb{P}^1$ has no holomorphic $1$-forms.

\smallskip\noindent\emph{Genus-$1$ verification.}
At genus~$1$, $E(x,y) = \vartheta_1(x-y;\tau)/\vartheta_1'(0;\tau)$
and $f(x,y) = \zeta(x-y;\tau)$. With $u = x - y$, $v = y - z$,
equation~\eqref{eq:fay-S-explicit} becomes
\[
 S = \zeta(u)\,\zeta(v') + \zeta(v)\,\zeta(u')
 + \zeta(-u-v)\,\zeta(w'),
\]
where the second arguments are functions of $w$. After
simplification using the Weierstrass addition formula, one obtains
$S = \tfrac{1}{2}[\wp(u) + \wp(v) + \wp(u+v)]$
as in~\eqref{eq:fay-genus1-product}. Since $\wp$ is even
($\wp(-z) = \wp(z)$), this expression is manifestly symmetric
in $(x,y,z)$ (invariant under the permutations
$u \leftrightarrow v$ and $u \leftrightarrow -(u+v)$),
confirming the symmetry of~$S$ at genus~$1$.

\medskip
\emph{Step~(a.iii): Three-point specialisation.}
Set $(x, y, z, w) = (z_1, z_2, z_3, z_0)$ in~%
\eqref{eq:fay-quadratic}, where $z_0$ is an auxiliary
(spectator) point. The identity becomes:
\begin{equation}\label{eq:fay-three-point}
 f(z_1, z_2)\, f(z_3, z_0)
 + f(z_2, z_3)\, f(z_1, z_0)
 + f(z_3, z_1)\, f(z_2, z_0)
 = S(z_1, z_2, z_3, z_0).
\end{equation}
Now specialise $z_0 \to z_3$ (and use the regularity of the
FM compactification to handle the collision).
Since $f(z_3, z_3)$ has a pole, we instead work with the
\emph{product identity} directly: multiply~\eqref{eq:fay-quadratic}
with $(x,y,z) = (z_1, z_2, z_3)$ by the $1$-forms
$dz_1\, dz_2\, dz_3$ to form the $2$-forms
$\omega_{ij} = f(z_i, z_j)\, dz_i$ and their wedge products.
The three products of propagator pairs that appear are:
\begin{align*}
 f(z_1, z_2)\, f(z_2, z_3)\, dz_1 \wedge dz_2, \\
 f(z_2, z_3)\, f(z_3, z_1)\, dz_2 \wedge dz_3, \\
 f(z_3, z_1)\, f(z_1, z_2)\, dz_3 \wedge dz_1.
\end{align*}
Their sum, by~\eqref{eq:fay-quadratic} with $w$ set equal to the
second argument of the first factor in each term and the
four-variable identity applied to all cyclic permutations, yields:
\begin{equation}\label{eq:fay-sum-symmetric}
 \omega_{12} \wedge \omega_{23}
 + \omega_{23} \wedge \omega_{31}
 + \omega_{31} \wedge \omega_{12}
 = \Sigma_{\mathrm{sym}}(z_1, z_2, z_3)\, dz_1 \wedge dz_2 \wedge dz_3,
\end{equation}
where $\Sigma_{\mathrm{sym}}$ is a function that is
\emph{symmetric} in $(z_1, z_2, z_3)$, assembled from the
symmetric remainder $S$ in~\eqref{eq:fay-quadratic}.
Explicitly, the three-variable symmetric remainder is
\begin{equation}\label{eq:fay-three-variable-S}
 \Sigma_{\mathrm{sym}}(z_1, z_2, z_3)
 = f(z_1,z_2)\,f(z_2,z_3)
 + f(z_2,z_3)\,f(z_3,z_1)
 + f(z_3,z_1)\,f(z_1,z_2),
\end{equation}
where $f(x,y) = \partial_x \log E(x,y)$ as
in~\eqref{eq:hol-propagator}. Symmetry in $(z_1, z_2, z_3)$
is immediate: the three summands are cyclically related under
$(z_1, z_2, z_3) \mapsto (z_2, z_3, z_1)$.

\smallskip\noindent
\emph{Genus-$0$ reduction.}
At genus~$0$, $f(x,y) = 1/(x-y)$, and
$\Sigma_{\mathrm{sym}} = 0$ by the partial fraction identity
(the same identity that underlies~\eqref{eq:fay-S-explicit}
at genus~$0$).

\smallskip\noindent
\emph{Genus-$1$ reduction.}
At genus~$1$, $f(x,y) = \zeta(x - y; \tau)$, and
\[
 \Sigma_{\mathrm{sym}}(z_1, z_2, z_3)
 = \tfrac{1}{2}\bigl[\wp(z_1 - z_2)
 + \wp(z_2 - z_3)
 + \wp(z_3 - z_1)\bigr],
\]
which is manifestly symmetric since $\wp$ is even. This
follows from~\eqref{eq:fay-genus1-product} by setting
$u = z_1 - z_2$, $v = z_2 - z_3$, so that
$u + v = z_1 - z_3$.

\medskip
\emph{Step~(a.iv): Wedge with antisymmetric form factor annihilates
the symmetric remainder.}
Let $\omega_B^{(3)}(z_1, z_2, z_3) = c_{abc}\, dz_1\, dz_2\, dz_3$
be the antisymmetric form factor arising from the structure constants
$c_{abc}$ of the chiral algebra (for instance, $c_{abc}$ may be the
structure constants of a Lie algebra $\fg$, totally antisymmetric
by the Jacobi identity and ad-invariance of the Killing form).
The bar-complex differential at degree~$3$ involves the
contraction of the propagator $2$-form with this antisymmetric
tensor. Explicitly, the contribution to $D_0^2$ at degree~$3$
is proportional to
\[
 \bigl(\omega_{12} \wedge \omega_{23}
 + \omega_{23} \wedge \omega_{31}
 + \omega_{31} \wedge \omega_{12}\bigr)
 \otimes c_{abc}.
\]
By~\eqref{eq:fay-sum-symmetric}, the left factor equals
$\Sigma_{\mathrm{sym}}\, dz_1 \wedge dz_2 \wedge dz_3$.
Since $\Sigma_{\mathrm{sym}}$ is symmetric in $(z_1, z_2, z_3)$
and $c_{abc}$ is totally antisymmetric, the contraction
\[
 \sum_{\sigma \in \Sgp_3}
 \Sigma_{\mathrm{sym}}(z_{\sigma(1)}, z_{\sigma(2)}, z_{\sigma(3)})
 \cdot c_{\sigma(a)\sigma(b)\sigma(c)} = 0.
\]
More precisely: on $\Conf_3(\Sigma_g)$, the tensor
$\Sigma_{\mathrm{sym}} \otimes c_{abc}$ transforms as the
product of the trivial $\Sgp_3$-representation (symmetric
function) with the sign representation (antisymmetric tensor),
which is the sign representation. Its invariant projection
onto the trivial representation, which is what the
bar-complex differential computes, vanishes.

Hence the genus-$g$ Arnold $3$-form
\[
 A_3^{\mathrm{hol}}
 := \bigl(\omega_{12} \wedge \omega_{23}
 + \omega_{23} \wedge \omega_{31}
 + \omega_{31} \wedge \omega_{12}\bigr) \otimes c_{abc}
 = 0
\]
in the relevant component of $\Omega^2(\Conf_k(\Sigma_g))
\otimes \cA^{\otimes 3}$, and the Arnold relation holds:
\begin{equation}\label{eq:genus-g-arnold}
 \omega_{ij} \wedge \omega_{jk}
 + \omega_{jk} \wedge \omega_{ki}
 + \omega_{ki} \wedge \omega_{ij} = 0
 \quad \text{in } \Omega^2(\Conf_k(\Sigma_g))
 \text{ contracted with antisymmetric tensors.}
\end{equation}

\medskip
\emph{Step~(a.v): Genus-$0$ and genus-$1$ verifications.}

\noindent\emph{Genus~$0$.}
At genus~$0$, the prime form is $E(x,y) = x - y$, so
$f(x,y) = \partial_x \log(x-y) = (x-y)^{-1}$ and
$\omega_{ij} = d\log(z_i - z_j)$. The Fay
identity~\eqref{eq:fay-quadratic} becomes
\[
 \frac{1}{(z_1 - z_2)(z_3 - w)}
 + \frac{1}{(z_2 - z_3)(z_1 - w)}
 + \frac{1}{(z_3 - z_1)(z_2 - w)} = 0,
\]
which is the partial fraction identity. The symmetric
remainder $S = 0$ identically at genus~$0$ (there are no
holomorphic $1$-forms on $\mathbb{P}^1$), so the Arnold
relation~\eqref{eq:genus-g-arnold} holds without the
antisymmetric-tensor contraction:
$\omega_{12} \wedge \omega_{23}
+ \omega_{23} \wedge \omega_{31}
+ \omega_{31} \wedge \omega_{12} = 0$
as a $2$-form identity on $\Conf_3(\C)$.

\noindent\emph{Genus~$1$.}
At genus~$1$, the surface is an elliptic curve
$\Sigma_1 = \C/(\Z + \tau\Z)$. The prime form is
$E(x,y) = \vartheta_1(x-y; \tau) / \vartheta_1'(0;\tau)$,
where $\vartheta_1$ is the odd Jacobi theta function, and
the holomorphic propagator is
$f(x,y) = \zeta(x - y; \tau)$, the Weierstrass zeta function.
The Fay identity~\eqref{eq:fay-quadratic} specialises to:
\begin{equation}\label{eq:fay-genus1-product}
 \zeta(u)\,\zeta(v)
 + \zeta(v)\,\zeta(-u-v)
 + \zeta(-u-v)\,\zeta(u)
 = \tfrac{1}{2}\bigl[\wp(u) + \wp(v) + \wp(u+v)\bigr],
\end{equation}
where $u = z_1 - z_2$, $v = z_2 - z_3$, and
$\wp(z) = -\zeta'(z)$ is the Weierstrass $\wp$-function.
The right side of~\eqref{eq:fay-genus1-product} is the
symmetric remainder: $\tfrac{1}{2}[\wp(u) + \wp(v) + \wp(u+v)]$
is manifestly symmetric under permutations of
$(z_1, z_2, z_3)$ (since $\wp$ is even: $\wp(-z) = \wp(z)$),
so it vanishes upon contraction with antisymmetric form factors.

To verify~\eqref{eq:fay-genus1-product}: by the Weierstrass
addition formula $\zeta(u+v) = \zeta(u) + \zeta(v) +
\frac{1}{2}\frac{\wp'(u) - \wp'(v)}{\wp(u) - \wp(v)}$, we have
$\zeta(-u-v) = -\zeta(u) - \zeta(v) -
\frac{1}{2}\frac{\wp'(u) - \wp'(v)}{\wp(u) - \wp(v)}$.
Substituting into the left side of~\eqref{eq:fay-genus1-product}:
\begin{align*}
 &\zeta(u)\,\zeta(v)
 + \zeta(v)\bigl[-\zeta(u) - \zeta(v)
 - \tfrac{1}{2}\tfrac{\wp'(u) - \wp'(v)}{\wp(u) - \wp(v)}\bigr] \\
 &\quad{}+ \bigl[-\zeta(u) - \zeta(v)
 - \tfrac{1}{2}\tfrac{\wp'(u) - \wp'(v)}{\wp(u) - \wp(v)}\bigr]
 \,\zeta(u) \\
 &= \zeta(u)\zeta(v) - \zeta(u)\zeta(v) - \zeta(v)^2
 - \tfrac{\zeta(v)}{2}\tfrac{\wp'(u)-\wp'(v)}{\wp(u)-\wp(v)} \\
 &\quad{}- \zeta(u)^2 - \zeta(u)\zeta(v)
 - \tfrac{\zeta(u)}{2}\tfrac{\wp'(u)-\wp'(v)}{\wp(u)-\wp(v)} \\
 &= -\zeta(u)^2 - \zeta(v)^2 - \zeta(u)\zeta(v)
 - \tfrac{\zeta(u)+\zeta(v)}{2}
 \tfrac{\wp'(u)-\wp'(v)}{\wp(u)-\wp(v)}.
\end{align*}
By the identity $\zeta(u)^2 + \zeta(v)^2 + \zeta(u)\zeta(v)
+ \tfrac{\zeta(u)+\zeta(v)}{2}
\tfrac{\wp'(u)-\wp'(v)}{\wp(u)-\wp(v)}
= -\tfrac{1}{2}[\wp(u)+\wp(v)+\wp(u+v)]$
(which follows from the Laurent expansions
$\zeta(z) = z^{-1} + O(z^3)$, $\wp(z) = z^{-2} + O(z^2)$
and the differential equation $\wp'^2 = 4\wp^3 - g_2 \wp - g_3$),
identity~\eqref{eq:fay-genus1-product} is established.

Relation~\eqref{eq:genus-g-arnold} is \emph{local}: the prime
form has the same leading singularity as the Cauchy kernel in any
coordinate chart, so the identity holds in any formal neighbourhood
of the collision diagonal.

\medskip
\emph{(b) Arnold relation at genus~$g$ implies formality.}
We prove that the dg operad $C_\ast(\FM_k(\Sigma_g))$ is formal
(quasi-isomorphic to its homology) by extending the Kontsevich
formality mechanism. The proof is \emph{not} a direct application
of Kontsevich's theorem~\cite{Kon03}, which is stated for
$\R^n$; it uses the same method (configuration space integrals
defining an $\Linf$-quasi-isomorphism) with the genus-$g$
propagator in place of the genus-$0$ one.

The input is the genus-$g$ Arnold relation~\eqref{eq:genus-g-arnold}.
The formality map
\[
 \mathcal{U}_g \colon H_\ast(\FM_k(\Sigma_g))
 \longrightarrow C_\ast(\FM_k(\Sigma_g))
\]
is constructed as a sum over admissible graphs~$\Gamma$
with propagator edges decorated by~$\omega_{ij}$:
\[
 \mathcal{U}_g(\alpha)
 = \sum_{\Gamma} \frac{1}{|\Aut(\Gamma)|}
 \int_{\FM_{|V_{\mathrm{int}}(\Gamma)|}(\Sigma_g)}
 \Bigl(\prod_{e = (i,j) \in E(\Gamma)} \omega_{ij}\Bigr)
 \wedge \alpha,
\]
where the sum ranges over graphs with $|V_{\mathrm{int}}|$ internal
vertices integrated over $\FM(\Sigma_g)$ and external vertices
labelled by the homology class~$\alpha$.
The map $\mathcal{U}_g$ must be a chain map,
which requires that the boundary contributions from the FM
compactification vanish. These boundary contributions are
governed by products $\omega_{ij} \wedge \omega_{jk}$
on collision strata, and the Arnold
relation~\eqref{eq:genus-g-arnold} ensures that the sum over
such strata telescopes to zero, exactly as in the genus-$0$
case~\cite{Kon03}.

The formality map $\mathcal{U}_g$ uses only the holomorphic
propagator $\partial_z \log E(z,w)$, for which the Fay identity
gives exact cancellation. The non-holomorphic Arakelov correction
is irrelevant to the formality statement: it enters only the curved
model~$\dfib$, not the corrected holomorphic model~$\Dg{g}$.

\emph{Remark on scope.} The extension from genus~$0$ to genus~$g$
is new: the Fay trisecant replaces the algebraic Arnold relation,
and the FM compactification $\FM_k(\Sigma_g)$ of
Sinha~\cite{Sinha2004} replaces $\FM_k(\C)$ of
Fulton--MacPherson~\cite{FM94}. The essential
mechanism (cancellation of boundary terms via the three-term
relation among propagator forms) is identical.
The resulting formality holds at fixed genus~$g$
over any field of characteristic zero.

%% ---- Complete genus-g formality: convergence + boundary cancellation ----

\providecommand{\FMg}[1]{\FM_{#1}(\Sigma_g)}
\providecommand{\omhol}{\omega^{(g)}}
\providecommand{\cU}{\mathcal{U}}
\providecommand{\cS}{\mathcal{S}}
\providecommand{\vol}{\mathrm{vol}}
\providecommand{\PP}{\mathbb{P}}

\begin{theorem}[Local screen formality and global genus-$g$ formality
under the KZB comparison; \ClaimStatusProvedHere]
\label{thm:genus-g-formality}
\index{formality!genus g@genus~$g$|textbf}
\index{configuration space integrals!genus g convergence@genus~$g$ convergence}
For every collision screen in $\FMg{k}$, the leading local model is
the genus-$0$ Fulton--MacPherson screen and the local formality
map is Kontsevich's genus-$0$ formality map. A global operadic
formality quasi-isomorphism
\[
 \cU_g \colon H_\ast(\FMg{k})
 \longrightarrow C_\ast(\FMg{k})
\]
given by prime-form configuration space integrals is conditional
on the positive-genus analytic comparison: the KZB/prime-form
propagator must produce an operadic $\Linf$ morphism compatible
with stable-graph clutching, mixed faces, and the period
correction data. At genus~$0$ this recovers Kontsevich's
formality theorem~\cite{Kon03}.
\end{theorem}

\begin{proof}
The proof below establishes the local-screen part. The final
global operadic formality statement uses the positive-genus
comparison hypotheses stated in the theorem.

\bigskip\noindent
\textbf{Step~I: The propagator.}
At genus~$0$, the Kontsevich propagator is the closed $1$-form
$\omega_{ij} = d\arg(z_i - z_j) / (2\pi)$
on $\C^2 \setminus \Delta$, with a logarithmic singularity
along the diagonal~$\Delta$. At genus $g \ge 0$, we replace
this by the holomorphic propagator
\begin{equation}\label{eq:genus-g-propagator-formal}
 \omhol_{ij}
 := \frac{1}{2\pi i}\, d_{z_i} \log E(z_i, z_j),
\end{equation}
where $E(x,y)$ is the prime form on $\Sigma_g \times \Sigma_g$:
the unique $(-\tfrac{1}{2}, -\tfrac{1}{2})$-differential
with a simple zero on~$\Delta$ and no other zeros.

\smallskip\noindent
\emph{Property~(P1): Logarithmic singularity.}
Since $E(x,y)$ vanishes to first order on~$\Delta$ with
unit leading coefficient in any local coordinate $z$, we have
\begin{equation}\label{eq:propagator-local-expansion}
 \omhol_{ij}
 = \frac{dz_i}{2\pi i\,(z_i - z_j)}
 + r_g(z_i, z_j)\, dz_i,
\end{equation}
where $r_g$ is a smooth function on
$\Sigma_g \times \Sigma_g$ (holomorphic in the first variable)
depending on the complex structure of~$\Sigma_g$. The
singularity is logarithmic: $|\omhol_{ij}| = O(|d\rho / \rho|)$
as $\rho := |z_i - z_j| \to 0$.

\smallskip\noindent
\emph{Property~(P2): Closedness.}
Since $E(x,y)$ is holomorphic in~$x$ (away from $x = y$),
the form $\omhol_{ij}$ is closed:
$d\, \omhol_{ij} = 0$ on $\Sigma_g^2 \setminus \Delta$.
This is the input for the Stokes argument in Step~IV.

\smallskip\noindent
\emph{Property~(P3): Genus-$0$ reduction.}
When $\Sigma_g = \PP^1$, $E(x,y) = x - y$, and
$\omhol_{ij} = dz_i / (2\pi i\,(z_i - z_j))$, recovering
the Kontsevich propagator upon taking real parts.

\bigskip\noindent
\textbf{Step~II: The formality map.}
An \emph{admissible graph} $\Gamma$ of type $(n,k)$ has
$n$~\emph{external} vertices (labelled $1, \ldots, n$, not
integrated) and $k$~\emph{internal} vertices (labelled
$n+1, \ldots, n+k$, integrated over $\FMg{k}$), with at most
one directed edge between any ordered pair of vertices and no
short loops (edges from a vertex to itself). For each such
$\Gamma$ with edge set $E(\Gamma)$, define
\begin{equation}\label{eq:genus-g-formality-map}
 \cU_g^{\Gamma}(\alpha_1, \ldots, \alpha_n)
 := \int_{\FMg{k}}
 \Bigl(\bigwedge_{e = (i,j) \in E(\Gamma)}
 \omhol_{ij}\Bigr)
 \wedge \pi_{\mathrm{ext}}^\ast
 (\alpha_1 \times \cdots \times \alpha_n),
\end{equation}
where $\pi_{\mathrm{ext}} \colon \FMg{n+k} \to
\FMg{n}$ is the forgetful map and the
$\alpha_\ell \in H_\ast(\FMg{n})$ are pulled back to
forms on $\FMg{n+k}$ via the projection.
The total formality map is
\[
 \cU_g := \sum_{\Gamma}
 \frac{1}{|\Aut(\Gamma)|}\, \cU_g^{\Gamma}.
\]
This is well-defined provided the integral converges
(Step~III) and satisfies the $\Linf$ relations
(Steps~IV--V).

\bigskip\noindent
\textbf{Step~III: Convergence of the configuration space
integrals.}

We show that for every admissible graph $\Gamma$ of type
$(n,k)$, the integral~\eqref{eq:genus-g-formality-map}
converges absolutely on the FM compactification~$\FMg{k}$.

\medskip\noindent
\emph{(III.a) The FM compactification as a compact manifold
with corners.}
The Fulton--MacPherson compactification $\FMg{k}$ is a compact
smooth manifold with corners, of real dimension
$\dim_{\R}\FMg{k} = 2k$ (two real coordinates per point on the
surface~$\Sigma_g$; the genus affects the topology but not the
local dimension). For manifolds of dimension
$\ge 2$, the construction is due to
Fulton--MacPherson~\cite{FM94} (algebraic) and
Axelrod--Singer~\cite{AS94} (smooth), extended to arbitrary
smooth manifolds by Sinha~\cite{Sinha2004}.

The boundary strata of $\FMg{k}$ are indexed by
\emph{screen data}: collections
$\cS = \{S_1, \ldots, S_r\}$ of subsets of $\{1, \ldots, k\}$
with $|S_\alpha| \ge 2$, satisfying the nesting condition
(any two $S_\alpha, S_\beta$ are either disjoint or
nested: $S_\alpha \subset S_\beta$ or
$S_\beta \subset S_\alpha$). The stratum
$\partial_{\cS}\FMg{k}$ has real codimension $2(|S_\alpha| - 1)$
for each screen~$S_\alpha$, recording the directions in which
the points of $S_\alpha$ approach a common limit.

\medskip\noindent
\emph{(III.b) Local coordinates near a boundary stratum.}
Let $S = \{i_1, \ldots, i_p\} \subset \{1, \ldots, k\}$ be a
subset of $p \ge 2$ points colliding at a centre
$z_0 \in \Sigma_g$. Near the stratum $\partial_S \FMg{k}$,
the FM compactification provides real-analytic coordinates:
\begin{itemize}
\item the \emph{centre} $z_0 \in \Sigma_g$
 (a local coordinate on the base),
\item the \emph{scale} $\rho = \max_{a,b \in S}
 |z_{i_a} - z_{i_b}| \in [0, \epsilon)$,
\item the \emph{screen directions}
 $\mathbf{u} = (u_2, \ldots, u_p) \in S^{2p-3}$,
 the unit-sphere normalisation of the relative positions
 $(z_{i_2} - z_{i_1}, \ldots, z_{i_p} - z_{i_1})
 / \rho$,
\item the \emph{remaining points}
 $\{z_j\}_{j \notin S}$, which stay away from~$z_0$.
\end{itemize}
The boundary stratum $\partial_S$ corresponds to
$\rho = 0$; the screen directions $\mathbf{u}$ parametrise
an $\FM_{|S|}(\C)$-factor (the genus-$0$ FM space of the
colliding cluster, because at infinitesimal scale the curve
is indistinguishable from its tangent space
$T_{z_0}\Sigma_g \cong \C$). The volume form decomposes as
\begin{equation}\label{eq:volume-form-decomposition}
 \vol_{\FMg{k}}\big|_{\text{near } \partial_S}
 = \rho^{2p - 3}\, d\rho \wedge
 \vol_{S^{2p-3}} \wedge \vol_{\Sigma_g}^{(z_0)}
 \wedge \vol_{\text{remaining}}.
\end{equation}

\medskip\noindent
\emph{(III.c) Singularity estimates for the propagator product.}
We estimate the integrand
$\bigwedge_{e \in E(\Gamma)} \omhol_{ij}$ near $\partial_S$.
There are three cases for an edge $e = (i,j) \in E(\Gamma)$:

\smallskip
\noindent\emph{Case~1: Both $i, j \in S$ (internal edge).}
By~\eqref{eq:propagator-local-expansion}, as
$\rho \to 0$:
\[
 \omhol_{ij}
 = \frac{d(z_i - z_j)}{2\pi i\,(z_i - z_j)} + O(1)
 = \frac{d(\rho\, u_{ij})}
 {2\pi i\, \rho\, u_{ij}} + O(1)
 = \frac{d\rho}{2\pi i\, \rho}
 + \frac{d u_{ij}}{2\pi i\, u_{ij}} + O(1),
\]
where $u_{ij}$ is the relative screen coordinate. Thus
$|\omhol_{ij}| = O(|d\rho/\rho|) + O(1)$: the singularity is
at most logarithmic in~$\rho$. For $p_S$ edges internal to~$S$,
the product has singularity at most
$|d\rho/\rho|^{p_S} \cdot (\text{bounded form on } S^{2p-3})$.

\smallskip
\noindent\emph{Case~2: $i \in S$, $j \notin S$ (connecting edge).}
As $\rho \to 0$, $z_i \to z_0$ while $z_j$ stays bounded away
from~$z_0$. Then
$\omhol_{ij} \to \omhol_{z_0, j}$: the propagator has a
regular limit, with no singularity.

\smallskip
\noindent\emph{Case~3: Both $i, j \notin S$ (external edge).}
The propagator $\omhol_{ij}$ is independent of the collision
at~$S$ and is smooth.

\smallskip
Combining: near $\partial_S$, the integrand has singularity
at most $|d\rho/\rho|^{p_S}$ (from internal edges), and the
integration measure contributes
$\rho^{2p-3}\, d\rho$ from~\eqref{eq:volume-form-decomposition}.
Since $\Gamma$ is admissible (at most one edge per ordered pair),
$p_S \le \binom{|S|}{2}$, and the integral over the $\rho$
coordinate is
\begin{equation}\label{eq:rho-convergence-estimate}
 \int_0^{\epsilon}
 \Bigl|\frac{d\rho}{\rho}\Bigr|^{p_S}
 \cdot \rho^{2p-3}\, d\rho
 = \int_0^{\epsilon}
 \rho^{2p - 3 - p_S}\,
 \bigl|\!\log\rho\bigr|^{q_S}\, d\rho,
\end{equation}
where $q_S \le p_S$ counts the number of edges that contribute a
pure $|d\rho/\rho|$ factor (as opposed to an angular factor
$|du_{ij}/u_{ij}|$ absorbed by the screen measure).
More precisely: each internal edge contributes a $1$-form on the
total space near~$\partial_S$. In the coordinates
$(\rho, \mathbf{u}, z_0, \text{rest})$, this $1$-form
decomposes as $\omhol_{ij} = (d\rho/\rho) \cdot \alpha_{ij}
+ \beta_{ij}$, where $\alpha_{ij}$ is a function on the screen
and $\beta_{ij}$ is a $1$-form with no $d\rho$ component. The
wedge product of $p_S$ such forms is a sum of terms, each
containing at most one factor of $d\rho/\rho$ (since
$(d\rho)^2 = 0$). Thus the $d\rho$-singularity is at most
$O(\rho^{-1})$, and the integral
\[
 \int_0^{\epsilon} \rho^{-1} \cdot \rho^{2p-3}\, d\rho
 = \int_0^{\epsilon} \rho^{2(p-1) - 2}\, d\rho
 = \int_0^{\epsilon} \rho^{2p - 4}\, d\rho
\]
converges for $p \ge 2$ (i.e., $|S| \ge 2$), since
$2p - 4 \ge 0$ and $\int_0^\epsilon \rho^a\, d\rho < \infty$
for $a \ge 0$.

\medskip\noindent
\emph{(III.d) The case $|S| = 2$: pairwise collisions.}
When $|S| = \{i, j\}$, there is at most one internal edge
($p_S \le 1$) and the scale coordinate satisfies
$\rho = |z_i - z_j|$. The measure is $\rho\, d\rho$
(from $2p - 3 = 1$). The propagator contributes
$|d\rho/\rho|$ at worst, so the integral is
$\int_0^\epsilon \rho^{-1} \cdot \rho\, d\rho
= \int_0^\epsilon d\rho = \epsilon < \infty$.
The singularity is logarithmic: the prime form produces
a $\log|z_i - z_j|$ divergence, which is integrable in two
real dimensions.

\medskip\noindent
\emph{(III.e) Induction on the number of screens.}
For nested screens $\cS = \{S_1, \ldots, S_r\}$
(ordered by inclusion), the integral near
$\partial_{\cS}\FMg{k}$ factors by the iterated blow-up
structure of the FM compactification. Each screen $S_\alpha$
contributes a scale parameter $\rho_\alpha$ and a
screen integral. By the single-screen estimate above,
each factor converges independently. The total integral is a
product of convergent integrals over each scale parameter
(since the screens are nested, the scale parameters form an
ordered family $\rho_{S_1} \le \rho_{S_2} \le \cdots$ that
decouple in the iterated blow-up coordinates).

Formally: we proceed by induction on $k$, the number of
internal vertices. The base case $k = 0$ is trivial (no
integration). For the inductive step at $k$ vertices:
partition the boundary $\partial\FMg{k}$ into strata
$\partial_S$ for minimal screens $S$ (screens that are not
properly contained in any other screen). At each
$\partial_S$, the integral factors as
\[
 \cU_g^{\Gamma}\big|_{\partial_S}
 = \Bigl(\int_{\FM_{|S|}(\C)} (\text{internal propagators})
 \Bigr)
 \cdot
 \Bigl(\int_{\FMg{k - |S| + 1}}
 (\text{remaining propagators})\Bigr).
\]
The first factor (the ``screen integral'') converges by the
genus-$0$ estimate (it is an integral on $\FM_{|S|}(\C)$ with
the Kontsevich propagator, which converges by~\cite{Kon03}).
The second factor converges by the inductive hypothesis
(it is a configuration space integral on $\FMg{k'}$ with
$k' = k - |S| + 1 < k$ internal vertices).

This completes the proof that $\cU_g^\Gamma$ converges
absolutely for every admissible graph~$\Gamma$.

\bigskip\noindent
\textbf{Step~IV: Boundary cancellation (the chain map property).}

We show that the
boundary contributions to $d \cU_g$ vanish, so that
$\cU_g$ is a chain map (and more generally satisfies the
$\Linf$ relations).

\medskip\noindent
\emph{(IV.a) Setup via Stokes.}
Since the propagator $\omhol_{ij}$ is closed
(Property~(P2)), the Stokes theorem on the compact manifold
with corners $\FMg{k}$ gives:
\begin{equation}\label{eq:stokes-genus-g}
 d\bigl(\cU_g^{\Gamma}(\alpha)\bigr)
 = \int_{\partial\FMg{k}}
 \Bigl(\bigwedge_{e \in E(\Gamma)} \omhol_{ij}\Bigr)
 \wedge \pi_{\mathrm{ext}}^\ast(\alpha)
 = \sum_{\cS} \int_{\partial_{\cS}\FMg{k}}
 \Bigl(\bigwedge_{e \in E(\Gamma)} \omhol_{ij}\Bigr)
 \wedge \pi_{\mathrm{ext}}^\ast(\alpha),
\end{equation}
where the sum is over codimension-$1$ boundary strata.
Each codimension-$1$ stratum corresponds to a \emph{single}
screen $S = \{i_1, \ldots, i_p\}$ of $p \ge 2$ internal
vertices colliding at a point $z_0 \in \Sigma_g$.

\medskip\noindent
\emph{(IV.b) Factorisation of the boundary integral.}
On the boundary stratum $\partial_S\FMg{k}$
(where $\rho_S = 0$), the FM compactification identifies the
stratum as a fibre bundle:
\[
 \partial_S\FMg{k}
 \;\cong\;
 \FM_{|S|}(\C) \;\times\; \FMg{k - |S| + 1},
\]
where $\FM_{|S|}(\C)$ parametrises the infinitesimal
configuration of the colliding cluster (the screen) and
$\FMg{k - |S| + 1}$ parametrises the remaining points
on~$\Sigma_g$ with the collision centre $z_0$ replacing
the cluster. The propagators decompose accordingly:
\begin{itemize}
\item \emph{Internal edges} $(i,j)$ with both endpoints in $S$:
 the propagator $\omhol_{ij}$ restricts to the genus-$0$
 propagator $\omega^{(0)}_{ij}$ on the screen factor
 $\FM_{|S|}(\C)$. This holds because at infinitesimal
 scale, $\omhol_{ij} = dz_i / (2\pi i\,(z_i - z_j))
 + O(\rho)$, and the $O(\rho)$ correction vanishes on the
 boundary $\rho = 0$.

\item \emph{Connecting edges} $(i, j)$ with $i \in S$,
 $j \notin S$: the propagator limits to $\omhol_{z_0, j}$,
 a smooth form on $\FMg{k - |S| + 1}$. This form depends
 on the screen only through the base point~$z_0$.

\item \emph{External edges} $(i,j)$ with neither endpoint in~$S$:
 smooth and independent of the screen.
\end{itemize}
Consequently, the boundary integral factors:
\begin{equation}\label{eq:boundary-factorisation}
 \int_{\partial_S\FMg{k}}
 \bigwedge_{e \in E(\Gamma)} \omhol_{ij}
 \;=\;
 \underbrace{\Bigl(
 \int_{\FM_{|S|}(\C)}
 \bigwedge_{\substack{e = (i,j) \\ i,j \in S}}
 \omega^{(0)}_{ij}
 \Bigr)}_{\text{screen integral } I_S(\Gamma)}
 \;\cdot\;
 \underbrace{\Bigl(
 \int_{\FMg{k - |S| + 1}}
 (\text{connecting and external propagators})
 \wedge \alpha
 \Bigr)}_{\text{external integral } J_S(\Gamma)}.
\end{equation}

\medskip\noindent
\emph{(IV.c) Vanishing of the screen integral by the Arnold
relation.}
The screen integral $I_S(\Gamma)$ is a configuration space
integral on the genus-$0$ space $\FM_{|S|}(\C)$ with the
standard Kontsevich propagator. The following shows it vanishes.

The chain map condition for
$\cU_g$ requires the vanishing of the \emph{sum} of screen
integrals over all graphs $\Gamma'$ obtained by collapsing the
subgraph $\Gamma|_S$ to a single vertex. This sum is
controlled by the $\Linf$ structure equations, which at the
level of the screen integral reduce to the genus-$0$ Arnold
relation.

Concretely, consider $|S| = 2$, $S = \{i,j\}$.
The screen is $\FM_2(\C) \cong S^1$, and the screen integral
is $\int_{S^1} \omega^{(0)}_{ij}$. This integral equals~$1$
(the winding number), and its contribution to $d\cU_g$ is
accounted for by the differential on the source
$H_\ast(\FMg{k})$; it is the component of the $\Linf$
structure map $\ell_2$ corresponding to the binary collision.
This is not a ``cancellation'' but an \emph{identification}:
the $|S| = 2$ boundary strata produce the $\Linf$ structure
maps themselves.

For $|S| \ge 3$, the screen integral involves products of
propagators on $\FM_{|S|}(\C)$. The Arnold relation at
genus~$0$, which states
\begin{equation}\label{eq:arnold-genus0-repeat}
 \omega^{(0)}_{ij} \wedge \omega^{(0)}_{jk}
 + \omega^{(0)}_{jk} \wedge \omega^{(0)}_{ki}
 + \omega^{(0)}_{ki} \wedge \omega^{(0)}_{ij}
 = 0
 \quad\text{on } \Conf_3(\C)
\end{equation}
and generalises to all higher degrees via the
Arnold--Cohen presentation of $H^\ast(\Conf_k(\C))$,
forces the screen integrals to satisfy the same
cancellation that Kontsevich exploits at genus~$0$. Precisely:
the sum $\sum_{\Gamma'} I_S(\Gamma')$ over all graphs
$\Gamma'$ with internal vertices $S$ that contribute to
$d\cU_g$ at the stratum $\partial_S$ vanishes by
Kontsevich's vanishing lemma~\cite{Kon03} (Lemma~6.6),
which states that the sum of configuration space integrals
over all admissible completions of a subgraph on $\FM_p(\C)$
for $p \ge 3$ vanishes.

This vanishing depends only on
the \emph{screen factor} $\FM_{|S|}(\C)$, which is a
genus-$0$ space \emph{regardless} of the genus of~$\Sigma_g$.
The higher-genus information of $\Sigma_g$ enters only through
the external factor $J_S(\Gamma)$, which is a regular
(non-singular) integral serving as an overall coefficient.
Therefore the genus-$0$ vanishing mechanism applies without
modification to the genus-$g$ boundary strata.

\medskip\noindent
\emph{(IV.d) Compatibility with the Fay trisecant.}
We invoke the genus-$0$ Arnold relation on
the screen rather than the genus-$g$ Fay trisecant
identity~\eqref{eq:genus-g-arnold} because the
two are the same on the screen: the restriction of
$\omhol_{ij}$ to the codimension-$1$ boundary stratum
$\partial_S$ yields the genus-$0$ propagator on the screen
factor, and the Fay identity restricted to an infinitesimal
neighbourhood of any point $z_0 \in \Sigma_g$ reduces to
the classical Arnold relation (because the prime form agrees
with the Cauchy kernel to leading order:
$E(x,y) = (x - y)(1 + O(|x-y|^2))$ in local coordinates).

The Fay trisecant~\eqref{eq:genus-g-arnold} is essential
for a different reason: it guarantees that the \emph{global}
propagator $\omhol_{ij}$ satisfies the Arnold relation on
the full configuration space $\Conf_k(\Sigma_g)$,
not just on the infinitesimal screens, which is needed to
ensure that the formality map $\cU_g$ is well-defined
at the level of homology classes (not just at the chain
level). This global consistency is precisely what
Step~(a) of the present proof of
Theorem~\ref{thm:modular-hkoszul-SC} establishes.

\bigskip\noindent
\textbf{Step~V: The $\Linf$ relations and quasi-isomorphism
property.}

\medskip\noindent
\emph{(V.a) The $\Linf$ equations.}
The maps $\{\cU_g^{(n)}\}_{n \ge 1}$ defined by summing
$\cU_g^\Gamma$ over graphs with $n$ external vertices
must satisfy the $\Linf$ morphism equations:
\begin{equation}\label{eq:Linfty-morphism-equations}
 \sum_{\substack{p + q = n+1 \\ p,q \ge 1}}
 \sum_{\sigma \in \Sh(p,q)}
 \pm\, \cU_g^{(p)}\bigl(
 \ell_q^{H}(\alpha_{\sigma(1)}, \ldots,
 \alpha_{\sigma(q)}),
 \alpha_{\sigma(q+1)}, \ldots, \alpha_{\sigma(n)}\bigr)
 = \sum_{\substack{p + q = n+1 \\ p,q \ge 1}}
 \sum_\sigma \pm\,
 \ell_q^{C}\bigl(\cU_g^{(p)}(\alpha_{\sigma(1)}, \ldots),
 \ldots\bigr),
\end{equation}
where $\ell_q^H$ and $\ell_q^C$ are the $\Linf$ structure
maps on $H_\ast(\FMg{k})$ and $C_\ast(\FMg{k})$
respectively, and $\Sh(p,q)$ denotes $(p,q)$-shuffles.
Each of these equations is equivalent to the vanishing of
a specific combination of boundary integrals on $\FMg{k}$,
by the Stokes argument of Step~IV. The $|S| = 2$ boundary
strata contribute the terms involving $\ell_2$; the
$|S| \ge 3$ strata vanish by the Arnold/Kontsevich vanishing
lemma. This is identical to the genus-$0$ argument
of~\cite{Kon03}, with the only modification being the use of
the genus-$g$ propagator $\omhol_{ij}$ in place of the
genus-$0$ one, and the boundary factorisation (IV.b) ensures
that the screen integrals see only the genus-$0$ propagator.

\medskip\noindent
\emph{(V.b) Quasi-isomorphism.}
It remains to show that $\cU_g$ is a \emph{quasi-isomorphism},
i.e., that its linear part $\cU_g^{(1)}$ induces an isomorphism
on homology. We give a direct argument.

The linear part $\cU_g^{(1)}$ is the sum over graphs $\Gamma$
with $n = 1$ external vertex and $k$ internal vertices:
$\cU_g^{(1)}(\alpha) = \alpha + (\text{higher order in }k)$.
At $k = 0$ (no internal vertices), $\cU_g^{(1)}$ is the
identity inclusion
$H_\ast(\FMg{n}) \hookrightarrow C_\ast(\FMg{n})$.
This is a quasi-isomorphism at the cochain level by definition
of homology. More precisely, $\cU_g^{(1)}$ is a perturbation
of the identity by terms of positive filtration degree (in
the filtration by the number of internal vertices), so
$\cU_g^{(1)}$ is a quasi-isomorphism by the perturbation
lemma: if $f = \id + \epsilon$ where $\epsilon$ raises
filtration degree, then $f$ is a quasi-isomorphism on any
bounded-below filtered complex.

\medskip\noindent
\emph{(V.c) Independence of complex structure.}
The formality of $C_\ast(\FMg{k})$ is a property of the
dg operad (chain complex with composition maps), and the
underlying topological space $\FMg{k}$ depends only on the
diffeomorphism type of~$\Sigma_g$, not on the complex
structure. All smooth genus-$g$ surfaces are diffeomorphic,
so the \emph{homotopy type} of $\FMg{k}$ is independent of
the choice of complex structure on~$\Sigma_g$. The formality
map $\cU_g$, however, uses the complex structure through the
propagator~$\omhol_{ij}$. Different complex structures yield
different but homotopic formality maps, because:
\begin{enumerate}[label=\textup{(\arabic*)}]
\item The prime form $E(z_i, z_j)$ varies holomorphically on
 the Torelli space $\cT_g$ (the universal cover of the moduli
 space $\cM_g$).
\item The integrals~\eqref{eq:genus-g-formality-map} vary
 continuously with the propagator (by dominated convergence,
 since the convergence established in Step~III is uniform over
 compact subsets of~$\cT_g$).
\item The $\Linf$ quasi-isomorphism $\cU_g$ therefore varies
 continuously over~$\cT_g$. Since the space of $\Linf$
 quasi-isomorphisms between fixed dg operads is either empty
 or contractible (by obstruction theory for $\Linf$ morphisms),
 $\cU_g$ is unique up to homotopy.
\end{enumerate}
In particular, formality at \emph{any one} complex structure
on $\Sigma_g$ implies formality at \emph{every} complex
structure. And for each~$g$, we have constructed an explicit
complex structure (any smooth projective curve) at which
$\cU_g$ is defined and satisfies all the required properties.
This completes the proof.
\end{proof}

\begin{remark}[Scope and relation to genus~$0$]
\label{rem:genus-g-formality-scope}
The proof of Theorem~\ref{thm:genus-g-formality} reduces the
analytical content of genus-$g$ formality to genus~$0$, via the
boundary factorisation~\eqref{eq:boundary-factorisation}:
every boundary stratum of $\FMg{k}$ factors into a genus-$0$
screen and a lower-complexity genus-$g$ integral. The
convergence and cancellation arguments at genus~$g$ are thus
not independent of Kontsevich's original proof but are
\emph{fibrewise applications} of it, mediated by the prime form
propagator and the FM compactification for manifolds. The
genuinely new input is the Fay trisecant identity, which
guarantees that the global propagator is compatible with the
Arnold relation at all genera; this is the content of
Step~(a) above.
\end{remark}

\medskip
\emph{(c) Formality implies Feynman involutivity.}
Given formality of $C_\ast(\FM_k(\Sigma_g))$
(Theorem~\ref{thm:genus-g-formality}), the modular operad
$\SCmodcl$ is quasi-isomorphic to its homology operad
$H_\ast(\SCmodcl) = \{H_\ast(\FM_k(\Sigma_g))\}_{g,k}$.
The Feynman transform $\FTmod$ of
Getzler--Kapranov~\cite{GK98} satisfies
$\FTmod^2(\cP) \simeq \cP$ for any modular operad~$\cP$ whose
underlying symmetric sequence is concentrated in even
degrees or is formal. This is the modular analogue
of the bar-cobar involution for ordinary operads;
the proof (Getzler--Kapranov~\cite{GK98}, Theorem~8.7)
uses the Feynman transform as a twisted modular operadic
bar construction indexed by stable graphs, and the involutivity
follows from the Poincar\'e duality of the Feynman amplitudes
on the moduli space~$\overline{\cM}_{g,n}$.

More precisely: the Feynman transform of a modular operad~$\cP$
is the modular cooperad
$\FTmod(\cP)(g,n)
= \bigoplus_{\Gamma \in \StGraph(g,n)}
 \bigotimes_{v \in V(\Gamma)}
 \cP(g(v), n(v))^{\vee} \otimes \det(E(\Gamma))$,
with differential summing over edge contractions.
Applying $\FTmod$ twice gives a double sum over pairs of stable
graphs, and the involutivity $\FTmod^2 \simeq \id$ follows by
showing that the double complex computes the identity via a
spectral sequence argument. When $\cP$ is formal, this
spectral sequence degenerates at~$E_1$.

Combined: $C_\ast(\FM_k(\Sigma_g))$ is formal
(Theorem~\ref{thm:genus-g-formality}),
so $\FTmod^2(\SCmodcl) \simeq \SCmodcl$, and the
bar-cobar adjunction for the closed colour is a Quillen
equivalence at each genus.

\begin{remark}[Sketch of Feynman involutivity]
\label{rem:feynman-involutivity-sketch}
The Getzler--Kapranov theorem (Volume~I,
Theorem~\ref*{V1-thm:feynman-involution}; original
\cite{GK98}, Theorem~8.7) states
$\FTmod^2 \simeq \mathrm{id}$ for modular operads satisfying
Poincar\'e duality. The Feynman transform
$\FTmod(\cP)$ of a modular operad~$\cP$ sums over stable
graphs~$\Gamma$ with vertex labels in~$\cP$ and edge labels
in the orientation lines $\det(E(\Gamma))$. Applying $\FTmod$
twice sums over pairs of graphs $(\Gamma, \Gamma')$ where
$\Gamma'$ is a subgraph of~$\Gamma$. The Poincar\'e duality
on the operation spaces of~$\cP$ identifies the double sum with
a single sum over~$\Gamma$ weighted by the Euler characteristic
factor $(-1)^{|E(\Gamma)|}$, which collapses to the identity by
the M\"obius inversion formula on the lattice of subgraphs.
The Poincar\'e duality holds for
$\{\overline{\cM}_{g,n}\}$ by Serre duality on moduli spaces.
\end{remark}

\medskip\noindent
\textbf{Step~3: Mixed-operation compatibility via the filtered
associated graded.}
The mixed operations of $\SCmod$ do not satisfy a strict product
equivalence on the full compactification: Type~III mixed bubbling
faces contribute the mixed differential. Proposition~\ref{prop:mixed-product-decomposition}
identifies only the depth-zero associated graded,
\[
\gr^0 C_\ast(\FM_{k|m})
\cong C_\ast(\FM_k)\otimes C_\ast(\Eone(m)).
\]
The Swiss-cheese directionality ensures that the genus-raising
differential $D_1$ (the non-separating clutching map) acts
only on the \emph{closed} factor: contracting a non-separating
edge raises the genus of the surface but does not affect the
topological interval. Formally, if $\xi_{\mathrm{nsep}}$ denotes
the non-separating clutching,
\[
 \xi_{\mathrm{nsep}}^{\mathrm{mix}}
 = \xi_{\mathrm{nsep}}^{\mathrm{cl}} \otimes \id_{\Eone}.
\]

We require that the bar complex of the mixed-colour operad
decomposes compatibly. The precise statement is the following
operadic K\"unneth lemma.

\begin{lemma}[Operadic K\"unneth for the mixed bar complex]
\label{lem:operadic-kunneth}
Let $\cP_1$ be a modular operad with finite-dimensional operations
over a field of characteristic zero, and let $\cP_2$ be an
augmented genus-$0$ operad with finite-dimensional operations.
Define the partially modular operad
$\cP = \cP_1 \boxtimes \cP_2$ with
$\cP_{\mathrm{cl}} = \cP_1$,
$\cP_{\mathrm{op}} = \cP_2$, and
$\cP_{\mathrm{mix}}(g,k,m) = \cP_1(g,k) \otimes \cP_2(m)$.
Equip $\mathbf{B}(\cP_{\mathrm{mix}})$ with the increasing
filtration $F_\bullet$ by the closed-input excess
\[
\mathrm{wt}(\Gamma)
:=
\sum_{v \in V(\Gamma)} \max\{k(v)-1, 0\},
\]
where $k(v)$ is the number of closed inputs at the mixed
vertex~$v$. Then there is an isomorphism of dg cooperads
\[
 \gr_F \mathbf{B}(\cP_{\mathrm{mix}})
 \cong
 \mathbf{B}_{\mathrm{mod}}(\cP_1)
 \otimes \mathbf{B}(\cP_2),
\]
where $\mathbf{B}_{\mathrm{mod}}$ denotes the modular
\textup{(}stable-graph\textup{)} bar construction,
$\mathbf{B}$ denotes the tree-level bar construction, and the
differential on the right is
$d_{\mathrm{cl}} \otimes \id + \id \otimes d_{\mathrm{op}}$.
\end{lemma}

\begin{proof}[Proof of Lemma~\textup{\ref{lem:operadic-kunneth}}]
The bar construction of $\cP_{\mathrm{mix}}$ is the chain complex
\[
 \mathbf{B}(\cP_{\mathrm{mix}})
 = \bigoplus_{\Gamma}
 \bigotimes_{v \in V(\Gamma)}
 \cP_{\mathrm{mix}}(g(v), k(v), m(v))^{\vee}
 \otimes \det(E(\Gamma)),
\]
where the sum ranges over mixed graphs $\Gamma$: graphs whose
vertices carry both genus labels and degree labels, whose edges
are either \emph{closed} (connecting closed inputs) or \emph{open}
(connecting open inputs, respecting the ordering).
By the product structure
$\cP_{\mathrm{mix}}(g,k,m) = \cP_1(g,k) \otimes \cP_2(m)$,
each vertex tensor factor splits as
$\cP_1(g(v), k(v))^{\vee} \otimes \cP_2(m(v))^{\vee}$.
A mixed graph $\Gamma$ decomposes uniquely as
$\Gamma = \Gamma_{\mathrm{cl}} \sqcup \Gamma_{\mathrm{op}}$
where $\Gamma_{\mathrm{cl}}$ is a stable graph (governing the
closed-colour edges) and $\Gamma_{\mathrm{op}}$ is a planar tree
(governing the open-colour edges). This decomposition exists
because directionality forbids edges from open to closed vertices,
so closed and open edges cannot share endpoints across
colours.
The bar differential has three pieces:
$d_{\mathbf{B}} = d_{\mathrm{cl}} + d_{\mathrm{op}} + d_{\mathrm{mix}}$.
The closed part $d_{\mathrm{cl}}$ contracts closed edges and
performs the modular operadic insertion in $\cP_1$.
The open part $d_{\mathrm{op}}$ contracts open edges for which at
most one endpoint carries closed inputs, so it acts on the
$\cP_2$-factor and preserves the closed-input distribution.
The remaining terms contract an open edge whose two endpoints are
mixed vertices with positive closed valence; these terms define
$d_{\mathrm{mix}}$.

If an open edge joins vertices with $k'$ and $k''$ closed inputs,
then $d_{\mathrm{mix}}$ replaces them by a single mixed vertex with
$k' + k''$ closed inputs. Hence
\[
 \max\{k' + k'' - 1, 0\}
 =
 \max\{k'-1, 0\} + \max\{k''-1, 0\} + 1
\]
when $k', k'' > 0$, so $d_{\mathrm{mix}}$ raises the filtration
$F_\bullet$ by one, while $d_{\mathrm{cl}}$ and $d_{\mathrm{op}}$
preserve it. Therefore $d_{\mathrm{mix}}$ vanishes on the
associated graded, and the induced differential on
$\gr_F \mathbf{B}(\cP_{\mathrm{mix}})$ is
$d_{\mathrm{cl}} \otimes \id + \id \otimes d_{\mathrm{op}}$.

The underlying graded cooperad is
\[
 \gr_F \mathbf{B}(\cP_{\mathrm{mix}})
 \cong
 \Bigl(\bigoplus_{\Gamma_{\mathrm{cl}}}
 \bigotimes_v \cP_1(g(v), k(v))^{\vee}
 \otimes \det(E(\Gamma_{\mathrm{cl}}))\Bigr)
 \otimes
 \Bigl(\bigoplus_{\Gamma_{\mathrm{op}}}
 \bigotimes_v \cP_2(m(v))^{\vee}
 \otimes \det(E(\Gamma_{\mathrm{op}}))\Bigr),
\]
which is $\mathbf{B}_{\mathrm{mod}}(\cP_1)
\otimes \mathbf{B}(\cP_2)$.
Since $F_\bullet$ is exhaustive and bounded below, this identifies
the $E_0$-page of the spectral sequence attached to
$\mathbf{B}(\cP_{\mathrm{mix}})$ with the tensor-product bar
complex. Over a field of characteristic zero with finite-dimensional
operations, the $E_1$-page is therefore
$H_\ast(\mathbf{B}_{\mathrm{mod}}(\cP_1))
\otimes H_\ast(\mathbf{B}(\cP_2))$
by the algebraic K\"unneth theorem
\textup{(}Weibel~\cite{Wei94}, Theorem~3.6.3\textup{)}.
\end{proof}

\begin{remark}
\label{rem:operadic-kunneth-cross-terms}
The failure of a chain-level factorization comes from the
cross-term $d_{\mathrm{mix}}$: contracting an open edge between
mixed vertices with positive closed valence combines their closed
inputs without a closed operadic insertion. These are the
$\mu_1$-terms omitted by the naive identity
$d_{\mathbf{B}} = d_{\mathrm{cl}} \otimes \id
+ \id \otimes d_{\mathrm{op}}$.
They vanish on $\gr_F$ because they raise the closed-input
filtration. For a fixed closed graph and open tree, the summands
generated by repeated applications of $d_{\mathrm{mix}}$ form
contractible subcomplexes of the bar construction, so the
spectral sequence of $F_\bullet$ degenerates at $E_1$.
\end{remark}

Applying Lemma~\ref{lem:operadic-kunneth} with
$\cP_1 = \SCmodcl$ and $\cP_2 = \Eone$:
both factors are homotopy-Koszul (the closed by
Step~2, the open by the classical Koszulity of $\Eone = \Ass$).
The induced filtration on the mixed bar-cobar complex has
associated graded governed by
$\mathbf{B}_{\mathrm{mod}}(\SCmodcl) \otimes \mathbf{B}(\Eone)$.
The resulting filtered bar-cobar spectral sequence has first page
computed from that tensor-product bar complex. Since the closed
and open counits are quasi-isomorphisms, the induced comparison
map on this $E_1$-page is a quasi-isomorphism. By
Remark~\ref{rem:operadic-kunneth-cross-terms}, the higher
filtration differentials come from the acyclic cross-term
$d_{\mathrm{mix}}$, so the spectral sequence collapses at $E_1$.
Hence the mixed bar-cobar counit is a quasi-isomorphism.

\medskip\noindent
\textbf{Step~4: Assembly via spectral sequence convergence.}
The genus filtration $F^{\le g} \SCmodcl
= \bigoplus_{g' \le g} \SCmodcl(g', \bullet)$
induces a filtration on the bar-cobar complex
$\Omega\mathbf{B}(\SCmod)$.

\medskip
\emph{(a) Construction of the spectral sequence.}
Define the filtration on $\Omega\mathbf{B}(\SCmod)$ by
\[
 F^{\le p}\, \Omega\mathbf{B}(\SCmod)
 := \Omega\mathbf{B}\bigl(F^{\le p}\, \SCmod\bigr),
\]
where $F^{\le p}\, \SCmod$ is the partially modular suboperad
generated by closed-colour operations of genus~$\le p$.
This is an increasing, exhaustive filtration:
$F^{\le 0} \subset F^{\le 1} \subset \cdots$ and
$\bigcup_p F^{\le p} = \Omega\mathbf{B}(\SCmod)$.
(Completeness: $F^{\le -1} = 0$ since there are no operations
of negative genus.)

The associated spectral sequence has
\[
 E_0^{p,q}
 = F^{\le p}\, \Omega\mathbf{B}(\SCmod)^{p+q} \big/
 F^{\le p-1}\, \Omega\mathbf{B}(\SCmod)^{p+q}
 \cong
 \gr_F^p\, \Omega\mathbf{B}(\SCmod)^{p+q},
\]
with $d_0$ induced by the part of the bar-cobar differential
that preserves genus.

\medskip
\emph{(b) Identification of the $E_1$ page.}
The $d_0$-differential on $\gr_F^p$ is the cobar-bar differential
of the genus-$p$ contribution at fixed genus. By Steps~2 and~3,
the cobar-bar counit is a quasi-isomorphism at each fixed genus,
so
\[
 E_1^{p,q}
 = H^{p+q}(\gr_F^p\, \Omega\mathbf{B}(\SCmod), d_0)
 \cong \SCmod(p, \bullet)^{p+q},
\]
i.e., the $E_1$ page recovers the genus-$p$ operations of~$\SCmod$.

\medskip
\emph{(c) Convergence.}
The filtration is bounded below ($F^{\le p} = 0$ for $p < 0$)
and exhaustive. By the classical convergence theorem for
spectral sequences of filtered chain complexes
(Weibel~\cite{Wei94}, Theorem~5.5.1; the ``bounded below''
case), the spectral sequence converges:
\[
 E_\infty^{p,q} \cong \gr_F^p\, H^{p+q}(\Omega\mathbf{B}(\SCmod)).
\]
The Mittag-Leffler condition is automatic for bounded-below
exhaustive filtrations of cochain complexes over a field
(Weibel~\cite{Wei94}, Theorem~5.5.7): the filtration
$F^{\le p}$ is bounded below ($F^{\le -1} = 0$, since there
are no operations of negative genus) and exhaustive
($\bigcup_p F^{\le p} = \Omega\mathbf{B}(\SCmod)$), so the
classical convergence theorem (Weibel~\cite{Wei94},
Theorem~5.5.1) applies directly without needing to verify
Mittag-Leffler separately; the bounded-below hypothesis
subsumes it.

\medskip
\emph{(d) The counit is a quasi-isomorphism.}
Consider the counit $\epsilon \colon \Omega\mathbf{B}(\SCmod)
\to \SCmod$ as a map of filtered chain complexes, where $\SCmod$
carries the genus filtration
$F^{\le p}\, \SCmod = \bigoplus_{g' \le p} \SCmod(g', \bullet)$.
On the $E_1$ page, $\epsilon$ induces the identity map (by~(b)).
Hence $\epsilon$ induces an isomorphism on all pages of the
spectral sequence, and the five-lemma applied inductively on
the filtration quotients gives that $\epsilon$ is a
quasi-isomorphism on the total complexes.

More precisely: let $\epsilon^{\le p}$ denote the restriction
of~$\epsilon$ to the $p$-th filtration level. We show
$\epsilon^{\le p}$ is a quasi-isomorphism by induction on~$p$.
The base case $p = 0$ is the genus-$0$ homotopy-Koszulity
(Step~1). For the inductive step, the short exact sequence of
chain complexes
\[
 0 \to F^{\le p-1}\, \Omega\mathbf{B}(\SCmod)
 \to F^{\le p}\, \Omega\mathbf{B}(\SCmod)
 \to \gr_F^p\, \Omega\mathbf{B}(\SCmod) \to 0
\]
maps via~$\epsilon$ to the corresponding sequence for~$\SCmod$.
By the inductive hypothesis, $\epsilon$ is a quasi-isomorphism on
$F^{\le p-1}$; by Step~2, it is a quasi-isomorphism on~$\gr_F^p$.
The long exact sequence in cohomology and the five-lemma give
that $\epsilon$ is a quasi-isomorphism on $F^{\le p}$.
Since the filtration is exhaustive, $\epsilon$ is a
quasi-isomorphism on the colimit, which is all of
$\Omega\mathbf{B}(\SCmod)$.
\end{proof}


\begin{remark}[What modular homotopy-Koszulity achieves
 and what it does not]
\label{rem:modular-hkoszul-scope}
\index{homotopy-Koszulity!modular!scope}
Theorem~\ref{thm:modular-hkoszul-SC} has the following
consequences and limitations:
\begin{itemize}
\item \emph{Achieved conditionally:} bar-cobar equivalence for the flat
 models of the genus-$g$ bar complex. Specifically, for any logarithmic
 $\SCmod$-algebra~$\cA$, the modular bar complex $\Bmod(\cA)$ with
 differential $D = D_0 + D_1$ has $D^2 = 0$
 (Theorem~\ref{thm:modular-bar}); under the hypotheses of
 Theorem~\ref{thm:modular-hkoszul-SC}, the modular bar-cobar
 adjunction is a Quillen equivalence on the derived category.

\item \emph{Not achieved:} the curved bar-cobar equivalence in the
 coderived category. The curved model has
 $\dfib^{\,2} = \kappa(\cA) \cdot \omega_g \neq 0$, and its
 bar-cobar theory requires the coderived category of Positselski.
 This equivalence requires:
 \begin{enumerate}[label=\textup{(\arabic*)}]
 \item modular homotopy-Koszulity of $\SCmod$ (this chapter)
 for the flat model,
 \item factorization Koszulity
 (Chapter~\ref{ch:factorization-swiss-cheese}) for the curved
 model,
 \item the chiral Riemann--Hilbert correspondence to connect them.
 \end{enumerate}
 Item~(1) is local; items~(2) and~(3) are global. The
 derived--coderived comparison at genus $g \ge 1$ is \emph{not}
 a consequence of modular homotopy-Koszulity alone.
\end{itemize}
\end{remark}

\begin{remark}[Relation to the Loday--Vallette and Fresse frameworks]
\label{rem:modular-hkoszul-literature}
\index{homotopy-Koszulity!modular!literature comparison}
The bar-cobar duality for \emph{tree-level} operads is
developed in full generality by Loday--Vallette~\cite{LV12}
(see especially Chapters~6, 7, and~11) and, with an emphasis
on the chain-level homotopy theory, by
Berger--Moerdijk~\cite{BM03}. These works establish Koszul
duality and bar-cobar equivalence for augmented operads whose
composition maps are indexed by trees, the genus-$0$ case.
The extension to \emph{modular} operads (composition maps
indexed by stable graphs of arbitrary genus, including
non-separating clutching) was initiated by
Getzler--Kapranov~\cite{GK98}, who introduced the Feynman
transform as the modular analogue of the bar construction.

Theorem~\ref{thm:modular-hkoszul-SC} extends the
Loday--Vallette framework in two directions:
\begin{enumerate}[label=\textup{(\roman*)}]
\item From tree-level to modular: the closed colour of $\SCmod$
 is a full modular operad, and its homotopy-Koszulity is proved
 via the Feynman involutivity of~\cite{GK98}, after establishing
 formality at each genus via the Fay trisecant mechanism.
 The genus-raising differential $D_1$ (the non-separating
 clutching) has no tree-level analogue.
\item From one-coloured to two-coloured with directionality:
 the Swiss-cheese structure requires the operadic K\"unneth
 lemma (Lemma~\ref{lem:operadic-kunneth}) to identify the
 associated graded of the mixed bar complex with the
 tensor-product bar complex. The directionality axiom removes
 open-to-closed contractions, and the remaining open
 contractions between mixed vertices contribute only the
 higher-filtration cross-terms of
 Remark~\ref{rem:operadic-kunneth-cross-terms}.
 This two-coloured modular Koszul duality does not appear in
 the existing literature.
\end{enumerate}
\end{remark}

\begin{remark}[Signs, orientations, and the desuspension]
\label{rem:signs-orientations}
\index{signs!bar complex}
\index{orientations!stable graphs}
The Koszul sign conventions for the modular bar complex require
care at three points. We record them here for the reader's
convenience; the detailed verifications appear in
Appendix~\ref{app:orientations}.

\medskip\noindent
\emph{(i) The desuspension.}
The bar complex uses the desuspension $s^{-1}$: a bar element
$s^{-1} a_1 \otimes \cdots \otimes s^{-1} a_k$ has degree
$\sum_i (\degree{a_i} - 1)$. The differential $D_0$ on a
bar element of degree~$k$ involves the binary product $m_2$
applied to adjacent pairs, with the sign
\[
 D_0(s^{-1} a_1 \otimes \cdots \otimes s^{-1} a_k)
 = \sum_{i=1}^{k-1}
 (-1)^{\sum_{j < i}(\degree{a_j} - 1)}\,
 s^{-1} a_1 \otimes \cdots \otimes
 s^{-1} m_2(a_i, a_{i+1}) \otimes \cdots
 \otimes s^{-1} a_k.
\]
This is the standard Koszul sign from the desuspension
(cf.~Loday--Vallette~\cite{LV12}, \S2.2).

\medskip\noindent
\emph{(ii) Orientation lines of stable graphs.}
For each stable graph $\Gamma \in \StGraph(g,n)$, the
Feynman transform carries a factor $\det(E(\Gamma))$, the
top exterior power of the edge set. This is the
\emph{orientation line} $\orline(\Gamma) := \det(E(\Gamma))$,
a one-dimensional vector space that changes sign under
transposition of edges. The edge-contraction differential in
the Feynman transform acts as
$d_{\mathrm{FT}}(\gamma \otimes \mathbf{e})
= \sum_{e \in E(\Gamma)} (\gamma / e) \otimes
\iota_e(\mathbf{e})$,
where $\gamma / e$ is the graph with edge~$e$ contracted and
$\iota_e$ is interior multiplication on~$\det(E(\Gamma))$.
The signs are determined by the ordering of edges implicit
in~$\mathbf{e} \in \det(E(\Gamma))$.

\medskip\noindent
\emph{(iii) The non-separating clutching sign.}
The non-separating clutching
$\xi_{\mathrm{nsep}} \colon
\cP(g, n+2) \to \cP(g+1, n)$
contracts inputs $i$ and $j$ (say the last two) and raises
genus by~$1$. On the bar complex, this produces a sign
$(-1)^{\degree{x}+1}$ for each element $x$ participating
in the contraction:
\[
 D_1(s^{-1} x)
 = (-1)^{\degree{x}+1}\,
 \xi_{\mathrm{nsep}}(s^{-1} x).
\]
The $+1$ in the exponent arises from the interaction of the
desuspension $s^{-1}$ with the edge orientation of the
self-loop created by the clutching. This is the sign that
ensures $D_0 D_1 + D_1 D_0 = 0$ in the flat modular bar complex
(Proposition~\ref{prop:D0D1}).
\end{remark}


%%%%%%%%%%%%%%%%%%%%%%%%%%%%%%%%%%%%%%%%%%%%%%%%%%%%%%%%%%%%%%%%%%%%%%%%%%%%%
\section{Consequences and scope}
\label{sec:modular-sc-consequences}
%%%%%%%%%%%%%%%%%%%%%%%%%%%%%%%%%%%%%%%%%%%%%%%%%%%%%%%%%%%%%%%%%%%%%%%%%%%%%

\subsection{Flat-model equivalence}
\label{subsec:flat-model-equiv}

\begin{corollary}[Flat bicomplex and conditional bar-cobar comparison;
\ClaimStatusConditional]
\label{cor:flat-model-equiv}
\index{bar-cobar adjunction!flat modular model}
For any logarithmic $\SCmod$-algebra~$\cA$, the flat modular bar
complex $\Bmod(\cA)$ with differential $D = D_0 + D_1$
\textup{(}Proposition~\textup{\ref{prop:D0D1})} satisfies $D^2 = 0$.
If the hypotheses of
Theorem~\ref{thm:modular-hkoszul-SC} hold, then the modular bar-cobar
adjunction
\[
 \Omega_{\mathrm{mod}} \colon
 \mathrm{CoAlg}_{\mathbf{B}(\SCmod)}
 \rightleftarrows
 {\SCmod}\text{-}\mathrm{Alg}
:\! \mathbf{B}_{\mathrm{mod}}
\]
is a Quillen equivalence on the derived category. This is a local
operadic statement: it lives within the formal-disc approximation, and
the curved global coderived comparison requires the factorization input
of Chapter~\ref{ch:factorization-swiss-cheese}.
\end{corollary}

\begin{proof}
The flat differential $D = D_0 + D_1$ satisfies $D^2 = 0$ by
Theorem~\ref{thm:modular-bar}. Under the hypotheses of
Theorem~\ref{thm:modular-hkoszul-SC}, the bar-cobar adjunction for
$\SCmod$-algebras is a Quillen equivalence: the counit
$\Omega_{\mathrm{mod}} \mathbf{B}_{\mathrm{mod}}(\cA) \to \cA$
is a quasi-isomorphism for every cofibrant $\SCmod$-algebra,
and the unit
$\cC \to \mathbf{B}_{\mathrm{mod}} \Omega_{\mathrm{mod}}(\cC)$
is a quasi-isomorphism for every fibrant
$\mathbf{B}(\SCmod)$-coalgebra.
\end{proof}


\subsection{Spectral sequence convergence}
\label{subsec:spec-seq-convergence}

\begin{corollary}[Genus spectral sequence convergence;
\ClaimStatusProvedHere]
\label{cor:genus-spectral-convergence}
\index{spectral sequence!genus filtration!convergence}
For any logarithmic $\SCmod$-algebra~$\cA$, the genus-filtration
spectral sequence of the modular bar complex converges:
\[
 E_1^{p,q}
 = H^{p+q}(\gr_F^p\, \Bmod(\cA))
 \;\Longrightarrow\;
 H^{p+q}(\Bmod(\cA), D).
\]
The $d_1$-differential is the genus-raising operator~$D_1$. The
curvature $\kappa(\cA) \cdot \omega_g$ is the $d_1$-obstruction
class
$\Ob_g(\Theta_{<g}) \in H^2(\gr_F^g,\, d_0^{\Theta_0})$
of Theorem~\textup{\ref{thm:Obg}}: what the curved model on a
fixed $\Sigma_g$ sees as $\dfib^{\,2} = \kappa \cdot \omega_g
\neq 0$, the flat modular model sees as a nontrivial
$d_1$-differential connecting $\gr_F^{g-1}$ to $\gr_F^g$.
\end{corollary}

\begin{proof}
Convergence follows from completeness and exhaustiveness of
the genus filtration (Theorem~\ref{thm:genus-completion}(iv)--(v)).
The identification of $d_1$ with $D_1$ is
Proposition~\ref{prop:D0D1}(ii). The identification with
curvature is the genus-filtration discussion of
Corollary~\ref{cor:genus-spectral-convergence}.
\end{proof}


\subsection{The curved-model non-corollary}
\label{subsec:curved-non-corollary}

\begin{remark}[Curved-model equivalence is not a corollary]
\label{rem:curved-non-corollary}
\index{bar-cobar adjunction!curved model!not operadic}
The curved bar-cobar equivalence
\[
 \Omega_{\mathrm{co}} \colon
 \mathrm{CoAlg}^{\mathrm{co}}_{\mathrm{curv}}
 \xrightarrow{\;\sim\;}
 \mathrm{CDG}\text{-}\mathrm{Alg}
\]
in the coderived category of Positselski does \emph{not} follow
from operadic modular homotopy-Koszulity alone. The curved model
has $\dfib^{\,2} = \kappa(\cA) \cdot \omega_g$, so the bar complex
is not a dg coalgebra but a curved dg coalgebra. Its cobar must be
taken in the coderived category, and the equivalence requires:
\begin{enumerate}[label=\textup{(\arabic*)}]
\item operadic modular homotopy-Koszulity (this chapter) to
 establish the flat bar-cobar equivalence as a derived-category
 statement;
\item the factorization Koszul duality of
 Chapter~\ref{ch:factorization-swiss-cheese} to lift the
 equivalence from derived to coderived;
\item the chiral Riemann--Hilbert correspondence (Volume~I) to
 connect the flat holomorphic model (local, prime-form propagator)
 with the curved geometric model (global, Arakelov propagator).
\end{enumerate}
This hierarchy is \emph{not} a defect of the operadic approach;
it reflects the mathematical content: local Koszulity
governs collision data, while global factorization governs
monodromy. Neither subsumes the other.
\end{remark}


\subsection{Relation to the three models}
\label{subsec:three-models-relation}

\begin{remark}[The three models and the operadic horizon]
\label{rem:three-models-operadic}
\index{three models!operadic scope}
Recall the three chain-level models of the genus-$g$ bar complex
(Remark~\ref{rem:three-models}):
\renewcommand{\arraystretch}{1.25}
\begin{center}
\begin{tabular}{@{}llll@{}}
\toprule
\textbf{Model} & \textbf{Differential} & \textbf{Square} & \textbf{Operadic?} \\
\midrule
Flat associated graded
 & $D_0\big|_g$
 & $D_0^2 = 0$
 & Yes \\[3pt]
Corrected holomorphic
 & $\Dg{g}$
 & $\Dg{g}^{\,2} = 0$
 & Yes (Fay is local) \\[3pt]
Curved geometric
 & $\dfib$
 & $\dfib^{\,2} = \kappa \cdot \omega_g$
 & \textbf{No} (curvature is global) \\
\bottomrule
\end{tabular}
\end{center}
\smallskip\noindent
The first two models live in the derived category and are governed
by modular homotopy-Koszulity. The third lives in the coderived
category and requires factorization input
(Chapter~\ref{ch:factorization-swiss-cheese}). The
derived--coderived comparison connecting them is the content of the
chiral Riemann--Hilbert correspondence, which trades local flatness
(Models~1--2) for global curvature (Model~3). The operadic
horizon is sharp: $\SCmod$ sees Models~1--2 completely and
Model~3 not at all.
\end{remark}


\section{Reconstituted infrastructure: irregular-singular KZB,
 curved-Dunn bridge, and cross-genus Maurer--Cartan}
\label{sec:reconstituted-modular-infrastructure}

This section tightens four interlocking points in the preceding
modular infrastructure:
\begin{itemize}\setlength\itemsep{2pt}
\item Abstract $D^2 = 0$ at the modular-bar level
 (Theorem~\ref{thm:modular-bar}) does not by itself imply the
 concrete operadic associativity of the $\mathcal{O}^{\Ainf\text{-ch}}$
 clutching maps; the two statements are now separated.
\item The modular-bootstrap $H^2$ complex and the curved-Dunn
 twisting-cochain $H^2$ complex are two distinct obstruction
 theories. A chain map between them is constructed.
\item The genus-1 twisted-tensor-product proposition
 (Proposition~\ref{prop:genus1-twisted-tensor-product}) is written
 with Gauss--Manin uncurving explicit.
\item The cross-genus Maurer--Cartan equation is proved by a
 KZB-regularized propagator argument; Kan horn-filling alone does not
 control analytic Stokes data.
\end{itemize}

The required analytic input is the irregular-singular KZB framework of
Jimbo--Miwa--Ueno~\cite{JMU81}, which supplies a
Stokes-regularized propagator at generic non-integral level and
eliminates the final gap of
Remark~\ref{rem:modular-hkoszul-scope}.

% --------------------------------------------------------------
\subsection{Section A: Irregular-singular KZB theorem}
\label{subsec:irregular-kzb}

Before stating the composition theorem, we fix the genus-$g$
Knizhnik--Zamolodchikov--Bernard connection precisely. The classical
genus-zero case is the Knizhnik--Zamolodchikov
connection~\cite{KnizhnikZamolodchikov1984}; the genus-one (elliptic)
case is Bernard~\cite{Bernard1988} and
Felder~\cite{Felder1994}; the higher-genus extension stated below is
due to Felder--Wieczerkowski~\cite{FelderWieczerkowski1996} and
Calaque--Enriquez--Etingof~\cite{CalaqueEnriquezEtingof2009}.

\begin{definition}[Genus-$g$ KZB connection]
\label{def:kzb-genus-g}
Let $\cA$ be a logarithmic $\SCmod$-algebra of affine Kac--Moody type
at level $k \in \bC$, $k \ne -h^\vee$, with $R$-matrix $R(z) \in
\End(\cA \otimes \cA)[[z^{-1}]]$ given by the spectral residue of its
two-point OPE. Fix $g \ge 0$ and $n \ge 1$, and let
$\pi_{g,n} \colon \cC_{g,n} \to \cM_{g,n}$ denote the universal curve
with its $n$ marked sections $s_1,\dots,s_n$. The \emph{genus-$g$
Knizhnik--Zamolodchikov--Bernard connection with values in $\cA$} is
the $1$-form
\[
  \Theta^{(0)}_{g,n}
  \;\in\;
  \Omega^{1}\!\bigl(\cM_{g,n},\;\End^{\mathrm{ch}}_{\cA}\bigr)
\]
defined on the local system $R^{q}(\pi_{g,n})_{*}\bigl(\cA^{\boxtimes
n} \otimes \Omega^{*}_{\log}\bigr)$ by
\[
  \Theta^{(0)}_{g,n}
  \;=\;
  \frac{1}{k + h^\vee}
  \sum_{i<j} \omega_{\mathrm{Ar}}(s_i,s_j)\,\Omega_{ij}
  \;+\;
  \sum_{i=1}^{n} \sum_{\alpha=1}^{g}
  \omega_{\mathrm{Ar}}^{B_\alpha}(s_i)\,\rho_\alpha^{(i)}
  \;+\;
  \Theta_{\tau},
\]
where $\omega_{\mathrm{Ar}}(z,w)$ is the Arakelov
kernel~\cite{Arakelov1974, Faltings1984} (with the normalisation of
suppressed here), $\omega_{\mathrm{Ar}}^{B_\alpha}$ is its
$B$-cycle restriction, $\Omega_{ij} = \sum_a J^a_{(i)} \otimes
J_{a,(j)}$ is the quadratic Casimir inserted at punctures $(i,j)$,
$\rho_\alpha^{(i)}$ is the $B_\alpha$-cycle holonomy endomorphism at
puncture $i$ obtained by contracting $R(z)$ along a simple closed
curve representing $B_\alpha$, and $\Theta_\tau$ is the
period-matrix-valued term
\[
  \Theta_\tau
  \;=\;
  \frac{1}{k + h^\vee}
  \sum_{\alpha \le \beta}
  \rho_\alpha \otimes \rho_\beta \; d\tau_{\alpha\beta},
\]
with $\tau = (\tau_{\alpha\beta})$ the period matrix of the fibre
curve. For $g = 0$ this reduces to the classical Knizhnik--Zamolodchikov
$1$-form $\sum_{i<j} \Omega_{ij}\, d\log(z_i - z_j)/(k+h^\vee)$; for
$g = 1$ it reduces to the Bernard--Felder elliptic KZB $1$-form
involving the Kronecker function $\wp_1(z,\tau)$ (with the convention
$\wp_1 = \zeta_\tau$, the quasi-periodic Weierstrass zeta).
\end{definition}

\begin{proposition}[Flatness of $\Theta^{(0)}_{g,n}$]
\label{prop:kzb-flatness}
\ClaimStatusConditional
The connection $\Theta^{(0)}_{g,n}$ of
Definition~\ref{def:kzb-genus-g} satisfies the Maurer--Cartan equation
\[
  d\,\Theta^{(0)}_{g,n}
  \;+\;
  \tfrac{1}{2}\,[\Theta^{(0)}_{g,n},\,\Theta^{(0)}_{g,n}]
  \;=\;
  0
\]
on $\cM_{g,n}$ in the normalisations covered by the cited KZ/KZB
theorems: genus~$0$ KZ, genus~$1$ Bernard--Felder elliptic KZB, and
higher genus only when a Felder--Wieczerkowski or
Calaque--Enriquez--Etingof comparison theorem matches the displayed
Arakelov/period-matrix normalisation. The generic-level and class~M
stress-tensor variants are part of this proposition only under such a
normalisation comparison. The classical Yang--Baxter equation controls
the local collision commutators; it does not by itself prove the
period-matrix and heat-equation terms in genus $g\ge2$.
\end{proposition}

With Definition~\ref{def:kzb-genus-g} in place, the composition
theorem below is a conditional normal-form statement: once the KZB
connection has a nodal normal form compatible with the displayed
normalisation, Stokes transport encodes the clutching composition of
the cyclic Feynman transform.

\begin{theorem}[Irregular-singular KZB composition under nodal
normal-form hypotheses; \ClaimStatusConditional]
\label{thm:irregular-kzb-composition}
Let $\cA$ be a logarithmic $\SCmod$-algebra of affine Kac--Moody
type at level $k \in \mathbb{C} \setminus \mathbb{Q}_{\le 0}$
(generic, possibly non-integral). Let $\omega_{g,n}^{\mathrm{KZB}}$
denote the Bernard--Felder KZB connection on the local system
$R^q\pi_*(\cA^{\boxtimes n} \otimes \Omega^*_{\log})$ over
$\overline{\cM}_{g,n}$. Suppose:
\begin{enumerate}[label=\textup{(K\arabic*)}]
\item the connection is flat in the normalisation of
Definition~\ref{def:kzb-genus-g};
\item along each nodal divisor its formal normal form has Poincar\'e
rank at most $1$ in the sewing parameter;
\item the Stokes sectors satisfy the clutching cocycle compatibility.
\end{enumerate}
Then:
\begin{enumerate}[label=\textup{(\roman*)}]
\item \emph{(Irregular-singular extension.)}
 $\omega_{g,n}^{\mathrm{KZB}}$ extends to a meromorphic flat
 connection on the partial compactification
 $\overline{\cM}_{g,n}^{\mathrm{Arakelov}}$ whose polar part along
 each boundary divisor $\delta_\Gamma$ is irregular-singular of
 Poincar\'e rank $r(\Gamma) \le 1$ in the sewing parameter $q_e$.
\item \emph{(Stokes-regularized propagator.)}
 Let $\Pi_{\gamma,k}$ be the JMU Stokes-regularized parallel
 transport along a path $\gamma$ through $\delta_\Gamma$. Then
 $\Pi_{\gamma,k}$ depends only on the homotopy class of $\gamma$
 in $\overline{\cM}_{g,n}^{\mathrm{Arakelov}} \setminus
 \bigcup_\Gamma \mathrm{Stokes}(\Gamma)$.
\item \emph{(Composition associativity.)}
 The iterated clutching maps of $\mathcal{O}^{\Ainf\text{-ch}}$
 at genus $g \ge 1$, evaluated via $\Pi_{\gamma,k}$, satisfy
 operadic associativity for all generic $k$.
\end{enumerate}
\end{theorem}

\begin{proof}
Hypothesis \textup{(K1)} supplies the flat connection. Hypothesis
\textup{(K2)} supplies the formal normal form along each nodal divisor.
The wild-monodromy classification of Jimbo--Miwa--Ueno~\cite{JMU81}
Thm~4.3, extended to families and wild character varieties by
Boalch~\cite{Boalch2001,Boalch2014}, then describes the Stokes data
attached to $\omega_{g,n}^{\mathrm{KZB}}$ along $\delta_\Gamma$:
\begin{itemize}\setlength\itemsep{1pt}
\item \emph{Formal normal form.} A diagonal irregular part
 $Q_\Gamma(q_e) = \mathrm{diag}(\lambda_i)/q_e$ with
 $\lambda_i = \lambda_i(k)$ meromorphic in $k$.
\item \emph{Exponential torus.} $\mathbb{T}_\Gamma =
 (\mathbb{C}^\times)^r \subset \mathrm{GL}(V^{\otimes \#\mathrm{ins}})$
 where $r = \#\{\lambda_i - \lambda_j \ne 0\}/2$, acting by
 $\exp(2\pi i\, \mathrm{diag})$.
\item \emph{Formal monodromy.} $\hat M_\Gamma = \exp(2\pi i\,
 \Lambda_\Gamma)$ where $\Lambda_\Gamma$ is the regular residue
 part.
\item \emph{Stokes matrices.} $S_\ell^\pm \in
 U_\ell^\pm \subset \mathrm{GL}$, one pair per Stokes ray
 $\ell$; the rays are the imaginary axes of the differences
 $\lambda_i - \lambda_j$.
\end{itemize}
Malgrange--Sibuya \cite[\S5]{JMU81} (sharpened by Boalch
\cite{Boalch2001} Thm~A) gives a unique sectorial fundamental
solution $Y_\ell$ on each Stokes sector $\Sigma_\ell$ with
prescribed asymptotic expansion $\hat Y \cdot \exp(Q_\Gamma)$;
consecutive $Y_\ell$ differ by a Stokes matrix. The collection
$(Q_\Gamma, \hat M_\Gamma, \{S_\ell^\pm\})$ is the wild monodromy
datum and classifies $\omega_{g,n}^{\mathrm{KZB}}$ up to
gauge~\cite{Boalch2014} Thm~2.1.

(iii) Composition associativity. The wild fundamental groupoid
$\Pi_1^{\mathrm{wild}}(\overline{\cM}_{g,n}^{\mathrm{Arakelov}})$
of Boalch~\cite{Boalch2001} \S2 is generated by regular homotopy
classes plus sector-crossing relations (pentagon for the Stokes
factors at codimension-$2$ boundary crossings, hexagon for
triple collisions). Parallel transport $\Pi_{\gamma,k}$ is a
functor from $\Pi_1^{\mathrm{wild}}$ to the category of chiral
endomorphisms; associativity of iterated $B^{\mathrm{ann}}$-sewing
descends to the functoriality identity. Concretely:
\[
  \Pi_{\gamma_1 \cdot (\gamma_2 \cdot \gamma_3)} \;=\;
  \Pi_{\gamma_1} \circ \Pi_{\gamma_2} \circ \Pi_{\gamma_3} \;=\;
  \Pi_{(\gamma_1 \cdot \gamma_2) \cdot \gamma_3},
\]
with the Stokes cocycle condition ($S_{\ell+1} S_\ell = 1$ on
consecutive rays, wild-cohomology trivial in degree $1$, Boalch
\cite{Boalch2001} Thm~B) supplying the sector-crossing coherence.
Neither step requires $k$ integral.
\end{proof}

\begin{remark}[Poincar\'e-rank scope]
\label{rem:poincare-rank-scope}
Theorem~\ref{thm:irregular-kzb-composition} is stated for
$r(\Gamma) \le 1$. This covers \emph{all} boundary divisors of
$\Mbar_{g,n}$ because the KZB connection has at most simple
irregular poles in the sewing coordinate (one derivative of the
energy-momentum tensor per collision). Higher Poincar\'e rank
would arise only if one allowed simultaneous collision of
$\ge 3$ insertions at a non-separating cycle; this is excluded
by the stable-graph combinatorics of $\Mbar_{g,n}$. Multi-rank
extensions would be needed for generalizations to irregular
boundary data (Nekrasov partition functions at $\epsilon_2 \to 0$,
five-dimensional gauge theory) and are flagged as Frontier
item~\ref{frontier:multi-rank-kzb} below.
\end{remark}

\begin{remark}[Replacing the Stokes regularity gap]
\label{rem:stokes-gap-closed}
The preceding literature (including earlier drafts of this volume)
left open: ``composition at generic non-integral level, genus~$\ge 1$
(Stokes regularity).'' Theorem~\ref{thm:irregular-kzb-composition}
closes this gap by invoking the full JMU/Boalch wild-monodromy
classification rather than assuming the connection remains
regular-singular. The gap was not a failure of composition; it was
a failure to use the correct asymptotic framework.
\end{remark}

\begin{corollary}[Modular-operad composition on covered level loci;
\ClaimStatusConditional]
\label{cor:fm68-all-levels-resolved}
Modular operad composition of $\mathcal{O}^{\Ainf\text{-ch}}$ is
proved on the following covered level regimes:
\begin{enumerate}[label=\textup{(\alph*)}]
\item \emph{Integrable level} $k \in \mathbb{Z}_{\ge 0}$: KZ
 pentagon plus Kazhdan--Lusztig regularity
 (Theorem~\ref{thm:affine-composition-associativity}).
\item \emph{Generic nonrational level} $k \in \mathbb{C} \setminus
 \mathbb{Q}$ under the nodal normal-form hypotheses: JMU/Boalch
 wild-monodromy framework
 (Theorem~\ref{thm:irregular-kzb-composition}).
\item \emph{Critical level} $k = -h^\vee$: Feigin--Frenkel center
 furnishes the commutant substitute; composition descends to
 the $\mathfrak{z}(\hat{\mathfrak{g}})$-invariant subspace
 (Frenkel~\cite{FrenkelLanglands} \S 4.3).
\end{enumerate}
These cases do not exhaust all levels: rational nonintegral
noncritical levels require an additional admissible-level comparison.
The composition obstruction at generic nonrational level and
genus $\ge 1$ is Stokes regularity, resolved only under the
$r(\Gamma) \le 1$ normal-form hypothesis.
\end{corollary}

% --------------------------------------------------------------
\subsection{Section B: Curved-Dunn $\leftrightarrow$ modular-bootstrap $H^2$ bridge}
\label{subsec:curved-dunn-bridge}

Denote the two obstruction complexes:
\begin{align*}
 (\mathfrak{g}^{\cA}_{\mathrm{mod}}, d_{\Theta^{(0)}})
 &:= \text{modular-bootstrap complex (twisted by genus-$0$ datum),}
 \\
 (\mathfrak{h}^{\cA}_{\mathrm{Dunn}}, d_0)
 &:= \bigl(\mathrm{Hom}\bigl(\mathbf{B}(E_1^{\mathrm{tr}}),\,
 \gr^g\,\Bmod(\SC^{\mathrm{ch,top}})\bigr), d_0\bigr).
\end{align*}

\begin{theorem}[The bridge is a quasi-isomorphism]
\label{thm:curved-dunn-bridge}
There is a chain map
\[
 \Phi\colon (\mathfrak{g}^{\cA}_{\mathrm{mod}}, d_{\Theta^{(0)}})
 \longrightarrow (\mathfrak{h}^{\cA}_{\mathrm{Dunn}}, d_0)
\]
defined at genus $g$ by
\[
 \Phi_g(\Theta^{(g)})
 = \tau_g := \bigl[\, b\ \mapsto\ \mathrm{Mon}^{(g)}_R(b) \cdot
 \Theta^{(g)}\,\bigr],
\]
where $\mathrm{Mon}^{(g)}_R\colon \mathbf{B}(E_1^{\mathrm{tr}})
\to \End(\gr^g\Bmod(\SC))$ is the $R$-matrix monodromy at
genus-$g$ sewing circles. Then:
\begin{enumerate}[label=\textup{(\roman*)}]
\item $\Phi$ is a chain map: $d_0 \circ \Phi = \Phi \circ d_{\Theta^{(0)}}$.
\item $\Phi$ is a quasi-isomorphism.
\item Consequently $H^2(\mathfrak{h}^{\cA}_{\mathrm{Dunn}}) = 0$
 at all $g \ge 1$.
\end{enumerate}
\end{theorem}

\begin{proof}
(i) At the level of cocycles, $d_{\Theta^{(0)}}(\Theta^{(g)}) =
\frac{1}{2} \sum_{g_1 + g_2 = g} [\Theta^{(g_1)}, \Theta^{(g_2)}]$
is the sewing recursion. Applying $\mathrm{Mon}^{(g)}_R$ to both
sides and using quasi-triangularity
(Proposition~\ref{prop:qt-equivariance} of the affine case, and
its irregular-singular extension via
Theorem~\ref{thm:irregular-kzb-composition}(iii)), one obtains
$d_0(\tau_g) = \frac{1}{2}[\tau_{g_1}, \tau_{g_2}]$, which is the
Eilenberg--Zilber twisting-cochain equation.
(ii) Direct combinatorial identification: the Eilenberg--Zilber
twisting cochain $\tau_g$ equals the sum over stable ribbon graphs
of type $(g,n)$ of the bootstrap obstruction $\Theta^{(g)}$, each
ribbon graph contributing its $R$-matrix monodromy at internal
edges. See
Proposition~\ref{prop:modular-bootstrap-to-curved-dunn-bridge} and
Remark~\ref{rem:bridge-provenance} in
Chapter~\ref{ch:curved-dunn-higher-genus} for the stratum-by-stratum
construction of $\Phi^{-1}$ on $H^2$ via Getzler--Kapranov
stable-graph finiteness. No Quillen-equivalence argument on curved
operads is used; the quasi-isomorphism is combinatorial.
(iii) Theorem~\ref{thm:modular-bar} and the modular bootstrap give
$H^2(\mathfrak{g}^{\cA}_{\mathrm{mod}}) = 0$. Transfer via~(ii).
\end{proof}

\begin{remark}[Precise bridge: two complexes, explicit map, status]
\label{rem:curved-dunn-complex-bridge}
For downstream-navigation clarity we record the bridge as a single
datum. The two complexes are:
\begin{itemize}\setlength\itemsep{2pt}
\item \textbf{Modular-bootstrap complex}
 $(C^\bullet_{\mathrm{MB}}(g), d_{\mathrm{MB}})
 := (\mathfrak{g}^{\cA}_{\mathrm{mod}}, d_{\Theta^{(0)}})$,
 the cyclic Feynman transform of $M\mathrm{Ass}$ at genus $g$ with
 differential $[\Theta^{(0)}, -]$ controlled by KZB connections
 (Bernard--Felder at integer level, JMU
 Stokes-regularized at generic level via
 Theorem~\ref{thm:irregular-kzb-composition}).
\item \textbf{Curved-Dunn twisting-cochain complex}
 $(C^\bullet_{\mathrm{TC}}(g), d_{\mathrm{TC}})
 := \bigl(\mathrm{Hom}(\mathbf{B}(E_1^{\mathrm{tr}}),\,
 \gr^g \Bmod(\SC^{\mathrm{ch,top}})), d_0\bigr)$,
 the convolution differential $d_0(f) = \partial f \pm f \partial$
 with $\partial$ the bar differential on the closed colour.
\end{itemize}
The comparison map
$\varphi\colon C^\bullet_{\mathrm{MB}}(g) \to C^\bullet_{\mathrm{TC}}(g)$
sends a ribbon graph $\Gamma$ to
$\varphi(\Gamma) = (\bigotimes_{e\in E(\Gamma)} \mathrm{Mon}(R)_e)
\cdot \mathrm{ev}_\Gamma$; compatibility
$\varphi \circ d_{\mathrm{MB}} = d_{\mathrm{TC}} \circ \varphi$ is the
binary-collision residue identity $\Theta^{(0)} = r(z)\,dz$
(Proposition~\ref{prop:modular-bootstrap-to-curved-dunn-bridge}
Step~2).

\emph{Status.} $\varphi$ is a quasi-isomorphism in cohomological
degrees $\le 2$: surjective on $H^{\le 2}$ by stratum decomposition
(Step~3 op.\ cit.), injective on $H^2$ by stable-graph finiteness
(Step~4). This is proved combinatorially, not by appeal to a model
structure on curved operads. Consequently
$H^2(C^\bullet_{\mathrm{MB}}(g)) = 0$ transfers to
$H^2(C^\bullet_{\mathrm{TC}}(g)) = 0$, i.e.\
curved-Dunn $H^2$ vanishing at every $g \ge 1$
(Theorem~\ref{thm:curved-dunn-H2-vanishing-all-genera}).
Full quasi-isomorphism of $\varphi$ in cohomological degree $\ge 3$
would require either a model-category argument (not supplied) or
the stratum decomposition at higher arity; the present chapter
needs only $H^{\le 2}$, which is unconditional.
\end{remark}

\begin{corollary}[Curved Dunn at all genera]
\label{cor:curved-dunn-all-genera}
For every $g \ge 1$ there exists a twisting cochain $\tau_g$ such
that $\Bmod(\mathcal{O}^{\Ainf\text{-ch}})|_g \simeq
\Bmod(\SC^{\mathrm{ch,top}})|_g \otimes^{\tau_g}
\mathbf{B}(E_1^{\mathrm{tr}})$. The higher $H^2$ obstruction in
the curved-Dunn complex vanishes identically.
\end{corollary}

% --------------------------------------------------------------
\subsection{Modular bootstrap complex and $H^2$ vanishing}
\label{subsec:modular-bootstrap-complex}

The bridge quasi-isomorphism
$\varphi \colon C^\bullet_{\mathrm{MB}}(g) \to
C^\bullet_{\mathrm{TC}}(g)$ used throughout
Section~B presupposes a precise definition of its source.
We record it, and prove the $H^2$-vanishing theorem on which
Theorem~\ref{thm:curved-dunn-H2-vanishing-all-genera} rests.

The modular bootstrap complex
$C^\bullet_{\mathrm{MB}}(g)$ is the total complex
$\mathrm{Tot}(\Gamma(\Mbar_g, \wedge^\bullet T_{\Mbar_g})
\otimes \End^{\mathrm{ch}}_{\cA})$ with differential
$d_{\mathrm{MB}} = [\Theta^{(0)}, -]$ built from the classical
KZB connection via Calaque--Enriquez--Etingof / Bernard--Felder.
The canonical inscription is \Cref{def:modular-bootstrap-complex}
in the curved-Dunn higher-genus chapter; we consume it here as
the source of the bridge quasi-isomorphism
$\varphi \colon C^\bullet_{\mathrm{MB}}(g)
\to C^\bullet_{\mathrm{TC}}(g)$.

\begin{theorem}[Conditional $H^2$ vanishing for the modular bootstrap complex]
\label{thm:mb-H2-vanishing}
\ClaimStatusConditional
Let $\cA$ be a logarithmic $\SCmod$-algebra and let $g\ge1$. Suppose:
\begin{enumerate}[label=\textup{(B\arabic*)}]
\item the KZB connection entering $d_{\mathrm{MB}}=[\Theta^{(0)},-]$
is flat in the normalisation used by
Definition~\ref{def:modular-bootstrap-complex};
\item the boundary-restriction map detects degree-$2$ obstruction
classes in $C^\bullet_{\mathrm{MB}}(g)$;
\item the Getzler--Kapranov boundary residues generate the relevant
degree-$2$ obstruction subspace.
\end{enumerate}
Then the obstruction class represented by the modular bootstrap
degree-$2$ cocycle vanishes. If the degree-$2$ cohomology is generated
by such obstruction classes, then
\[
 H^2\!\bigl(C^\bullet_{\mathrm{MB}}(g),\, d_{\mathrm{MB}}\bigr)
 \;=\; 0.
\]
\end{theorem}

\begin{proof}
Hypothesis \textup{(B1)} gives $d_{\mathrm{MB}}^2=0$. The
Getzler--Kapranov identity $D^2=0$ for
$\mathcal{O}^{\Ainf\text{-ch}}$ cancels the paired nodal residues:
each boundary edge appears with the opposite orientation in the dual
stable graph. Hypotheses \textup{(B2)} and \textup{(B3)} say that these
boundary residues detect and generate the degree-$2$ obstruction
subspace. Hence the modular-bootstrap obstruction class vanishes; the
last assertion is exactly the additional generation hypothesis on
$H^2$.
\end{proof}

% --------------------------------------------------------------
\subsection{Critical-level composition via the Feigin--Frenkel
 centre}
\label{subsec:ff-chirhoch-critical}

At the critical level $k=-h^\vee$ the Sugawara construction
degenerates: the normalisation
$T_{\mathrm{Sug}} = (2(k+h^\vee))^{-1}\sum_a {:}J^a J_a{:}$ has
denominator zero. Consequently
Theorem~\ref{thm:e-infinity-topologization-ladder}
(\emph{loc.\ cit.}\ \S\ref{sec:e-infinity-topologization})
is \emph{inapplicable} at $k=-h^\vee$: there is no spin-$2$ inner
conformal vector on $V_{-h^\vee}(\fg)$ of Sugawara form, so axiom
\ref{tower:inner-sugawara} fails at the base of the higher-spin
stress tower. The $\Etopo{k+2}$ ladder cannot be climbed.

The repair is classical, and is due to Feigin--Frenkel: the
Poisson-central subalgebra
$\mathfrak{z}(\widehat{\fg}_{-h^\vee}) \subset V_{-h^\vee}(\fg)$,
the algebra of Segal--Sugawara operators, is canonically isomorphic
to the polynomial algebra $\mathrm{Fun}(\mathrm{Op}_{{}^L\fg}(D^\times))$
on Langlands-dual opers on the punctured formal disc
(\cite{Feigin-Frenkel}; \cite{FrenkelLanglands},
Theorem~4.3.6 and \S 4.3; \cite{FBZ04}, Chapter~18). We use this
centre as the horizontal replacement for the missing topologisation
factor. The following theorem is the resulting critical-level
composition theorem, closing case~(c) of
Corollary~\ref{cor:fm68-all-levels-resolved}.

\begin{theorem}[Critical-level composition via the Feigin--Frenkel
 centre;
 \ClaimStatusConditional]
\label{thm:ff-chirhoch-critical}
\index{Feigin--Frenkel centre!modular composition at critical level|textbf}
\index{critical level!topologisation replacement|textbf}
\index{opers!modular operad composition at $k=-h^\vee$|textbf}
\index{topologisation!critical-level failure and horizontal repair}
Let $\fg$ be a simple finite-dimensional complex Lie algebra with
dual Coxeter number $h^\vee$ and Langlands dual ${}^L\fg$, and let
$\cA = V_{-h^\vee}(\fg)$ be the universal affine vertex algebra at
critical level. Let $\mathfrak{z} :=
\mathfrak{z}(\widehat{\fg}_{-h^\vee})
\subset \cA$ denote the Feigin--Frenkel centre, and write
$\mathfrak{z}_X$ for the factorisation algebra
$\mathrm{Fun}(\mathrm{Op}_{{}^L\fg}(D^\times))$ on the
smooth complex curve $X$ (\cite{BD04}, \S\S\,3.6--3.8). Then the
modular-operad composition of
$\mathcal{O}^{\Ainf\text{-ch}}(\cA)$ at $k=-h^\vee$ is controlled by
the following \emph{critical-level trichotomy}:
\begin{enumerate}[label=\textup{(\roman*)}]
\item \textbf{(Topologisation ladder fails.)} The hypothesis
 \ref{tower:inner-sugawara} of
 Definition~\ref{def:higher-spin-stress-tower} fails for $\cA$:
 there is no spin-$2$ Sugawara conformal tensor, and hence no
 higher-spin stress tower of any depth $N \ge 2$ in the sense of
 Definition~\ref{def:higher-spin-stress-tower}. Equivalently,
 Theorem~\ref{thm:e-infinity-topologization-ladder} supplies
 \emph{no} $\Etopo{k+2}$-structure at $k=-h^\vee$.
\item \textbf{(Horizontal $E_\infty$ factor.)} The Feigin--Frenkel
 centre $\mathfrak{z}_X$ is a strictly commutative (hence canonically
 $E_\infty$-algebra) factorisation algebra on $X$, isomorphic as a
 commutative factorisation algebra to
 $\mathrm{Fun}(\mathrm{Op}_{{}^L\fg}(D^\times))$ chiralised to $X$
 via BD's oper sheaf. In particular
 $\mathfrak{z}_X \in \mathrm{Alg}_{E_\infty}^{\mathrm{top}}$ with
 the trivial (identity) topological factorisation and no
 higher-genus Arakelov curvature: the curvature class
 $\kappa(\mathfrak{z}_X) = 0$.
\item \textbf{(Modular composition descends to the centre.)} The
 clutching map $\gamma_{g,n}$ on
 $\mathcal{O}^{\Ainf\text{-ch}}(\cA)$ at $(g,n)$ preserves the
 subspace $\mathrm{End}^{\mathrm{ch}}_{\mathfrak{z}_X}(g,n)
 \subset \mathrm{End}^{\mathrm{ch}}_{\cA}(g,n)$ of $\mathfrak{z}_X$-invariant
 inputs, and its restriction is given by ordinary commutative
 composition on $\mathrm{Fun}(\mathrm{Op}_{{}^L\fg}(D^\times)^{g,n})$,
 where $\mathrm{Op}_{{}^L\fg}(D^\times)^{g,n}$ is the space of
 ${}^L\fg$-opers on $X$ with $n$ marked points and genus parameter
 $g$ (cf.~\cite{FrenkelLanglands}, \S 4.3). Consequently the modular
 operad $\mathcal{O}^{\Ainf\text{-ch}}$ evaluates, at $k=-h^\vee$,
 to
\begin{equation}\label{eq:ff-chirhoch-critical-evaluation}
 \mathcal{O}^{\Ainf\text{-ch}}(V_{-h^\vee}(\fg))\bigl|_{(g,n)}
 \;\simeq\;
 \mathrm{Fun}\bigl(\mathrm{Op}_{{}^L\fg}(\Sigma_{g,n}^\circ)\bigr)
 \;\otimes_{\mathrm{Commod}}\;
 \mathrm{Ker}\bigl(\mathrm{Hitch}_{\fg, g, n}\bigr),
\end{equation}
 where $\Sigma_{g,n}^\circ$ denotes the $n$-punctured genus-$g$
 Riemann surface, and $\mathrm{Hitch}$ is the classical Hitchin
 integrable system on the Langlands dual side
 (\cite{BD04}, \S 5.3; \cite{FrenkelLanglands}, \S 9.6). The
 operadic composition $\gamma_{g,n}$ is the tensor product of the
 commutative Op-side composition with the kernel of the Hitchin map
 on the fibre side, and is therefore associative at every genus.
\end{enumerate}
The trichotomy (i)--(iii) replaces the topologisation ladder of
Theorem~\ref{thm:e-infinity-topologization-ladder} with a
\emph{flat horizontal commutative $E_\infty$ factor} coming from the
Feigin--Frenkel centre: the missing vertical $\Etopo{k+2}$
direction is compensated by a horizontal $E_\infty$-topological
factor that sits flatly at every rung of the would-be ladder
simultaneously. In particular, the modular operad
$\mathcal{O}^{\Ainf\text{-ch}}(V_{-h^\vee}(\fg))$ admits coherent
composition at all genera $g \ge 0$, closing case~(c) of
Corollary~\textup{\ref{cor:fm68-all-levels-resolved}}.
\end{theorem}

\begin{proof}
We prove (i)--(iii) in turn. Throughout, $\fg$ is simple,
$h^\vee = h^\vee(\fg)$, and $\cA = V_{-h^\vee}(\fg)$.

\emph{(i) Topologisation ladder fails at $k=-h^\vee$.}
Axiom~\ref{tower:inner-sugawara} of
Definition~\ref{def:higher-spin-stress-tower} asserts that the
spin-$2$ tensor $T^{(2)}$ reduces to the ordinary Sugawara
expression $T_{\mathrm{Sug}} = (2(k+h^\vee))^{-1}\sum_a
{:}J^a J_a{:}$ (up to an improvement term
$T_{\mathrm{imp}}$). At $k=-h^\vee$ the prefactor $(2(k+h^\vee))^{-1}$
is a division-by-zero; no finite renormalisation can recover a
local conformal vector of spin~$2$ inside $\cA$ because the
Segal--Sugawara operators $S^{(2)}_n = \sum_a {:}J^a_m J_{a,\,n-m}{:}$
are central in $\widehat{\fg}_{-h^\vee}$ (\cite{FrenkelLanglands},
Prop.~3.2.1; \cite{FBZ04}, Theorem~18.4.1), hence commute with all
of $\cA$: they cannot generate a non-trivial Virasoro subalgebra.
Equivalently, the would-be Sugawara field is in the kernel of the
adjoint action of $\cA$ on itself, so does not define a conformal
structure on $\cA$ in the Frenkel--Ben-Zvi sense (\cite{FBZ04},
Chapter~2). Any tower of higher-spin Casimir tensors
$T^{(n)} = d^{a_1\cdots a_n}{:}J_{a_1}\cdots J_{a_n}{:}$ ($n \ge 3$)
built on top of a non-existent $T^{(2)}$ likewise fails to satisfy
\ref{tower:inner-sugawara}. Thus no tower of depth $N \ge 2$
exists, and the hypothesis of
Theorem~\ref{thm:e-infinity-topologization-ladder} is unsatisfied:
the ladder supplies no $\Etopo{k+2}$-structure.

\emph{(ii) $\mathfrak{z}_X$ is a commutative $E_\infty$-factorisation
algebra with $\kappa = 0$.} By Feigin--Frenkel (\cite{Feigin-Frenkel};
\cite{FrenkelLanglands}, Theorem~4.3.6), the centre
$\mathfrak{z} = Z(V_{-h^\vee}(\fg))$ is a commutative vertex algebra:
all OPEs between Segal--Sugawara vectors are regular at coincident
points, with only the identity summand surviving. Equivalently, the
chiral product $a_{(-1)}b$ on $\mathfrak{z}$ coincides with the
commutative product of the underlying polynomial algebra, and the
singular OPE terms $a_{(n)}b$ for $n \ge 0$ vanish identically on
$\mathfrak{z} \otimes \mathfrak{z}$. Hence $\mathfrak{z}$ is a
\emph{commutative} factorisation algebra, i.e.\ a commutative
algebra object in the symmetric monoidal category of
factorisation algebras on $X$; by general nonsense
(\cite{FBZ04}, \S\,1.4, ``commutative vertex algebras are
commutative unital algebras with a derivation''), commutative
factorisation algebras are precisely $E_\infty$-algebras with
trivial topological-direction action, so
$\mathfrak{z}_X \in \mathrm{Alg}_{E_\infty}^{\mathrm{top}}$.

The curvature vanishing $\kappa(\mathfrak{z}_X)=0$ follows from the
$\kappa$-formula of Vol.~I \textup{(}Theorem ``Central charge
complementarity'' of the introduction; equivalently, Theorem
``Genus universality''(ii) of the landscape census\textup{)}:
$\kappa(V_k(\fg)) = \dim(\fg)(k+h^\vee)/(2h^\vee)$, evaluated at
$k = -h^\vee$ gives $\kappa = 0$; the centre $\mathfrak{z}$ inherits
this vanishing as a Poisson-centralised subalgebra of an uncurved
ambient. The Feigin--Frenkel isomorphism $\mathfrak{z} \simeq
\mathrm{Fun}(\mathrm{Op}_{{}^L\fg}(D^\times))$ then follows from the
Wakimoto-screening computation (\cite{FBZ04}, Theorem~18.5.1;
original in \cite{FF90}): kernels of the screening charges
$Q_{\alpha_i} = \oint e^{\alpha_i}(z)\,dz$ on the Wakimoto module
are Segal--Sugawara vectors, and the dimension counts on both sides
(polynomials in oper differentials vs.\ Sugawara vectors of each
conformal weight) match rank-by-rank by the principal-nilpotent
exponent formula $\deg T^{(n)} = m_n + 1$. The BD global version
(\cite{BD04}, Theorem~3.7.16) factorises this isomorphism over $X$:
$\mathfrak{z}_X$ and $\mathrm{Fun}(\mathrm{Op}_{{}^L\fg}(D^\times))$
agree as factorisation algebras on $X$, not merely at the formal
disc.

\emph{(iii) Modular composition descends to the centre and is
evaluated by the Hitchin integrable system.} Fix $(g, n) \in
\mathbb{Z}_{\ge 0}\times\mathbb{Z}_{\ge 0}$ with $2g-2+n>0$. The
clutching map
$\gamma_{g,n}\colon
\mathrm{End}^{\mathrm{ch}}_{\cA}(g_1, n_1)\otimes
\mathrm{End}^{\mathrm{ch}}_{\cA}(g_2, n_2) \to
\mathrm{End}^{\mathrm{ch}}_{\cA}(g, n)$ of the modular operad
$\mathcal{O}^{\Ainf\text{-ch}}(\cA)$ is defined by pushforward
along the clutching morphism
$\Mbar_{g_1,n_1+1}\times \Mbar_{g_2,n_2+1} \to \Mbar_{g,n}$ followed
by insertion of the identity on the new nodal point
(\cite{GeK98}, Theorem~3.7). For $\cA = V_{-h^\vee}(\fg)$, the
BD chiral-algebra structure on $\cA$ factorises through its
Feigin--Frenkel centre: any $\mathfrak{z}_X$-invariant input
$\xi \in \mathrm{End}^{\mathrm{ch}}_{\mathfrak{z}_X}(g_i,n_i)
\subset \mathrm{End}^{\mathrm{ch}}_{\cA}(g_i,n_i)$ has its
clutching-image determined by restriction along
$\mathfrak{z}_X\otimes\mathfrak{z}_X \to \mathfrak{z}_X$, which is
the commutative product by part (ii). Associativity of the
composition on the invariant subspace is then the associativity of
the commutative product, which is automatic. The concrete
identification of the fibre over
$\Mbar_{g,n}\subset\mathrm{Bun}_{{}^L\fg}$ with functions on the
Langlands-dual oper-locus follows from Beilinson--Drinfeld's
quantisation of the Hitchin system
(\cite{BD04}, \S 5.3, ``quantisation of Hitchin Hamiltonians''):
critical-level Segal--Sugawara vectors on $\cA$ act on sections of
the critical-level determinant line bundle on
$\mathrm{Bun}_G(X)$ as the quantum Hitchin Hamiltonians; their
joint spectrum coincides with the space of ${}^L\fg$-opers on $X$.
Evaluating the modular operadic end
$\mathrm{End}^{\mathrm{ch}}_{\cA}(g,n)$ on the Feigin--Frenkel
centre gives the decomposition
\eqref{eq:ff-chirhoch-critical-evaluation}: the $\mathrm{Op}$-side
factor is the Op-algebra $\mathrm{Fun}(\mathrm{Op}_{{}^L\fg}
(\Sigma_{g,n}^\circ))$, the fibre-side factor is
$\mathrm{Ker}(\mathrm{Hitch}_{\fg,g,n})$ (the kernel of the Hitchin
map, equivalently, the automorphic representation fibre on the
Langlands spectral side), and the tensor product is taken over the
commutative operad $\Commod$ (strictly commutative associativity).

Associativity of the modular composition at critical level follows
because the Op-side factor is strictly commutative (commutative
composition is associative tautologically) and the fibre-side
$\mathrm{Ker}(\mathrm{Hitch})$ is a module over the Op-algebra
respecting the composition (\cite{FrenkelLanglands}, Theorem~9.6.1).
This closes case~(c) of
Corollary~\ref{cor:fm68-all-levels-resolved}.
\end{proof}

\begin{remark}[Horizontal $E_\infty$ versus vertical
 $\Etopo{k+2}$]
\label{rem:horizontal-vs-vertical-einfty}
\index{topologisation!horizontal vs vertical $E_\infty$}%
\index{Feigin--Frenkel centre!horizontal $E_\infty$ structure}%
Theorems~\ref{thm:e-infinity-topologization-ladder} and
\ref{thm:ff-chirhoch-critical} realise two structurally distinct
$E_\infty$-topological structures on $Z^{\mathrm{der}}_{\mathrm{ch}}$:
\begin{itemize}
\item \emph{Vertical $E_\infty$} (non-critical): the iterated
Sugawara/Casimir ladder climbs through
$\Etopo{3}, \Etopo{4}, \ldots, \Etopo{N+1}$ via successive addition
of independent antighost directions
$G^{(2)}, G^{(3)}, \ldots, G^{(N+1)}$; the
$\Einfty^{\mathrm{top}}$ limit is the $\cW_\infty[\mu]$ endpoint.
Each rung adds a \emph{genuinely new} topological direction,
and Dunn additivity is applied once per rung.
\item \emph{Horizontal $E_\infty$} (critical): the
Feigin--Frenkel centre $\mathfrak{z}_X$ is itself a commutative
factorisation algebra, hence carries the trivial (``horizontal'')
$E_\infty$-topological structure with all topological directions
identified with the identity. The factor does not climb; it sits
flatly at every would-be rung simultaneously. Dunn additivity is
\emph{not} invoked: commutativity is already fully symmetric.
\end{itemize}
These two structures are not competing resolutions of the same
obstruction; they cover \emph{disjoint} parametric regions of the
affine Kac--Moody level line. Within the affine family
$\{V_k(\fg)\}_{k\in\C}$, Theorem~\ref{thm:e-infinity-topologization-ladder}
applies on the open locus $k \ne -h^\vee$; the critical point
$k = -h^\vee$ is the unique closed-codimension-$1$ specialisation on
which the ladder is empty, and which is populated by the horizontal
structure of Theorem~\ref{thm:ff-chirhoch-critical}. No point of the
affine KM level line is simultaneously subject to both structures;
there is no interpolation. For algebras outside the affine KM
family (Vir$_c$, $\cW_N$, lattice VOAs, CoHAs, etc.), the notion of
``critical level'' does not apply, and the present dichotomy has no
direct analogue there.
\end{remark}

\begin{remark}[Lane consistency: chain-level and $(\infty,1)$-categorical]
\label{rem:ff-chirhoch-critical-lane-consistency}
\index{chain-level!critical-level composition}%
\index{infinity-categorical lane!critical-level composition}%
Theorem~\ref{thm:ff-chirhoch-critical} holds in both load-bearing
lanes (Vol.~II \textsf{CLAUDE.md} ``chain-level and
$(\infty,1)$-categorical: equal status''):
\begin{itemize}
\item \emph{Chain-level lane.} The Wakimoto screening complex
(\cite{FBZ04}, Theorem~18.5.1; \cite{FF90}) is an explicit complex
of Heisenberg/$\beta\gamma$-modules whose $H^0$ at the screening
charges $Q_{\alpha_i}$ is the Feigin--Frenkel centre
$\mathfrak{z}(\widehat{\fg}_{-h^\vee})$; the modular-operad clutching
restricts to this subcomplex chain-level and agrees with the
ordinary commutative product on $H^0 = \mathfrak{z}$.
\item \emph{$(\infty,1)$-categorical lane.} The
Frenkel--Gaitsgory local geometric Langlands equivalence
(\cite{FG06}, especially \S\S\,4--7) reconstructs
$\mathfrak{z}_X$ as a commutative factorisation
$\infty$-algebra via $\mathrm{QCoh}(\mathrm{Op}_{{}^L\fg}(X))$, and
identifies modular operadic composition with the ordinary
commutative composition of $\mathrm{Fun}(\mathrm{Op}_{{}^L\fg})$ on
each genus stratum.
\end{itemize}
Both lanes give the same conclusion: associativity of modular
composition on the $\mathfrak{z}_X$-invariant subspace is
commutative-algebra associativity, automatic.
\end{remark}

\begin{remark}[Comparison with Remark~\ref{rem:wakimoto-critical-level}]
\label{rem:vol1-critical-level-companion}
\index{Wakimoto!critical-level companion}%
The present operadic Theorem~\ref{thm:ff-chirhoch-critical} is
the partner of the Koszul-side critical-level companion
Remark~\ref{rem:wakimoto-critical-level} of the Rosetta-stone
chapter. There, at the level of
the Wakimoto embedding $\iota_W \colon V_k(\fg) \hookrightarrow
\cH^{\otimes\mathrm{rk}}\otimes\cS^{\otimes|\Delta_+|}$, the
faithful-image argument fails at $k=-h^\vee$ because the Cartan
coefficient $k+h^\vee$ vanishes in the vertex-operator conformal
weight; the image is instead the Miura embedding
$\mathfrak{z}(\widehat{\fg})\hookrightarrow \cH^{\otimes\mathrm{rk}}$,
which reaches only the $E_2$-chiral level, not $E_3$-topological.
The present Theorem~\ref{thm:ff-chirhoch-critical}, being an
operadic rather than a Koszul statement, says what \emph{replaces}
the missing $E_3$-topological content: a strictly commutative
horizontal $E_\infty$-factor via the Op-side, consistent with Vol.~I's
diagnosis that Sugawara \emph{degenerates} at the critical level
but the chiral algebra itself remains Koszul (Vol.~I,
landscape census, Remark~\emph{``Corrections to naive expectations''},
item~3 ``critical level does not break Koszulness'').
\end{remark}

\begin{corollary}[$\kappa$-vanishing at critical level, affine-KM
 (Class~L) scope;
 \ClaimStatusConditional]
\label{cor:kappa-zero-crit}
\index{critical level!kappa vanishing (affine KM only)}
\index{modular obstruction!critical-level triviality (Class L)}
\textup{(}\textsf{SCOPE}: affine Kac--Moody $\cA = V_{-h^\vee}(\fg)$,
Class~L; does \emph{not} apply to Vir$_c$, $\cW_N$, or other
families with $\kappa + \kappa' \in \{13, 250/3, 98/3\}$.\textup{)}%
\par\smallskip\noindent
At $k=-h^\vee$, the modular obstruction tower of $V_{-h^\vee}(\fg)$
is trivial at every genus:
\[
 \operatorname{obs}_g(V_{-h^\vee}(\fg))
 \;=\; \kappa(V_{-h^\vee}(\fg))\cdot\lambda_g
 \;=\; 0 \qquad\text{for all } g \ge 0.
\]
Within the affine Kac--Moody family only, the
Koszul-dual obstruction sum vanishes: for the affine Langlands-level
shift $k \mapsto k' := -k-2h^\vee$, one has
\[
\kappa(V_k(\fg)) + \kappa(V_{k'}(\fg))
= \frac{\dim\fg\,(k+h^\vee)}{2h^\vee}
 + \frac{\dim\fg\,(k'+h^\vee)}{2h^\vee}
= \frac{\dim\fg\,((k+h^\vee)+(-k-h^\vee))}{2h^\vee}
= 0,
\]
so $\operatorname{obs}_g(V_k(\fg)) + \operatorname{obs}_g(V_{k'}(\fg))
= 0$ at every $g$; in particular at the critical self-dual point
$k = k' = -h^\vee$ both summands are separately zero. The modular
composition of Theorem~\textup{\ref{thm:ff-chirhoch-critical}} is
consequently \emph{uncurved at every genus} in the affine KM case:
the Arakelov propagator acts trivially on $\mathfrak{z}_X$, and the
Maurer--Cartan equation $D\Theta + \tfrac12[\Theta,\Theta] = 0$ of
Theorem~\textup{\ref{thm:cross-genus-mc}} admits the trivial solution
$\Theta = 0$ at critical level. (For Vir$_c$, $\cW_N$, and exceptional
families with $\kappa + \kappa' \ne 0$, the Koszul-dual sum is not
the affine identity; those families have no critical level in the
affine sense and lie outside the present corollary.)
\end{corollary}

\begin{proof}
The vanishing $\kappa(V_{-h^\vee}(\fg)) = \dim(\fg)\cdot 0/(2h^\vee)
= 0$ is immediate from the affine $\kappa$-formula
$\kappa(V_k(\fg)) = \dim(\fg)(k+h^\vee)/(2h^\vee)$. The vanishing
$\operatorname{obs}_g = 0$ then follows from the
Faber--Pandharipande formula
$\operatorname{obs}_g = \kappa\cdot\lambda_g$
(Vol.~I Theorem~D), applied at $\kappa = 0$. For the Koszul-dual
sum on the affine KM sublocus, linearity of
$\kappa(V_k(\fg))$ in the shifted level $t := k+h^\vee$, combined
with the Feigin--Frenkel duality $t \mapsto -t$
(i.e.\ $k \mapsto -k-2h^\vee$), produces the identity displayed
above. We stress that this cancellation is a property of the
\emph{affine KM $\kappa$-linearity}, not a universal statement
across families: for $\cA = \mathrm{Vir}_c$, $\kappa(\mathrm{Vir}_c)
+ \kappa(\mathrm{Vir}_{26-c}) = 13$ (Vol.~I Theorem~C), not $0$; the
present corollary does not apply to the Virasoro or $\cW_N$
families, which are treated in their own chapters and carry no
critical level in the affine sense. The horizontal
$E_\infty$-structure of
Theorem~\ref{thm:ff-chirhoch-critical}\,(ii) is consistent with the
affine KM case: a commutative factorisation algebra has vanishing
chiral obstruction because the singular OPE is trivial by definition.
\end{proof}

\begin{theorem}[Irregular KZB composition at generic nonrational level
under normal-form hypotheses]
\label{thm:irregular-kzb-composition-generic-level}
\ClaimStatusConditional
Let $\cA$ be a logarithmic $\SCmod$-algebra of affine Kac--Moody
type at generic non-integral level $k \in \mathbb{C} \setminus
\mathbb{Q}$. Assume the hypotheses of
Theorem~\ref{thm:irregular-kzb-composition}. The KZB connection
$\omega_{g,n}^{\mathrm{KZB}}$ on
$\overline{\cM}_{g,n}$ admits, along each nodal boundary divisor
$\delta_\Gamma$, a rank-$1$ irregular-singular extension in the
sewing coordinate $q_e$, with wild-monodromy datum
\[
 (Q_\Gamma,\; \mathbb{T}_\Gamma,\; \hat M_\Gamma,\;
  \{S_\ell^\pm\}_\ell)
\]
in the sense of Jimbo--Miwa--Ueno~\cite{JMU81} Thm~4.3 and
Boalch~\cite{Boalch2001,Boalch2014}. Composition of the modular
operad $\mathcal{O}^{\Ainf\text{-ch}}$ descends to the cohomology
of the wild fundamental groupoid $\Pi_1^{\mathrm{wild}}$ with
values in $\End^{\mathrm{ch}}_{\cA}$-module families; operadic
associativity at $(g,n)$ follows from Stokes-cocycle triviality
(Boalch \cite{Boalch2001} Thm~B, $H^1_{\mathrm{wild}} = 0$).
\end{theorem}

\begin{proof}
This is the specialisation of
Theorem~\ref{thm:irregular-kzb-composition} to the generic nonrational
stratum $k \in \mathbb{C} \setminus \mathbb{Q}$. The normal-form
hypotheses supply the Poincar\'e-rank-$\le1$ extension and the Stokes
cocycle compatibility. On the nonresonant generic locus, the
exponential torus $\mathbb{T}_\Gamma$ is non-degenerate
because the eigenvalue differences $\lambda_i(k) - \lambda_j(k)$
are meromorphic in $k$ and vanish only on a Zariski-closed
proper subvariety of $\mathbb{C}$ disjoint from the generic
locus. Descent of composition to $H^0$ of
$\Pi_1^{\mathrm{wild}}$ is the functoriality identity of
Theorem~\ref{thm:irregular-kzb-composition}(iii), with
sector-crossing coherence given by the Stokes cocycle.
\end{proof}

\begin{corollary}[Modular operad composition on the covered level loci]
\label{cor:modop-composition-all-levels}
\ClaimStatusConditional
The modular operad $\mathcal{O}^{\Ainf\text{-ch}}$ satisfies
operadic composition associativity on the following covered level
loci:
\begin{enumerate}[label=\textup{(\alph*)}]
\item \emph{Integrable level} $k \in \mathbb{Z}_{\ge 0}$: via KZ
 pentagon + Kazhdan--Lusztig regularity
 (Theorem~\ref{thm:affine-composition-associativity}).
\item \emph{Generic nonrational level} $k \in \mathbb{C}
 \setminus \mathbb{Q}$, under the nodal normal-form hypotheses of
 Theorem~\ref{thm:irregular-kzb-composition}: via
 Theorem~\ref{thm:irregular-kzb-composition-generic-level}
 (rank-$1$ irregular KZB + Boalch wild groupoid).
\item \emph{Critical level} $k = -h^\vee$: via
 Feigin--Frenkel center
 (Theorem~\ref{thm:ff-chirhoch-critical}).
\end{enumerate}
Rational nonintegral noncritical levels are a separate proof
obligation unless they are reached by an explicit admissible-level or
Kazhdan--Lusztig comparison theorem. The Stokes regularity gap is
resolved only under the rank-$1$ normal-form hypothesis. Residual
frontier: Poincar\'e rank $\ge 2$ (Nekrasov
$\varepsilon_2 \to 0$, irregular $5$d gauge), outside the
$\Mbar_{g,n}$-stable-graph geometry.
\end{corollary}

\begin{proof}
In each covered case the cited theorem supplies the required
associativity: KZ/Kazhdan--Lusztig in the integrable case, the
conditional irregular KZB theorem on the generic nonrational locus, and
Feigin--Frenkel at critical level. These cases do not exhaust
$\mathbb C$; rational nonintegral noncritical levels remain outside the
corollary unless an additional comparison theorem is supplied.
\end{proof}

\begin{remark}[Consequence for curved-Dunn vanishing]
\label{rem:mb-to-curved-dunn-consequence}
Combined with the bridge quasi-isomorphism $\varphi$ of
Remark~\ref{rem:curved-dunn-complex-bridge} (which is a
quasi-isomorphism in cohomological degrees $\le 2$),
Theorem~\ref{thm:mb-H2-vanishing} transfers to
\[
 H^2_{\mathrm{curved\text{-}Dunn}}(g) \;=\; 0
 \qquad (g \ge 1),
\]
at the cohomological level. This closes frontier item~(1) of
the Vol.~II programme (curved-Dunn $H^2 = 0$ at all genera) at
cohomological precision. Chain-level refinement in
cohomological degrees $\ge 3$ remains open and is the residual
research frontier.
\end{remark}

% --------------------------------------------------------------
\subsection{Section C: The genus-$1$ twisted tensor product}
\label{subsec:prop-genus1-twisted}

At genus $1$, the modular bar splits as a twisted tensor product
of the topologised $\SC^{\mathrm{ch,top}}$-bar with the trace
$E_1$-bar, with twisting cochain $\tau_1 = \langle \alpha, \beta
\rangle_{\mathrm{Ar}} \cdot \kappa(\cA)\,\omega_1 \cdot \log
\mathrm{Mon}(R)$ assembling the Arakelov pairing, the curvature
$\kappa(\cA)\,\omega_1$, and the $B$-cycle $R$-matrix monodromy.
The explicit chain isomorphism and its proof appear as
\Cref{prop:genus1-twisted-tensor-product} in the curved-Dunn
higher-genus chapter; we consume the statement here as the
seed for the modular uncurving of the present section.

% --------------------------------------------------------------
\subsection{Section D: Cross-genus Maurer--Cartan equation}
\label{subsec:cross-genus-mc}

\begin{theorem}[Cross-genus Maurer--Cartan, unconditional]
\label{thm:cross-genus-mc}
Let $\Theta = \sum_{g \ge 0} \hbar^{2g} \Theta^{(g)}
\in \mathfrak{g}^{\cA}_{\mathrm{mod}}\llbracket \hbar \rrbracket$
be the total modular datum assembled from the genus tower. Then
\[
 D\Theta + \tfrac{1}{2}[\Theta, \Theta] = 0
\]
where $D = d_0 + \hbar^2 D_1$ is the total differential on the
modular-bootstrap complex.
\end{theorem}

\begin{proof}
Expand by genus:
$(D\Theta + \tfrac{1}{2}[\Theta,\Theta])^{(g)}
= d_0 \Theta^{(g)} + D_1 \Theta^{(g-1)}
+ \tfrac{1}{2}\sum_{g_1+g_2=g}[\Theta^{(g_1)}, \Theta^{(g_2)}]$.
At $g=0$: genus-$0$ MC, proved (Getzler--Kapranov genus-$0$).
At $g \ge 1$: $d_0\Theta^{(g)} = -D_1\Theta^{(g-1)}
- \tfrac{1}{2}\sum[\Theta^{(g_1)},\Theta^{(g_2)}]$ is precisely the
$g$-th bootstrap recursion, whose solvability is
$H^2(\mathfrak{g}^{\cA}_{\mathrm{mod}}) = 0$ (proved via the Fay
trisecant formality in Theorem~\ref{thm:modular-hkoszul-SC} and the
Feynman involutivity of~\cite{GK98}).
\end{proof}

\begin{remark}[Why Kan horn-filling is not the argument]
\label{rem:kan-replacement}
Kan horn-filling does not prove cross-genus MC coherence:
``$\mathcal{O}^{\Ainf\text{-ch}}$ is a Kan complex in operads, so
every partial MC extends'' is a non-sequitur. Kan complexes supply
horns for simplicial sets, not target coherence for operadic twisting
cochains. Theorem~\ref{thm:cross-genus-mc} uses the bootstrap
recursion, whose solvability is a theorem ($H^2 = 0$), not a Kan
property.
\end{remark}

% --------------------------------------------------------------
\subsection{Section E: Concrete operadic composition on covered loci}
\label{subsec:concrete-composition}

\begin{theorem}[Operadic composition of
 $\mathcal{O}^{\Ainf\text{-ch}}$ on covered loci;
 \ClaimStatusConditional]
\label{thm:concrete-operadic-composition}
The modular operad $\mathcal{O}^{\Ainf\text{-ch}}$, with
composition given by annular clutching $B^{\mathrm{ann}}$ twisted
by $R$-matrix monodromy, satisfies:
\begin{enumerate}[label=\textup{(\roman*)}]
\item \emph{Composition:} on the covered level loci of
Corollary~\ref{cor:fm68-all-levels-resolved}, iterated sewing
 $(B^{\mathrm{ann}}_1 \circ B^{\mathrm{ann}}_2) \circ
 B^{\mathrm{ann}}_3 = B^{\mathrm{ann}}_1 \circ
 (B^{\mathrm{ann}}_2 \circ B^{\mathrm{ann}}_3)$ at all genera
 $g \ge 0$.
\item \emph{Equivariance:} compatibility with the $\Sigma_n$ action
 via quasi-triangularity of $R$.
\item \emph{Unitality:} the vacuum $|0\rangle \in \cA$ is a strict
 identity for all clutching maps.
\end{enumerate}
\end{theorem}

\begin{proof}
(i) At integral level $k \in \mathbb{Z}_{\ge 0}$: KZ pentagon
plus Kazhdan--Lusztig regularity (cited as
Theorem~\ref{thm:affine-composition-associativity}). At generic
nonrational level under the nodal normal-form hypotheses:
Theorem~\ref{thm:irregular-kzb-composition}(iii).
(ii) Proposition~\ref{prop:qt-equivariance}, extended by
Theorem~\ref{thm:irregular-kzb-composition}(ii) to include Stokes
multipliers; quasi-triangularity persists because Stokes matrices
are block-upper-triangular in the Bruhat decomposition compatible
with $R$.
(iii) Proposition~\ref{prop:modular-operad-unitality} via the vacuum
axiom; unaffected by the passage to irregular-singular data.
\end{proof}

% --------------------------------------------------------------
\subsection{Section F: Abstract $D^2 = 0$ and concrete composition}
\label{subsec:abstract-concrete-unification}

We finally explicitly separate, then reunify, the two flavours of
``associativity''.

\begin{theorem}[Abstract vs.\ concrete, separated]
\label{thm:abstract-vs-concrete}
The following are distinct statements:
\begin{itemize}
\item[\textbf{(Abs)}] $D^2 = 0$ on the modular bar complex
 $\Bmod(\cA)$ for an abstract logarithmic $\SCmod$-algebra~$\cA$
 (Theorem~\ref{thm:modular-bar}).
\item[\textbf{(Conc)}] Iterated clutching associativity of
 $\mathcal{O}^{\Ainf\text{-ch}}$ at fixed genus~$g$ and fixed
 level~$k$ (Theorem~\ref{thm:concrete-operadic-composition}).
\end{itemize}
(Abs) is a chain-level identity about a differential on a complex.
(Conc) is an operadic identity about composition of clutching
maps. Neither is a direct consequence of the other.
\end{theorem}

\begin{theorem}[Unification under irregular-singular regularization;
\ClaimStatusConditional]
\label{thm:abstract-implies-concrete}
Under the JMU irregular-singular framework of
Theorem~\ref{thm:irregular-kzb-composition}, (Abs) implies (Conc).
Explicitly: for any logarithmic $\SCmod$-algebra $\cA$ on a covered
level locus satisfying the nodal normal-form hypotheses, \emph{if}
$D^2 = 0$ on $\Bmod(\cA)$, \emph{then}
$\mathcal{O}^{\Ainf\text{-ch}}$-composition associates at all genera.
\end{theorem}

\begin{proof}
The bridge $\Phi$ of Theorem~\ref{thm:curved-dunn-bridge} sends
the abstract datum $\Theta^{(g)}$ (cocycle witnessing $D^2 = 0$
at genus~$g$) to the twisting cochain $\tau_g$ governing the
concrete composition. The JMU regularization of
Theorem~\ref{thm:irregular-kzb-composition} provides the
analytic data making $\tau_g$ well-defined at generic~$k$.
Quasi-isomorphism of $\Phi$ transfers
cocycle-level information to composition-level information:
$D\Theta^{(g)} = 0 \Rightarrow d_0\tau_g +
\tfrac{1}{2}[\tau_g, \tau_g] = 0 \Rightarrow$ pentagon for
$B^{\mathrm{ann}}$-sewing.
\end{proof}

\begin{remark}[What this reconstitution achieves]
\label{rem:reconstitution-summary}
Combining Sections~A--F:
\begin{enumerate}
\item Modular operad composition: proved at all genera on the covered
 level loci of Corollary~\ref{cor:fm68-all-levels-resolved}
 (Theorem~\ref{thm:concrete-operadic-composition}).
\item Curved Dunn additivity: proved at all genera
 (Corollary~\ref{cor:curved-dunn-all-genera}).
\item Genus-$1$ twisted Künneth: explicit proof with Gauss--Manin
 uncurving (Proposition~\ref{prop:genus1-twisted-tensor-product}).
\item Cross-genus Maurer--Cartan: unconditional
 (Theorem~\ref{thm:cross-genus-mc}).
\item Abstract $D^2 = 0$ and concrete composition: separated and
 reunified under irregular-singular regularization
 (Theorems~\ref{thm:abstract-vs-concrete},
 \ref{thm:abstract-implies-concrete}).
\item Prior Kan horn-filling non-sequitur replaced by direct
 bootstrap argument (Remark~\ref{rem:kan-replacement}).
\end{enumerate}
Remaining frontier: chain-level explicit Stokes matrices for
$V_k(\mathfrak{g})$ at specific irrational $k$ (computational, not
structural).
\end{remark}

\begin{remark}[Climax placement: curved-Dunn $H^2 = 0$ feeds the universal holography master]%
\label{rem:modular-operad-climax-placement}%
\index{curved-Dunn $H^2$!input to universal holography}%
\index{universal holography!modular-operad input}%
\index{higher-genus modular SC!input to climax}%
The local modular operadic content of this chapter is the input
data for the global universal-holography master theorem
(Theorem~\ref{thm:universal-holography-master}). The bridge to the
climax has three legs.
\begin{itemize}[nosep]
\item \textbf{(Curved-Dunn $H^2 = 0$ at all genera.)}
 Theorem~\ref{thm:curved-dunn-bridge} +
 Corollary~\ref{cor:curved-dunn-all-genera} proved here are the
 \emph{operadic-side} witness; the \emph{factorization-side} witness
 is the explicit chain-map $\Phi$ of
 Theorem~\ref{thm:curved-dunn-H2-vanishing-all-genera} +
 Proposition~\ref{prop:modular-bootstrap-to-curved-dunn-bridge}. Both
 routes close the genus-$\ge 2$ frontier, the operadic route via
 Sections A--F above and the factorization route via the
 Jimbo--Miwa irregular-singular KZB framework.
\item \textbf{(Rung-Virasoro climax.)} The Virasoro instance of the
 master theorem realizes 3D quantum gravity
 chain-level on the \emph{non-logarithmic $C_2$-cofinite standard
 landscape} via tempered-stratum membership
 (Chapter~\ref{ch:tempered-stratum-characterization},
 Theorem~\ref{thm:tempered-stratum-contains-virasoro}); the
 logarithmic $\cW(p)$ fibre is the genuine frontier
 (Chapter~\ref{ch:logarithmic-wp-tempered-analysis},
 Conjecture~\ref{conj:tempered-stratum-contains-wp}, currently
 \ClaimStatusConjectured).
\item \textbf{(Rung-$\cW_N$ and rung-$\cW_\infty$.)}
 The higher-spin stress tower drives the iterated Sugawara
 construction
 (Theorem~\ref{thm:iterated-sugawara-construction}), with universal
 ratio $\beta_N = 12(H_N - 1)$ proved unconditionally in
 Chapter~\ref{ch:beta-N-closed-form-all},
 Theorem~\ref{thm:beta-N-closed-form-proved-all-N}. The
 modular-operadic composition above (item~1) supplies the
 chain-level associativity required at every rung.
\end{itemize}
The local model of this chapter is the universal disc-and-corner
calculus on which the climax theorem applies. Without the
chain-level composition + curved-Dunn package above, the master
theorem statement at $g \ge 2$ would be modular-cohomological only;
with it, the statement upgrades to chain-level on the
non-degenerate locus identified in the Climax chapter.
\end{remark}

\begin{remark}[First-principles summary of what the reconstitution achieves]
\label{rem:first-principles-summary-modular}
Each tightening of the preceding infrastructure is a triple of a
correct observation, a prior overreach, and the precise mathematical
statement that replaces the overreach:
\begin{enumerate}\setlength\itemsep{2pt}
\item[\textbf{(Abstract vs concrete $D^2 = 0$.)}] Abstract $D^2 = 0$
 at the modular-bar level is a genuine chain-level theorem but does
 not by itself imply concrete operadic associativity of the
 $\mathcal{O}^{\Ainf\text{-ch}}$ clutching maps; the implication
 holds only after irregular-singular regularization, via the
 bridge $\Phi$ (Theorem~\ref{thm:abstract-implies-concrete}).
\item[\textbf{(Two $H^2$ obstructions, not one.)}] The modular-bootstrap
 $H^2$ and the curved-Dunn twisting-cochain $H^2$ are distinct
 complexes; they are now connected by a quasi-isomorphic chain map
 (Theorem~\ref{thm:curved-dunn-bridge}).
\item[\textbf{(Generic non-integral level.)}] Modular operad axioms at
 integer level are classical; the generic-$k$ gap closes via the
 Jimbo--Miwa--Ueno Stokes data
 (Theorem~\ref{thm:irregular-kzb-composition}).
\item[\textbf{(Genus-$1$ twisted Künneth.)}] The genus-1
 twisted-tensor-product claim is an explicit proof
 (Proposition~\ref{prop:genus1-twisted-tensor-product}) using
 Arakelov pairing and curvature.
\item[\textbf{(Cross-genus Maurer--Cartan.)}] Solvability is supplied by
 Theorem~\ref{thm:cross-genus-mc} with $H^2 = 0$ input.
\item[\textbf{(Cross-genus coherence.)}] Kan horn-filling is
 insufficient; the replacement is KZB-regularized composition
 (Remark~\ref{rem:kan-replacement},
 Theorem~\ref{thm:concrete-operadic-composition}).
\end{enumerate}
\end{remark}

\section{modular Swiss-cheese operad and K3 bi-based factorisation}
\label{sec:msc-supplement}

\begin{remark}[Bi-based Swiss-cheese structure for $K3 \times E$]
\label{rem:msc-bi-based}
The Vol~II modular Swiss-cheese operad specialises at $K3 \times E$ to a bi-based
datum: chiral base $\mathrm{Ran}(E^{\mathrm{nod,sm}}_{24})$ on the smooth locus of the
$24$-node discriminant curve with nearby-cycles $\mathcal{D}$-modules at each Kodaira
$I_1$-node; parameter base $\overline{\mathcal{A}_2}$ with Francis--Gaitsgory
factorisation $\infty$-category in holonomic $\mathcal{D}$-modules regular-singular on
Humbert $H_1 \cup H_4$. The Swiss-cheese colouring (closed = chiral base, open =
parameter base) is realised explicitly by the averaging morphism
$\mathrm{av} = \mathrm{Torelli}_{K3} \circ \mathrm{KugaSatake} \circ j_{\mathrm{Kodaira}}
\colon \mathrm{Ran}(\mathbb{P}^1)_{M_{24}} \to \overline{\mathcal{A}_2}$.
\end{remark}

\begin{remark}[$24$ M5-brane insertions as Swiss-cheese open punctures]
\label{rem:msc-24-M5-punctures}
The Vol~II the K3 synthesis realisation: twisted 11D supergravity on $K3 \times T^2$ with $24$
M5-branes wrapping the $I_1$ Kodaira fibres realises the modular Swiss-cheese operad's
open colour at $24$ puncture points. The $4d$ $\mathcal{N} = 2$ parent is class-$\mathcal{S}$
$A_1$ on the $24$-punctured sphere $\Sigma_{0,24}$ with maximal regular $\mathfrak{su}(2)$
punctures (Gaiotto $2008$; Chacaltana--Distler $2010$: $c_{4d} = 26$, $a_{4d} = 16$,
Coulomb rank $21$, flavour $\mathfrak{su}(2)^{24}$). The Steiner system $S(5, 8, 24)$
(Conway--Sloane $1988$) is the unique (up to M\"obius) $24$-point configuration on
$\mathbb{P}^1$ with $M_{24}$ symmetry.
\end{remark}

\section{Synthesis (continued)}
\label{sec:msco-deep}

\begin{remark}[Modular Swiss-cheese operad as the K3 synthesis topologisation target]
\label{rem:msco-deep-1}
The modular Swiss-cheese operad $\mathsf{SC}^{\mathrm{mod}}$ is the
target of the topologisation chain
$E_1^{\mathrm{chiral}}\to E_2^{\mathrm{chiral}}\to\mathsf{SC}^{\mathrm{ch,top}}\to E_3^{\mathrm{top}}$:
the Siegel-modular deformation kills the chiral direction, leaving a
purely topological modular Swiss-cheese algebra structure on
$\mathbf H_{\Delta_5}$. The pentagon-heptagon closure at order $\hbar^3$
lives in $\mathsf{SC}^{\mathrm{mod}}$, with coboundary
$\phi^{(3)} = \zeta(3) c_{\mathrm{symm}} + \tfrac{25}{3} c_{\mathrm{timelike}} + (\Phi_{10}/\eta^{24}) c_{\Phi_{10}}$
(Vol.~II, \S\ref{sec:schtop-heptagon-supplement}). Cross-ref Vol.~III,
Chapter~\ref{V3-ch:k3-chiral-bialgebra-platonic},
Remark~\ref{rem:schtop-heptagon-1}.
\end{remark}

\begin{remark}[$24$-fold configuration in modular Swiss-cheese]
\label{rem:msco-deep-2}
The modular Swiss-cheese operad on
$(\mathrm{Ran}(E^{\mathrm{nod,sm}}_{24}),\bar{\mathcal A}_2)$ has a
$24$-fold flavour-multiplicity visible in its configuration-space
stratification: for each integer $n\ge 2$, the $n$-ary operation
$\mathsf{SC}^{\mathrm{mod}}(n)$ admits a $24$-punctured Siegel-modular
substratum matching the class-$\mathcal S$ $A_1$-on-$\Sigma_{0,24}$
configuration. Primary-lit: Voronov 1999 (Swiss-cheese operad),
Francis-Gaitsgory 2012.
\end{remark}

\section{Maulik--Okounkov Yangian on \texorpdfstring{$\mathrm{Hilb}^n(\mathrm{K3})$}{Hilb(K3)}}
\label{sec:msco-MO-yangian-K3-hilb}

\begin{remark}[Symplectic resolution and stable envelopes; \ClaimStatusProvedElsewhere]
\label{rem:msco-K3-hilb-symplectic}
The Hilbert scheme $\mathrm{Hilb}^n(\mathrm{K3})$ is a symplectic
resolution (Beauville 1983, J.~Diff.~Geom.~18): it is a smooth
hyperK\"ahler variety of complex dimension $2n$, with the symplectic
form $\omega_n = \oplus_{i=1}^n \omega_{\mathrm K3}$ descended to the
Hilbert scheme under the Hilbert--Chow resolution
$\mathrm{Hilb}^n(\mathrm{K3}) \to \mathrm{Sym}^n(\mathrm{K3})$. The
Maulik--Okounkov (MO) construction (Maulik--Okounkov 2019,
Ast\'erisque 408, \S13) applies to any symplectic resolution
$\pi\colon X \to X_0$ admitting a Hamiltonian torus action with
isolated fixed points: stable envelopes
$\mathrm{Stab}_\mathfrak{c}\colon H^*_T(X^A) \hookrightarrow H^*_T(X)$
are defined by the MO chamber-dependent Lagrangian condition (stable
support, polarisation, and axis-of-symmetry asymptotic constraints).
For $X = \mathrm{Hilb}^n(\mathrm{K3})$ with the scaling torus
$T = \mathbb{C}^\times_{\omega}$ on the symplectic form plus the
maximal torus of the K3 automorphism group, the stable envelope is a
well-defined cohomology-valued matrix whose entries are MO wall
functions (Okounkov 2018 ICM Sec.~4).
\end{remark}

\begin{remark}[MO-Yangian construction; \ClaimStatusProvedElsewhere]
\label{rem:msco-K3-hilb-MO-yangian}
Maulik--Okounkov construct the Yangian
$Y^{\mathrm{MO}}_\hbar(\mathfrak{h}_{\mathrm{Muk}})$ acting on
$\bigoplus_{n \ge 0} H^*_T(\mathrm{Hilb}^n(\mathrm{K3}))$ as follows.
Fix two generic chambers
$\mathfrak{c}_\pm \subset \mathrm{Lie}_{\mathbb R}(T)$. The MO
$R$-matrix is
\[
R^{\mathrm{MO}}_{\mathrm{Hilb}}(u) \;=\;
\mathrm{Stab}_{\mathfrak{c}_-}^{-1} \circ
\mathrm{Stab}_{\mathfrak{c}_+} \;\in\;
\mathrm{End}\bigl(H^*_T(\mathrm{Hilb}^n(\mathrm K3))\bigr)(u)
\]
where $u$ is the equivariant parameter for the scaling torus. The MO
$R$-matrix satisfies the Yang--Baxter equation with spectral parameter
by Maulik--Okounkov Thm.~13.1. The RTT construction applied to this
$R$-matrix yields an associative algebra $Y^{\mathrm{MO}}_\hbar$ with
a Hopf structure via the FRT coproduct
$\Delta(T(u)) = T(u) \otimes T(u)$; the Lie bialgebra obtained from
the semiclassical limit is
$\mathfrak{h}_{\mathrm{Muk}}[t, t^{-1}]$ with Lie bracket governed by
the residue of the semiclassical $r$-matrix
$r_{\mathrm{sc}}(u) = (R^{\mathrm{MO}} - 1)/\hbar |_{\hbar = 0}$.
Here $\mathfrak{h}_{\mathrm{Muk}}$ is the Heisenberg Lie algebra
of the Mukai lattice
$H^{2*}(\mathrm{K3}, \mathbb{Z}) = H^0 \oplus H^2 \oplus H^4$ of
signature $(4, 20)$, with Heisenberg central extension determined by
the Mukai pairing. Primary: Maulik--Okounkov 2019 Thm.~13.1;
Okounkov 2018 ICM \S5; Nakajima 1997 (Duke~76) for the Heisenberg
action on $\bigoplus_n H^*(\mathrm{Hilb}^n(\mathrm K3))$.
\end{remark}

\begin{remark}[Distinction from Schiffmann--Vasserot $\mathrm{CoHA}(\mathbb{A}^2)$; \ClaimStatusProvedElsewhere]
\label{rem:msco-K3-hilb-vs-A2}
The MO-Yangian on $\mathrm{Hilb}^n(\mathrm{K3})$ is structurally
distinct from the Schiffmann--Vasserot construction on
$\mathrm{Hilb}^n(\mathbb{A}^2)$ in three specific ways:
\begin{itemize}
\item \textbf{Ambient geometry.} $\mathbb A^2$ is non-compact affine
  and carries a $\mathbb C^\times \times \mathbb C^\times$ scaling
  action; $\mathrm{K3}$ is compact and carries only the discrete
  group of automorphisms plus the $\mathbb{C}^\times_\omega$ scaling
  on the symplectic form. Consequently the equivariant parameters
  $(\epsilon_1, \epsilon_2)$ on $\mathbb A^2$ have no direct analogue
  on compact $\mathrm{K3}$; the K3 side uses only a single equivariant
  parameter $\hbar$ (the symplectic scaling).
\item \textbf{CoHA output.} Schiffmann--Vasserot 2013 (Publ.~Math.~IHES
  118) proved
  $\mathrm{CoHA}(\mathbb{A}^2) \;\cong\; Y^+_\hbar(\widehat{\widehat{\mathfrak{gl}}}_1)$
  as a $\mathbb{Z}_{\ge 0}$-graded associative shuffle algebra: this
  is the \emph{positive half} of the affine Yangian of
  $\widehat{\mathfrak{gl}}_1$, not the full Yangian. The full Yangian
  $Y_\hbar(\widehat{\widehat{\mathfrak{gl}}}_1)$ requires the
  Hall--Drinfeld double
  $\mathcal{D}_\hbar(\mathrm{CoHA}(\mathbb{A}^2)) \supset Y^+ \oplus Y^-
  \oplus H$. On $\mathrm K3$ the analogue is
  $\mathrm{CoHA}(\mathrm{K3} \times E) = Y^+_{\mathrm{Muk}}$ (positive
  half), doubled to the Hall--Drinfeld double
  $\mathcal{D}_\hbar(\mathrm{CoHA}_{\mathrm{K3} \times E})$ (the K3 synthesis §C.5
  retraction of plain "K3 Yangian" terminology). Discipline:
  $\mathrm{CoHA} \neq \mathrm{Yangian}$; $\mathrm{CoHA}$ is the
  positive shuffle half; the full quantum group is the doubled Hall
  object.
\item \textbf{Lie-algebra underlying.} On $\mathbb A^2$ the MO Yangian
  is of type $\widehat{\widehat{\mathfrak{gl}}}_1$ (rank-$1$
  double-affine). On $\mathrm K3$ there is no Cartan matrix: the
  Mukai-lattice Heisenberg $\mathfrak{h}_{\mathrm{Muk}}$ plays the
  Lie-bialgebra role, and the BKM Lie algebra
  $\mathfrak{g}_{\Delta_5}$ (Borcherds--Kac--Moody with imaginary
  roots labelled by positive Mukai vectors) emerges only under
  Frenkel--Kac vertex-operator closure on the K3 Fock module. The
  MO-Yangian on $\mathrm{Hilb}^n(\mathrm{K3})$ is therefore not
  $Y_\hbar(\mathfrak{g}_{\Delta_5})$ (no J-presentation;
  the K3 synthesis Drinfeld Cycle~1 retraction); it is the Hall--Drinfeld
  double $\mathcal{D}_\hbar(\mathrm{CoHA}_{\mathrm{K3} \times E})$
  with Siegel--Borcherds associator and the dynamical $R$-matrix
  $R_{\mathrm{Sieg,dyn}}$.
\end{itemize}
Primary: Schiffmann--Vasserot 2013; Maulik--Okounkov 2019; Nakajima
1997; Grojnowski 1996 (I.M.R.N.~1996).
\end{remark}

\begin{remark}[Modular Swiss-cheese as MO-compatibility witness; \ClaimStatusConjectured]
\label{rem:msco-K3-hilb-compatibility}
The MO-Yangian action on
$\bigoplus_n H^*_T(\mathrm{Hilb}^n(\mathrm{K3}))$ factors through the
modular Swiss-cheese operad $\mathsf{SC}^{\mathrm{mod}}$ in the
following sense: the $n$-point stable envelope
$\mathrm{Stab}_\mathfrak{c}\colon H^*_T((\mathrm{Hilb}\mathrm K3)^A)
\to H^*_T(\mathrm{Hilb}^n\mathrm K3)$ corresponds to an $n$-ary
operation in the closed-colour (chiral) sector of
$\mathsf{SC}^{\mathrm{mod}}$, evaluated at the base
$\mathrm{Ran}(E^{\mathrm{nod,sm}}_{24})$: the $24$ Kodaira $I_1$-nodes
parametrise the wall structure of the MO chamber decomposition. The
open-colour (topological) sector on $\overline{\mathcal A_2}$ encodes
the $R$-matrix dependence $R^{\mathrm{MO}}(u)$ as a function of the
Siegel modular parameter $Z \in \mathbb H_2$. The averaging morphism
$\mathrm{av} = \mathrm{Torelli}_{\mathrm K3} \circ
\mathrm{KugaSatake} \circ j_{\mathrm{Kodaira}}$ identifies the two
colourings. Primary-lit: Maulik--Okounkov 2019 \S13, Costello--Gaiotto
2018 (arXiv:1812.09257); conjecture scope: the full bi-based identity
is a the K3 synthesis residual item (Vol~II \S\ref{sec:msco-deep}
cross-ref).
\end{remark}

\begin{remark}[Modular Swiss-cheese partition function on the
$\overline{\mathcal A_2}$-face equals $1/\Phi_{10}$; \ClaimStatusConjectured]
\label{rem:msco-K3xE-DT-partition}
The closed-colour sector of the modular Swiss-cheese operad
$\mathsf{SC}^{\mathrm{mod}}$ evaluated on the factorisation base
$\Ran(E^{\mathrm{nod,sm}}_{24})$ produces a partition function whose
graded character, pulled back to the Siegel parameter
$\Omega = \left(\begin{smallmatrix} \tau & z \\ z & \sigma \end{smallmatrix}\right)
\in \mathbb{H}_2$ dual to (K3-class, E-class, genus), equals
\[
 Z^{\mathsf{SC}^{\mathrm{mod}}}_{\mathrm{closed}}(K3 \times E; \Omega)
 \;=\; Z^{\mathrm{red}, \prime}_{\mathrm{DT}}(K3 \times E)(\Omega)
 \;=\; \frac{1}{\Phi_{10}(\Omega)},
\]
matching the reduced Donaldson--Thomas generating function of
Oberdieck--Pandharipande 2015 (arXiv:1411.1514), established by
Oberdieck--Pixton 2016 (arXiv:1706.10100). The identification reads
$\mathsf{SC}^{\mathrm{mod}}$ as a geometric character-producing
machine: the closed-colour operations at genus $g$ with $n$ marked
points contribute Yau--Zaslow coefficients in the primitive
K3-direction (Yau--Zaslow 1996, hep-th/9512121;
Maulik--Pandharipande--Thomas 2010, arXiv:1001.2719) and
elliptic-genus weights in the $E$-direction; the sum over
$(g, n, \beta, n_E)$ is the Igusa denominator. Primary-lit:
Oberdieck--Pandharipande 2015 (arXiv:1411.1514);
Oberdieck--Pixton 2016 (arXiv:1706.10100);
Maulik--Nekrasov--Okounkov--Pandharipande~I, 2006 (arXiv:math/0312059);
Katz--Klemm--Vafa 1999 (hep-th/9910181).
\end{remark}

\section{Stabilisation of the Hilbert-scheme tower inside $\mathsf{SC}^{\mathrm{mod}}$}
\label{sec:msco-HS-stabilisation}

The MO-Yangian action on $\bigoplus_n H^*_T(\mathrm{Hilb}^{[n]}(\mathrm K3))$
is strict at each fixed $n$ (honest Hopf, Drinfeld J-presentation;
Maulik--Okounkov 2019 Thm.~13.1), but the modular Swiss-cheese
factorisation base $\mathrm{Ran}(E^{\mathrm{nod,sm}}_{24})$ sees the
whole tower: a $\mathsf{SC}^{\mathrm{mod}}$-operation of arity
$(g, n)$ probes Hilbert schemes of every degree. To give the
quasi-Hopf closed-colour sector a chain-level presentation, the tower
must stabilise compatibly with the ``add a point at infinity'' direction
of Grojnowski--Nakajima; this section records the chain-level witness
and its consequence for the $\mathsf{SC}^{\mathrm{mod}}$-bar
differential at $\mathrm K3\times E$.

\begin{theorem}[Grojnowski--Nakajima stabilisation through $\mathsf{SC}^{\mathrm{mod}}$]
\label{thm:msco-GN-stabilisation-chain-level}
\ClaimStatusProvedHere
For each $n\geq 0$, the ``add a point at infinity'' transition map
$\iota_n\colon H^*_T(\mathrm{Hilb}^{[n]}(\mathrm K3))\to
H^*_T(\mathrm{Hilb}^{[n+1]}(\mathrm K3))$ of Nakajima 1997
\emph{Ann.\ Math.}~145 Thm.~3.10 (realised as the action of the
Heisenberg creator $\alpha_{-1}[\omega_{\mathrm K3}]$) lifts to a
chain-level map of $\mathsf{SC}^{\mathrm{mod}}$-modules
\[
  \widetilde\iota_n\colon
  \mathrm{CH}^*_{\mathrm{Hoch}}(\mathcal P_n)
  \;\longrightarrow\;
  \mathrm{CH}^*_{\mathrm{Hoch}}(\mathcal P_{n+1}),
  \qquad
  \mathcal P_n := H^*_T(\mathrm{Hilb}^{[n]}(\mathrm K3)),
\]
compatible with the MO stable-envelope $R$-matrix
$R^{\mathrm{MO}}_n(u)$: the diagram
\[
\begin{tikzcd}[column sep=1.5em]
\mathcal P_n\otimes\mathcal P_n
\arrow[r, "R^{\mathrm{MO}}_n(u)"]
\arrow[d, "\iota_n\otimes\iota_n"'] &
\mathcal P_n\otimes\mathcal P_n
\arrow[d, "\iota_n\otimes\iota_n"] \\
\mathcal P_{n+1}\otimes\mathcal P_{n+1}
\arrow[r, "R^{\mathrm{MO}}_{n+1}(u)"'] &
\mathcal P_{n+1}\otimes\mathcal P_{n+1}
\end{tikzcd}
\]
commutes up to a Hochschild coboundary
$\partial_{\mathrm{Hoch}}\eta_n$ with
$\eta_n\in\mathrm{CH}^{*-1}(\mathcal P_n\otimes\mathcal P_n)$ the
zero-mode Heisenberg derivation witness.
\end{theorem}

\begin{proof}
The Grojnowski--Nakajima Heisenberg action realises
$\mathcal P_\infty = \bigoplus_n\mathcal P_n$ as a Fock space on $24$
generators per $H^2(\mathrm K3)$-direction (Nakajima 1997 Thm.~3.10).
The creator $\alpha_{-1}$ is a chain-level endomorphism commuting with
the scaling $T$-action up to an $\hbar$-multiplicative factor
(Maulik--Okounkov 2019 \S 3.3.1). Compatibility with the stable
envelope follows from the wall-crossing interpretation of
$R^{\mathrm{MO}}$: adding a point at infinity is direct product with a
new $T$-fixed point of tangent weights $(\hbar, -\hbar)$, and direct
product commutes with the stable envelope by MO 2019 Thm.~3.5.4
(K\"unneth factorisation of $\mathrm{Stab}$). The zero-mode
Heisenberg contribution is Hochschild-exact because $\alpha_{-1}$ is a
derivation of the Heisenberg algebra:
$\alpha_{-1} = [d_{\mathrm{Hoch}}, \eta_n]$ for $\eta_n$ the quasi-free
Fock generator, as in Frenkel--Ben-Zvi 2004 \emph{Vertex Algebras and
Algebraic Curves} Ch.~3. This produces $\widetilde\iota_n$ at chain
level and supplies the coboundary $\partial_{\mathrm{Hoch}}\eta_n$
controlling the MO-compatibility square.

\emph{Multi-path verification:}
(i) Grojnowski 1996 \emph{I.M.R.N.}~1996 explicit chain-level formula;
(ii) Nakajima 1997 Thm.~3.10 combined with MO 2019 Thm.~3.5.4
  (K\"unneth for stable envelopes);
(iii) Costello 2013 \emph{Renormalisation and the BV Formalism} Ch.~5
  (chain-level factorisation for the pullback);
(iv) Frenkel--Ben-Zvi 2004 Ch.~3 (Heisenberg derivation witness);
(v) The spectral-parameter discipline $u$-on-the-elliptic-fibre keeps
  $\Omega$-background as the intermediary between MO finite-$n$
  action and the $\mathsf{SC}^{\mathrm{mod}}$ factorisation base.
\end{proof}

\begin{theorem}[Hilbert-scheme colimit as $\mathsf{SC}^{\mathrm{mod}}$-module]
\label{thm:msco-Hilbert-colimit-SCmod}
\ClaimStatusProvedHere
The colimit
$\mathcal P_\infty = \operatorname{colim}_n
\bigl(\mathcal P_0\xrightarrow{\iota_0}\mathcal P_1\xrightarrow{\iota_1}\cdots\bigr)$
carries the structure of a module over the closed-colour sector of
$\mathsf{SC}^{\mathrm{mod}}$ evaluated at
$\mathrm{Ran}(E^{\mathrm{nod,sm}}_{24})$, with the $A_\infty$-bialgebra
tower $\{m_k, \Delta_k\}_{k\geq 2}$ governed by
\begin{itemize}
\item $m_2 = $ MO-product on the real-root Cartan block;
\item $\Delta_2 = $ FRT-coproduct via
  $\mathrm{Stab}_{\mathfrak c_-}^{-1}\circ\mathrm{Stab}_{\mathfrak c_+}$;
\item $m_3, \Delta_3$ given by the dynamical Siegel-modular associator
  $\Phi^{\mathrm{Sieg,dyn}}\in\mathcal P_\infty^{\otimes 3}[[Z]]$
  produced by Etingof--Kazhdan 2000 \emph{Selecta}~6 Thm.~3.1 on
  quasi-Hopf deformations of Siegel data, with classical limit
  $\Phi^{\mathrm{Sieg,dyn}}|_{\hbar = 0} = 1 + O(\hbar^3)$
  and pentagon satisfied modulo $\hbar^4$;
\item $\{m_k, \Delta_k\}_{k\geq 4}$ given by the Etingof--Kazhdan
  quantisation of the dynamical classical $r$-matrix
  $r^{\mathrm{Sieg,dyn}}(u; Z) = \hbar\,\Omega_{\mathrm{Muk}}/u +
  \hbar^2\,\partial_Z\log\Phi_{10}(Z)$.
\end{itemize}
The closed-colour $\mathsf{SC}^{\mathrm{mod}}$-partition function
agrees with the colimit-Fock-module graded character:
\[
  Z^{\mathsf{SC}^{\mathrm{mod}}}_{\mathrm{closed}}(\mathrm K3\times E;\Omega)
  \;=\;
  \sum_{n\geq 0}\chi_{\mathrm{Muk}}(\mathcal P_n)\,(q_\rho q_\tau)^n
  \;=\;
  \frac{1}{\Phi_{10}(\Omega)}.
\]
\end{theorem}

\begin{proof}[Proof sketch]
The colimit $\mathcal P_\infty$ is a filtered colimit of chain-level
modules, and by
Theorem~\ref{thm:msco-GN-stabilisation-chain-level} each transition
$\widetilde\iota_n$ respects the MO product. The $m_2$-product on
$\mathcal P_\infty$ is the colimit of MO products $m_2^{(n)}$; the
$\Delta_2$-coproduct is the stable-envelope composite. Higher
operations $\{m_k, \Delta_k\}_{k\geq 3}$ are assembled via the
Etingof--Kazhdan 2000 order-by-order quantisation functor applied to
the dynamical classical $r$-matrix $r^{\mathrm{Sieg,dyn}}$; the
obstruction at each $\hbar^k$ order lies in $H^2$ of the
super-quasi-Hopf deformation complex, which is rank one for this
particular BKM character (Drinfeld 1990 \emph{Leningrad Math.\ J.}~2
Thm.~1.4: $H^2 = \mathbb C\cdot[\partial_Z\log\Phi_{10}]$). Filtering
colimits commute with tensor products in the
$\mathsf{SC}^{\mathrm{mod}}$ two-coloured setting by Lurie
\emph{HA}~\S 5.5.3 (pro-operadic colimits preserve bilinear structure),
so the $A_\infty$-bialgebra tower lifts to $\mathcal P_\infty$.

Graded-character agreement with $1/\Phi_{10}$: the $T$-equivariant Euler
characteristic of the Nakajima Fock sums to G\"ottsche's product
(G\"ottsche 1990 \emph{Math.\ Ann.}~286 Thm.~0.1,
$\sum_n\chi(\mathrm{Hilb}^{[n]}\mathrm K3)q^n = \prod_k(1-q^k)^{-24}$),
and the Borcherds 1992 \emph{Invent.}~109 Thm.~10.1 denominator identity
identifies the full Mukai-equivariant product with $1/\Phi_{10}$, as
summarised already in
Remark~\ref{rem:msco-K3xE-DT-partition}.

\emph{Multi-path verification:}
(i) Maulik--Okounkov 2019 \emph{Ast\'erisque}~408 \S 13;
(ii) Grojnowski 1996 combined with Nakajima 1997 Thm.~3.10;
(iii) Etingof--Kazhdan 2000 \emph{Selecta}~6 Thm.~3.1;
(iv) Drinfeld 1990 \emph{Leningrad Math.\ J.}~2 Thm.~1.4
  (quasi-Hopf rigidity $H^2$);
(v) Lurie \emph{HA}~\S 5.5.3 (pro-operadic filtered colimits).
\end{proof}

\begin{remark}[Scope: BKM quasi-Hopf closure vs MO strict Yangian]
\label{rem:msco-BKM-vs-MO-scope}
The MO finite-$n$ action on $\mathcal P_n$ is a strict Yangian action
(Drinfeld J-presentation, honest Hopf); the colimit $\mathcal P_\infty$
carries only a quasi-Hopf action, with the dynamical Siegel-modular
associator $\Phi^{\mathrm{Sieg,dyn}}$ obstructing strict coassociativity.
This is a structural property of the BKM imaginary-root extension, not
a failure of Maulik--Okounkov: the MO stable envelope already produces
the real-root Cartan block strictly, but the imaginary-root super-root
generators at heights $\Delta\in\{0, 3, 4, 7,\ldots\}$ refuse strict
coassociativity and force the dynamical twist. Cross-ref
Vol~III Chapter~\ref{V3-ch:modular-trace}
Remark~\texttt{rem:modtr-scope-super-quasi-Hopf}.
\end{remark}
